% Economic growth and Development — Economics Upper Sixth — Cours signé OnBuch+
% Official MINESEC syllabus. Compiler avec : tectonic ECO07-economic-growth-and-development.tex
\newcommand{\DOCMATIERE}{Economics}
\newcommand{\DOCNIVEAU}{Upper Sixth}
\newcommand{\DOCMODULE}{Upper Sixth — Economics}
\newcommand{\DOCLECON}{7}
\newcommand{\DOCTITRE}{Economic growth and Development}
% OnBuch+ — préambule « cours signés » ENRICHI (compilé avec tectonic / XeLaTeX)
% Système visuel « Pop » d'OnBuch+ : fond crème (jamais blanc), bordures encre
% épaisses, ombres « dures » décalées, titres Archivo Black en capitales.
% Version MATHÉMATIQUES : graphes (pgfplots/tikz), boîtes pédagogiques
% (définition, propriété, méthode, attention, le savais-tu, exercice résolu,
% exercice, TP, bilan), page de garde et signature OnBuch+ ancrée en bas.
%
% Macros attendues dans main.tex AVANT \input{preamble} :
%   \DOCMATIERE  \DOCNIVEAU  \DOCMODULE  \DOCLECON (numéro)  \DOCTITRE
\documentclass[11pt]{article}

\usepackage{fontspec}
\usepackage[english]{babel}
\usepackage[a4paper,margin=2cm,top=3.4cm,bottom=2.8cm,headheight=1.4cm,footskip=1.3cm]{geometry}
\usepackage{amsmath,amssymb,amsfonts}
\usepackage{mathtools} % \xmapsto, \coloneqq... souvent utilisés par les modèles
\usepackage{cancel} % simplifications barrées (bilans de forces)
% \abs{z} (module, valeur absolue) et \norm{u} (norme), utilisés par les modèles
\providecommand{\abs}{}\let\abs\relax\DeclarePairedDelimiter{\abs}{\lvert}{\rvert}
\providecommand{\norm}{}\let\norm\relax\DeclarePairedDelimiter{\norm}{\lVert}{\rVert}
\usepackage[table]{xcolor} % [table] : fournit aussi \rowcolors (alternance de lignes)
\usepackage{pagecolor}
\usepackage{tikz}
\usetikzlibrary{shapes.symbols,circuits.logic.US,circuits.logic.IEC,arrows.meta,positioning,calc,shapes.geometric,shapes.misc,shapes.arrows,decorations.pathmorphing,decorations.pathreplacing,decorations.markings,patterns,shadows,mindmap,trees,fit,backgrounds,matrix,intersections,angles,quotes}
\usepackage{pgfplots}
\usepackage[version=4]{mhchem} % équations nucléaires \ce{^{235}_{92}U}
\usepackage[european,siunitx]{circuitikz} % schémas électriques
\pgfplotsset{compat=1.18}
\usepgfplotslibrary{fillbetween,groupplots,polar,statistics,dateplot} % aires entre courbes (calcul intégral)
\usepackage{enumitem}
\usepackage{fancyhdr}
% [most] charge toutes les bibliothèques tcolorbox (listings, minted, raster,
% documentation...) et gonfle inutilement la mémoire TeX sur un document long
% (~300 boîtes) : on ne charge que ce qui est réellement utilisé.
\usepackage[skins,breakable]{tcolorbox}
\usepackage{graphicx}
\usepackage{colortbl}
\usepackage{booktabs}
\usepackage{tabularx}
\usepackage{diagbox} % case d'en-tête diagonale des tableaux à double entrée
% Colonnes extensibles centrée / gauche / droite (C, L, R), souvent utilisées par les modèles
\newcolumntype{C}{>{\centering\arraybackslash}X}
\newcolumntype{L}{>{\raggedright\arraybackslash}X}
\newcolumntype{R}{>{\raggedleft\arraybackslash}X}
\usepackage{array}
\usepackage{multicol}
\usepackage{siunitx}
% parse-numbers=false : accepte aussi \SI{2,5 \times 10^{-3}}{...}
\sisetup{locale=UK,output-decimal-marker={.},group-minimum-digits=4,parse-numbers=false}
\DeclareSIUnit\ohm{\ensuremath{\symup{\Omega}}} % U+2126 absent de la police
\DeclareSIUnit\division{div} % réglages d'oscilloscope (ms/div, V/div)
\DeclareSIUnit\year{yr}\DeclareSIUnit\annum{yr}\DeclareSIUnit\Ma{Ma}\DeclareSIUnit\tr{tr} % tours (vitesse de rotation, ex. \SI{1500}{\tr\per\minute})
\usepackage{qrcode}
% \degree : commande fréquemment utilisée par les modèles pour un angle ou une
% température en dehors de \SI{}{} (non fournie par les paquets chargés).
\providecommand{\degree}{^{\circ}}
% \qed (fin de démonstration), utilisé par les modèles sans amsthm
\providecommand{\qed}{\hfill$\square$}
% \barre : double barre des valeurs interdites dans un tableau de variations (tkz-tab)
\providecommand{\barre}{\ensuremath{\|}}
% environnement proof (amsthm non chargé) utilisé par les modèles
\ifdefined\proof\else\newenvironment{proof}[1][Proof]{\par\noindent\textit{#1.}\ }{\qed\par}\fi
\usepackage{lastpage}
\usepackage{hyperref}
\hypersetup{hidelinks}

% ============================================================
% Palette « Pop » OnBuch+ — mêmes hex que la classe P (plus_ui.dart)
% ============================================================
\definecolor{poporange}{HTML}{F59321}
\definecolor{popdark}{HTML}{E07A0C}
\definecolor{popcream}{HTML}{FFF6E8}
\definecolor{popink}{HTML}{1C1714}
\definecolor{poppurple}{HTML}{7A5AE0}
\definecolor{popgreen}{HTML}{2E9E6A}
\definecolor{poppink}{HTML}{E0457B}
\definecolor{popblue}{HTML}{2F7DE1}
\definecolor{popgold}{HTML}{C98A12}
\definecolor{popmuted}{HTML}{8A7F76}
\definecolor{popline}{HTML}{EADFD2}
\definecolor{popgray}{HTML}{8A7F76}
\colorlet{popgrey}{popgray}
% Alias tolérants (le modèle nomme parfois les couleurs différemment)
\colorlet{popwhite}{white}
\colorlet{popblack}{popink}
\colorlet{popred}{poppink}
\colorlet{popyellow}{popgold}
% Teintes claires pour fonds de tableaux / zones
\colorlet{popblueL}{popblue!10!white}
\colorlet{poporangeL}{poporange!14!white}
\colorlet{popgreenL}{popgreen!12!white}
\colorlet{poppurpleL}{poppurple!12!white}
\colorlet{poppinkL}{poppink!12!white}
\colorlet{popgoldL}{popgold!14!white}

\pagecolor{popcream}

% ============================================================
% Typographie
% ============================================================
\defaultfontfeatures{Ligatures=TeX}
\setmainfont{PlusJakartaSans-Medium.ttf}[Path=../../../fonts/,BoldFont=PlusJakartaSans-Bold.ttf]
% Maths en sans-serif (Fira Math) avec chiffres et lettres droites de Plus
% Jakarta, pour que les formules se fondent dans le texte.
\usepackage[math-style=ISO,bold-style=ISO]{unicode-math}
\setmathfont{FiraMath-Regular.otf}[Path=../../../fonts/]
\setmathfont{PlusJakartaSans-Medium.ttf}[Path=../../../fonts/,range={up/num,up/{Latin,latin},\mathup}]
% (icomma retiré : décimales avec point en anglais)
% Caractères mathématiques Unicode absents des polices de texte (Plus Jakarta,
% Archivo Black) : rendus par leur équivalent mathématique, en texte comme en
% formule — sinon ils s'affichent en carré vide (ex. « Arithmétique dans ℤ »).
\usepackage{newunicodechar}
\newunicodechar{ℝ}{\ensuremath{\mathbb{R}}}
\newunicodechar{ℤ}{\ensuremath{\mathbb{Z}}}
\newunicodechar{ℂ}{\ensuremath{\mathbb{C}}}
\newunicodechar{ℕ}{\ensuremath{\mathbb{N}}}
\newunicodechar{ℚ}{\ensuremath{\mathbb{Q}}}
\newunicodechar{𝔻}{\ensuremath{\mathbb{D}}}
\newunicodechar{α}{\ensuremath{\alpha}}
\newunicodechar{β}{\ensuremath{\beta}}
\newunicodechar{γ}{\ensuremath{\gamma}}
\newunicodechar{δ}{\ensuremath{\delta}}
\newunicodechar{ε}{\ensuremath{\varepsilon}}
\newunicodechar{θ}{\ensuremath{\theta}}
\newunicodechar{λ}{\ensuremath{\lambda}}
\newunicodechar{μ}{\ensuremath{\mu}}
\newunicodechar{σ}{\ensuremath{\sigma}}
\newunicodechar{τ}{\ensuremath{\tau}}
\newunicodechar{φ}{\ensuremath{\varphi}}
\newunicodechar{ψ}{\ensuremath{\psi}}
\newunicodechar{ω}{\ensuremath{\omega}}
\newunicodechar{Σ}{\ensuremath{\Sigma}}
% majuscules produites par \MakeUppercase dans les titres (α-aminés -> Α)
\newunicodechar{Α}{\ensuremath{\alpha}}
\newunicodechar{Β}{\ensuremath{\beta}}
\newunicodechar{Γ}{\ensuremath{\Gamma}}
\newunicodechar{Δ}{\ensuremath{\Delta}}
\newunicodechar{Ε}{\ensuremath{\varepsilon}}
\newunicodechar{Θ}{\ensuremath{\Theta}}
\newunicodechar{Λ}{\ensuremath{\Lambda}}
\newunicodechar{Μ}{\ensuremath{\mu}}
\newunicodechar{Τ}{\ensuremath{\tau}}
\newunicodechar{Φ}{\ensuremath{\Phi}}
\newunicodechar{Ψ}{\ensuremath{\Psi}}
\newunicodechar{Ω}{\ensuremath{\Omega}}
\newunicodechar{↦}{\ensuremath{\mapsto}}
\newunicodechar{∈}{\ensuremath{\in}}
\newunicodechar{∉}{\ensuremath{\notin}}
\newunicodechar{∧}{\ensuremath{\wedge}}
\newunicodechar{⇔}{\ensuremath{\Leftrightarrow}}
\newunicodechar{⟺}{\ensuremath{\Longleftrightarrow}}
\newunicodechar{⇒}{\ensuremath{\Rightarrow}}
\newunicodechar{⟹}{\ensuremath{\Longrightarrow}}
\newunicodechar{∘}{\ensuremath{\circ}}
\newunicodechar{∪}{\ensuremath{\cup}}
\newunicodechar{∩}{\ensuremath{\cap}}
\newunicodechar{≡}{\ensuremath{\equiv}}
\newunicodechar{⊂}{\ensuremath{\subset}}
\newunicodechar{⊥}{\ensuremath{\perp}}
\newunicodechar{∃}{\ensuremath{\exists}}
\newunicodechar{∀}{\ensuremath{\forall}}
\newunicodechar{∛}{\ensuremath{\sqrt[3]{\,}}}
\newunicodechar{∜}{\ensuremath{\sqrt[4]{\,}}}
\newunicodechar{⃗}{\ensuremath{{}^{\to}}}
\newunicodechar{ⁿ}{\ifmmode^{n}\else\textsuperscript{\ensuremath{n}}\fi}
\newunicodechar{ˣ}{\ifmmode^{x}\else\textsuperscript{\ensuremath{x}}\fi}
\newunicodechar{ᵇ}{\ifmmode^{b}\else\textsuperscript{\ensuremath{b}}\fi}
\newunicodechar{ʳ}{\ifmmode^{r}\else\textsuperscript{\ensuremath{r}}\fi}
\newunicodechar{ᵘ}{\ifmmode^{u}\else\textsuperscript{\ensuremath{u}}\fi}
\newunicodechar{ʸ}{\ifmmode^{y}\else\textsuperscript{\ensuremath{y}}\fi}
\newunicodechar{ᵃ}{\ifmmode^{a}\else\textsuperscript{\ensuremath{a}}\fi}
\newunicodechar{ᵉ}{\ifmmode^{e}\else\textsuperscript{\ensuremath{e}}\fi}
\newunicodechar{ᵏ}{\ifmmode^{k}\else\textsuperscript{\ensuremath{k}}\fi}
\newunicodechar{ᵖ}{\ifmmode^{p}\else\textsuperscript{\ensuremath{p}}\fi}
\newunicodechar{ᴺ}{\ifmmode^{N}\else\textsuperscript{\ensuremath{N}}\fi}
\newunicodechar{ᶜ}{\ifmmode^{c}\else\textsuperscript{\ensuremath{c}}\fi}
\newunicodechar{ʰ}{\ifmmode^{h}\else\textsuperscript{\ensuremath{h}}\fi}
\newunicodechar{ᵠ}{\ifmmode^{q}\else\textsuperscript{\ensuremath{q}}\fi}
\newunicodechar{ₙ}{\ifmmode_{n}\else\textsubscript{\ensuremath{n}}\fi}
\newunicodechar{ₐ}{\ifmmode_{a}\else\textsubscript{\ensuremath{a}}\fi}
\newunicodechar{ᵢ}{\ifmmode_{i}\else\textsubscript{\ensuremath{i}}\fi}
\newunicodechar{ᵣ}{\ifmmode_{r}\else\textsubscript{\ensuremath{r}}\fi}
\newunicodechar{ₖ}{\ifmmode_{k}\else\textsubscript{\ensuremath{k}}\fi}
\newunicodechar{ᵦ}{\ifmmode_{\beta}\else\textsubscript{\ensuremath{\beta}}\fi}
\newunicodechar{ₓ}{\ifmmode_{x}\else\textsubscript{\ensuremath{x}}\fi}
\newfontfamily\popdisplay{ArchivoBlack-Regular.ttf}[Path=../../../fonts/]
\newfontfamily\poplogo{SpaceGrotesk-Bold.ttf}[Path=../../../fonts/]
\newfontfamily\popsemi{PlusJakartaSans-SemiBold.ttf}[Path=../../../fonts/]
\newfontfamily\popxbold{PlusJakartaSans-ExtraBold.ttf}[Path=../../../fonts/]
\newfontfamily\popxboldls{PlusJakartaSans-ExtraBold.ttf}[Path=../../../fonts/,LetterSpace=3.0]

\setlist[enumerate]{leftmargin=1.7em,itemsep=5pt,topsep=4pt}
\setlist[itemize]{leftmargin=1.7em,itemsep=4pt,topsep=4pt,label={\color{poporange}\textbullet}}
\setlength{\parindent}{0pt}
\setlength{\parskip}{5pt}
\renewcommand{\arraystretch}{1.25}

% pgfplots : style Pop par défaut
\pgfplotsset{
  every axis/.append style={axis line style={popink,line width=1pt},tick style={popink},
    grid style={popline,line width=0.5pt},label style={font=\small},tick label style={font=\footnotesize}},
  popcurve/.style={poporange,line width=1.8pt,no markers,smooth},
}
% Même style disponible dans un \draw TikZ simple (hors pgfplots)
\tikzset{popcurve/.style={poporange,line width=1.8pt}}
\tikzset{popgrid/.style={popline,very thin}}
\colorlet{popcurve}{poporange} % le nom du style est parfois utilisé comme couleur

% ============================================================
% Pill / squiggle
% ============================================================
\newcommand{\poppill}[3]{%
  \tikz[baseline=-0.55ex]\node[draw=popink,line width=1.2pt,rounded corners=2.6pt,%
    fill=#2,inner xsep=6.5pt,inner ysep=3pt]{%
    \popxboldls\fontsize{7}{7}\selectfont\color{#3}\MakeUppercase{#1}};%
}
\newcommand{\popsquiggle}[1]{%
  \tikz[baseline=-0.4ex]{
    \draw[#1,line width=2.6pt,line cap=round]
      plot [smooth] coordinates {(0,0.09) (0.34,0.01) (0.68,0.21) (1.02,0.01) (1.36,0.10)};
  }%
}
% Mot-clé surligné (marqueur orange sous le texte)
\newcommand{\cle}[1]{{\popxbold\color{popdark}#1}}

% ============================================================
% En-tête / pied
% ============================================================
\pagestyle{fancy}
\fancyhf{}
\renewcommand{\headrulewidth}{0pt}
\renewcommand{\footrulewidth}{0pt}
\lhead{\raisebox{-0.15ex}{\poplogo\fontsize{13}{13}\selectfont\color{popink}OnBuch\color{poporange}+}}
\rhead{\poppill{\DOCMATIERE\ \textbullet\ \DOCNIVEAU\ \textbullet\ Chapter \DOCLECON}{white}{popink}}
\lfoot{\popxbold\fontsize{7.6}{7.6}\selectfont\color{popmuted}\MakeUppercase{Signed course OnBuch+}\ \textbullet\ {\popsemi\DOCTITRE}}
\rfoot{\tikz[baseline=-0.6ex]\node[rounded corners=8.5pt,fill=popink,minimum height=17pt,minimum width=17pt,inner xsep=5pt,inner sep=0]{\popxbold\fontsize{8}{8}\selectfont\color{popcream}\thepage\,/\,\pageref*{LastPage}};}

% ============================================================
% Titres
% ============================================================
\newcommand{\chaptitre}[2]{%
  \begin{tcolorbox}[enhanced,colback=poporange,colframe=popink,boxrule=2.6pt,arc=11pt,
      drop shadow=popink,left=15pt,right=15pt,top=12pt,bottom=11pt,before skip=3pt,after skip=18pt]
    \poppill{#2}{popink}{popcream}\par\vspace{9pt}
    {\raggedright\hyphenpenalty=10000\exhyphenpenalty=10000\popdisplay\fontsize{22}{24}\selectfont\color{popcream}\MakeUppercase{#1}\par}
  \end{tcolorbox}
}
% Section numérotée : pastille orange avec le numéro + titre Archivo
\newcounter{popsec}
\newcommand{\coursec}[1]{%
  \stepcounter{popsec}\setcounter{popsub}{0}%
  \par\vspace{16pt}\noindent
  \begin{minipage}{\linewidth}
  \tikz[baseline=-0.7ex]\node[draw=popink,line width=1.6pt,fill=poporange,rounded corners=4pt,minimum size=22pt,inner sep=0]{\popdisplay\fontsize{12}{12}\selectfont\color{popcream}\thepopsec};%
  \hspace{8pt}{\popdisplay\fontsize{15}{17}\selectfont\color{popink}\MakeUppercase{#1}}\par
  \vspace{4pt}\noindent{\color{poporange}\rule{36pt}{3.6pt}}
  \end{minipage}\par\nopagebreak\vspace{8pt}
}
\newcounter{popsub}
\newcommand{\courssub}[1]{%
  \stepcounter{popsub}%
  \par\vspace{9pt}\noindent{\popxbold\fontsize{11.5}{13}\selectfont\color{popink}{\color{poporange}\thepopsec.\thepopsub}\hspace{6pt}#1}\par\nopagebreak\vspace{4pt}
}

% ============================================================
% Boîtes pédagogiques — même recette : fond blanc, bordure épaisse colorée,
% ombre dure encre, badge plein. Titre optionnel [#1].
% ============================================================
\tcbset{popbase/.style={breakable,enhanced,colback=white,boxrule=2.3pt,arc=9pt,
  shadow={3pt}{-3pt}{0pt}{popink},left=14pt,right=14pt,top=11pt,bottom=11pt,before skip=11pt,after skip=15pt,
  fontupper=\popsemi\fontsize{10.6}{14.5}\selectfont\color{popink},
  fontlower=\popsemi\fontsize{10.6}{14.5}\selectfont\color{popink}}}
% Coche dessinée (le glyphe ✓ n'existe pas dans les polices du document)
\newcommand{\popcheck}[1]{\tikz[baseline=-0.5ex]\draw[#1,line width=1.6pt,line cap=round,line join=round](-0.13,0.0)--(-0.04,-0.09)--(0.13,0.1);}
\providecommand{\coche}{\popcheck{popgreen}}
\providecommand{\resultat}[1]{\par\smallskip\noindent\fcolorbox{popgreen}{popgreenL}{\parbox{\dimexpr\linewidth-2\fboxsep-2\fboxrule\relax}{\textbf{Result.} #1}}\par\smallskip}
\newcommand{\popboxhead}[3]{\poppill{#1}{#2}{white}\ifx\relax#3\relax\else\hspace{6pt}{\popxbold\fontsize{10.6}{12}\selectfont\color{popink}#3}\fi\par\vspace{7pt}}

\newtcolorbox{aretenir}[1][]{popbase,colframe=popgreen,before upper={\popboxhead{Key points}{popgreen}{#1}}}
\newtcolorbox{exemplebox}[1][]{popbase,colframe=poporange,before upper={\popboxhead{Exemple}{poporange}{#1}}}
\newtcolorbox{definition}[1][]{popbase,colframe=popblue,before upper={\popboxhead{Definition}{popblue}{#1}}}
\newtcolorbox{propriete}[1][]{popbase,colframe=poppurple,before upper={\popboxhead{Property}{poppurple}{#1}}}
\newtcolorbox{methode}[1][]{popbase,colframe=popgold,before upper={\popboxhead{Method}{popgold}{#1}}}
\newtcolorbox{attention}[1][]{popbase,colframe=poppink,before upper={\popboxhead{Warning}{poppink}{#1}}}
\newtcolorbox{experience}[1][]{popbase,colframe=popblue,colback=popblueL,before upper={\popboxhead{Experiment}{popblue}{#1}}}
% « Le savais-tu ? » : carte violette pleine (contexte, culture, Cameroun)
\newtcolorbox{savaistu}[1][]{popbase,colback=poppurple,colframe=popink,
  fontupper=\popsemi\fontsize{10.4}{14.2}\selectfont\color{white},
  before upper={\poppill{Did you know?}{white}{poppurple}\ifx\relax#1\relax\else\hspace{6pt}{\popxbold\color{white}#1}\fi\par\vspace{7pt}}}
% Exercice résolu : énoncé en haut, \tcblower puis la solution sur fond crème
\newcounter{popexo}\newcounter{popexr}
\newtcolorbox{exoresolu}[1][]{popbase,colframe=poporange,colbacklower=popcream,
  segmentation style={popink,line width=1.2pt,dashed},
  before upper={\stepcounter{popexr}\popboxhead{Worked example \thepopexr}{poporange}{#1}},
  before lower={\poppill{Solution}{popink}{popcream}\par\vspace{6pt}}}
% Exercice (non résolu en place) — la correction est dans la partie Corrigés
\newtcolorbox{exercice}[1][]{popbase,colframe=popink,
  before upper={\stepcounter{popexo}\popboxhead{Exercise \thepopexo}{popink}{#1}}}
% Corrigé
\newtcolorbox{corrige}[1][]{popbase,colframe=popgreen,colback=popgreenL,
  before upper={\popboxhead{Solution}{popgreen}{#1}}}
% Objectifs / prérequis / bilan
\newtcolorbox{objectifs}{popbase,colframe=popink,colback=poporangeL,
  before upper={\popboxhead{Chapter objectives}{popink}{}}}
\newtcolorbox{prerequis}{popbase,colframe=popink,colback=popblueL,
  before upper={\popboxhead{Prerequisites}{popblue}{}}}
\newtcolorbox{bilan}{popbase,colframe=popink,colback=white,boxrule=2.8pt,
  before upper={\popboxhead{Fiche bilan}{poporange}{L'essentiel à connaître pour l'examen}}}
% Figure encadrée avec légende
\newtcolorbox{popfigure}[1][]{enhanced,breakable=false,colback=white,colframe=popline,boxrule=1.4pt,arc=8pt,
  left=10pt,right=10pt,top=10pt,bottom=8pt,before skip=10pt,after skip=12pt,halign=center,
  lower separated=false,halign lower=center,
  fontlower=\popsemi\fontsize{9}{11.5}\selectfont\color{popmuted},
  #1}
\newcommand{\legende}[1]{\par\vspace{6pt}{\popsemi\fontsize{9}{11.5}\selectfont\color{popmuted}#1}}

% ============================================================
% Signature OnBuch+
% ============================================================
\newcommand{\poplogoinline}[1]{{\poplogo\fontsize{#1}{#1}\selectfont\color{popink}OnBuch\color{poporange}+}}
\newcommand{\signatureonbuch}{%
  \begin{tcolorbox}[enhanced,colback=popink,colframe=popink,boxrule=2.2pt,arc=10pt,
      drop shadow=poporange,left=14pt,right=14pt,top=10pt,bottom=10pt,before skip=0pt,after skip=0pt]
    \tikz[baseline=-0.7ex]\node[circle,fill=popgreen,draw=popcream,line width=1.3pt,minimum size=22pt,inner sep=0]{\popcheck{white}};%
    \hspace{9pt}\begin{minipage}[c]{0.62\linewidth}
      {\popxbold\fontsize{12}{13}\selectfont\color{popcream}Signed\ \textbullet\ {\poplogo OnBuch\color{poporange}+}}\par\vspace{2pt}
      {\raggedright\popsemi\fontsize{8}{10.5}\selectfont\color{popline}\MakeUppercase{OnBuch+ teaching team}\\\MakeUppercase{Follows the official MINESEC syllabus}\par}
    \end{minipage}\hfill
    {\poplogo\fontsize{22}{22}\selectfont\color{popcream}OnBuch\color{poporange}+}
  \end{tcolorbox}%
}

% ============================================================
% Page de garde
% #1 titre · #2 matière · #3 niveau · #4 module · #5 n° leçon · #6 durée ·
% #7 description / objectifs (texte ou itemize)
% ============================================================
\newcommand{\pagegarde}[7]{%
  \thispagestyle{empty}
  \newgeometry{a4paper,margin=1.7cm,top=1.6cm,bottom=1.4cm}%
  \begin{tikzpicture}[remember picture,overlay]
    \fill[poporange!18] ([xshift=-3.2cm,yshift=-3.6cm]current page.north east) circle (4.6cm);
    \fill[poppurple!14] ([xshift=2.4cm,yshift=7.6cm]current page.south west) circle (3.8cm);
  \end{tikzpicture}%
  {\poplogo\fontsize{30}{30}\selectfont\color{popink}OnBuch\color{poporange}+}\hfill
  \raisebox{0.6ex}{\poppill{Signed course \textbullet\ Premium}{popink}{popcream}}\par
  \vspace{1pt}\popsquiggle{poppurple}\par
  \vspace{8pt}
  \begin{tcolorbox}[enhanced,colback=poporange,colframe=popink,boxrule=2.8pt,arc=13pt,
      drop shadow=popink,left=18pt,right=18pt,top=12pt,bottom=13pt,after skip=10pt]
    \poppill{#2 \textbullet\ #3}{popink}{popcream}\hspace{5pt}\poppill{#4}{white}{popink}\par\vspace{8pt}
    {\popdisplay\fontsize{14}{14}\selectfont\color{popink}CHAPTER #5}\par\vspace{4pt}
    {\raggedright\hyphenpenalty=10000\exhyphenpenalty=10000\popdisplay\fontsize{25}{27}\selectfont\color{popcream}\MakeUppercase{#1}\par}
    \vspace{8pt}
    {\popxbold\fontsize{9.5}{9.5}\selectfont\color{popink}\MakeUppercase{Suggested time: #6}}
  \end{tcolorbox}
  \begin{tcolorbox}[enhanced,colback=white,colframe=popink,boxrule=2.2pt,arc=10pt,
      drop shadow=popink,left=16pt,right=16pt,top=10pt,bottom=10pt,after skip=8pt]
    \poppill{In this chapter}{poppurple}{white}\par\vspace{5pt}
    \popsemi\fontsize{9.3}{12.6}\selectfont\color{popink}#7
  \end{tcolorbox}
  \noindent\begin{minipage}[t]{0.487\linewidth}
    \begin{tcolorbox}[enhanced,colback=popgreen,colframe=popink,boxrule=2.2pt,arc=10pt,
        drop shadow=popink,left=11pt,right=11pt,top=10pt,bottom=10pt,height=3.1cm,valign=center]
      \tcbox[colback=white,colframe=popink,boxrule=1.3pt,arc=5pt,boxsep=0pt,nobeforeafter,
        left=3pt,right=3pt,top=3pt,bottom=3pt,valign=center]{\qrcode[height=1.9cm]{https://play.google.com/store/apps/details?id=com.luvvix.onbuch&hl=en}}%
      \hspace{8pt}\begin{minipage}[c]{0.5\linewidth}
        {\popdisplay\fontsize{11}{12.5}\selectfont\color{white}\MakeUppercase{Download OnBuch}}\par\vspace{4pt}
        {\raggedright\popsemi\fontsize{8.4}{11}\selectfont\color{white}Free \textbullet\ Google Play\\AI tutor, past papers, results\par}
      \end{minipage}
    \end{tcolorbox}
  \end{minipage}\hfill
  \begin{minipage}[t]{0.487\linewidth}
    \begin{tcolorbox}[enhanced,colback=poppurple,colframe=popink,boxrule=2.2pt,arc=10pt,
        drop shadow=popink,left=11pt,right=11pt,top=10pt,bottom=10pt,height=3.1cm,valign=center]
      \tikz[baseline=-0.7ex]\node[circle,fill=white,draw=popink,line width=1.3pt,minimum size=1.6cm,inner sep=0]{\popdisplay\fontsize{24}{24}\selectfont\color{poppurple}+};%
      \hspace{8pt}\begin{minipage}[c]{0.6\linewidth}
        {\popdisplay\fontsize{11}{12.5}\selectfont\color{white}\MakeUppercase{Go OnBuch+}}\par\vspace{4pt}
        {\raggedright\popsemi\fontsize{8.4}{11}\selectfont\color{white}All signed courses, unlimited AI tutor, PDF sheets — OnBuch+ tab in the app\par}
      \end{minipage}
    \end{tcolorbox}
  \end{minipage}
  \vspace{8pt}
  \popcontenu
  \vfill
  \signatureonbuch
  \restoregeometry
  \newpage
}

% Rangée « Ce cours contient » (page de garde)
\newcommand{\poptile}[3]{% #1 couleur #2 gros texte #3 légende
  \begin{tcolorbox}[enhanced,colback=#1,colframe=popink,boxrule=2pt,arc=9pt,shadow={2.5pt}{-2.5pt}{0pt}{popink},
      left=6pt,right=6pt,top=9pt,bottom=9pt,halign=center,valign=center,height=1.75cm,width=\linewidth,nobeforeafter]
    {\popdisplay\fontsize{12.5}{13}\selectfont\color{white}\MakeUppercase{#2}}\par\vspace{3pt}
    {\centering\popsemi\fontsize{7.8}{9.5}\selectfont\color{white}#3\par}
  \end{tcolorbox}}
\newcommand{\popcontenu}{%
  \par\noindent{\popxboldls\fontsize{7.5}{7.5}\selectfont\color{popmuted}THIS SIGNED COURSE CONTAINS}\par\vspace{7pt}
  \noindent\begin{minipage}[t]{0.235\linewidth}\poptile{popblue}{Course}{complete, illustrated, step by step}\end{minipage}\hfill
  \begin{minipage}[t]{0.235\linewidth}\poptile{poporange}{Methods}{and worked examples}\end{minipage}\hfill
  \begin{minipage}[t]{0.235\linewidth}\poptile{poppink}{Exercises}{exam style + solutions}\end{minipage}\hfill
  \begin{minipage}[t]{0.235\linewidth}\poptile{popgreen}{Summary sheet}{the essentials to remember}\end{minipage}\par}

% Sommaire
\newtcolorbox{sommaire}{enhanced,colback=white,colframe=popink,boxrule=2.2pt,arc=10pt,
  drop shadow=popink,left=18pt,right=18pt,top=13pt,bottom=13pt,before skip=6pt,after skip=20pt,
  before upper={\poppill{Contents}{poppurple}{white}\par\vspace{9pt}\popsemi\fontsize{10.6}{14.5}\selectfont\color{popink}}}

% Signature de fin de cours — poussée tout en bas de la dernière page
\newcommand{\findecours}{%
  \par\vfill\nopagebreak
  \begin{center}\popsquiggle{poporange}\end{center}\vspace{4pt}
  \signatureonbuch
}
\AtBeginDocument{\def\llbracket{\mathopen{[\mkern-3.5mu[}}\def\rrbracket{\mathclose{]\mkern-3.5mu]}}} % intervalles d'entiers (glyphe ⟦ absent de la police)
% Glyphes absents de Fira Math : redéfinis à partir de glyphes disponibles
\AtBeginDocument{%
  \def\setminus{\mathbin{\backslash}}\let\smallsetminus\setminus
  \def\vdots{\mathord{\vbox{\baselineskip3.2pt\lineskiplimit0pt\kern4pt\hbox{.}\hbox{.}\hbox{.}}}}%
  \def\ddots{\mathinner{\mkern1mu\raise6.5pt\hbox{.}\mkern2mu\raise3.6pt\hbox{.}\mkern2mu\raise0.7pt\hbox{.}\mkern1mu}}%
  \def\triangle{\mathord{\scalebox{0.9}{$\symup{\Delta}$}}}%
  \def\nabla{\mathord{\raisebox{\depth}{\scalebox{1}[-1]{$\symup{\Delta}$}}}}%
  \def\checkmark{\popcheck{popdark}}%
  \def\star{\mathbin{\ast}}\def\bigstar{\mathord{\ast}}%
  \def\diamond{\mathbin{\scalebox{0.6}{$\Diamond$}}}%
  \def\bigcup{\mathop{\mathchoice{\raisebox{-0.3em}{\scalebox{1.7}{$\cup$}}}{\scalebox{1.25}{$\cup$}}{\cup}{\cup}}\displaylimits}%
  \def\bigcap{\mathop{\mathchoice{\raisebox{-0.3em}{\scalebox{1.7}{$\cap$}}}{\scalebox{1.25}{$\cap$}}{\cap}{\cap}}\displaylimits}%
  \def\lesseqqgtr{\mathrel{\lessgtr}}\def\bigoplus{\mathop{\oplus}\displaylimits}%
}
\providecommand{\ssi}{if and only if} % abréviation fréquente
\AtBeginDocument{\providecommand{\wideparen}{\overparen}} % arc de cercle

% Fonctions usuelles que les modèles utilisent sans que amsmath les fournisse
\providecommand{\arsinh}{\operatorname{arsinh}}\providecommand{\arcosh}{\operatorname{arcosh}}
\providecommand{\artanh}{\operatorname{artanh}}\providecommand{\arsech}{\operatorname{arsech}}
\providecommand{\arcsch}{\operatorname{arcsch}}\providecommand{\arcoth}{\operatorname{arcoth}}
\providecommand{\arcsinh}{\operatorname{arsinh}}\providecommand{\arccosh}{\operatorname{arcosh}}
\providecommand{\arctanh}{\operatorname{artanh}}\providecommand{\sech}{\operatorname{sech}}
\providecommand{\csch}{\operatorname{csch}}\providecommand{\arcsec}{\operatorname{arcsec}}
\providecommand{\arccsc}{\operatorname{arccsc}}\providecommand{\arccot}{\operatorname{arccot}}
\providecommand{\sgn}{\operatorname{sgn}}\providecommand{\lcm}{\operatorname{lcm}}
\providecommand{\Var}{\operatorname{Var}}\providecommand{\Cov}{\operatorname{Cov}}
\providecommand{\permil}{\ensuremath{{}^{0}\!/_{00}}}\providecommand{\textperthousand}{\ensuremath{{}^{0}\!/_{00}}}

%% FIGFIX-BEGIN v5 — correctifs globaux des figures (docs/onbuch-generation/tools/figfix)
\usepackage{adjustbox}\usepackage{etoolbox}\tcbuselibrary{raster}
% toute figure TikZ/pgfplots plus large que la ligne est réduite à la largeur disponible
\BeforeBeginEnvironment{tikzpicture}{\begin{adjustbox}{max width=\linewidth}}
\AfterEndEnvironment{tikzpicture}{\end{adjustbox}}
% légendes pgfplots lisibles par-dessus les courbes
\pgfplotsset{every axis/.append style={legend style={fill=white,fill opacity=0.92,text opacity=1,draw=black!15,inner sep=3pt}}}
% étiquettes d'arêtes (syntaxe edge["..."]) avec fond blanc
\tikzset{every edge quotes/.append style={fill=white,inner sep=1.2pt}}
%% FIGFIX-END
\begin{document}

\pagegarde{Economic growth and Development}{Economics}{Upper Sixth}{Upper Sixth — Economics}{7}{4 periods}{This chapter distinguishes economic growth from economic development, explains how growth is measured and what drives or hinders it, and presents the classical growth models (Harrod–Domar, Rostow's stages) used to analyse the take-off of developing economies like Cameroon. It closes with the meaning, indicators and strategies of economic development, including human development and the GCE exam's favourite debates.
\vspace{4pt}
\begin{multicols}{2}\small
\begin{itemize}
\item By the end of this chapter I can define economic growth and distinguish it clearly from economic development.
\item By the end of this chapter I can calculate and interpret the growth rate of GDP/GNP, in nominal and real terms, and per capita.
\item By the end of this chapter I can explain the sources and factors of economic growth (labour, capital, natural resources, technology, institutions).
\item By the end of this chapter I can identify the benefits and costs (limits) of economic growth.
\item By the end of this chapter I can present the Harrod–Domar model, state its assumptions and equation, and use it to compute a warranted growth rate.
\item By the end of this chapter I can describe Rostow's five stages of growth and apply them critically to Cameroon's situation.
\end{itemize}\end{multicols}}

\begin{sommaire}
\begin{multicols}{2}
\begin{enumerate}
\item Economic Growth: Meaning, Measurement and Sources
\item The Harrod–Domar Growth Model
\item Rostow's Stages of Growth and Other Perspectives
\item Economic Development: Meaning, Indicators and Strategies
\item Integration activity
\item Methods and worked examples
\item Exercises
\item Solutions to the exercises
\item Summary sheet
\end{enumerate}
\end{multicols}
\end{sommaire}

\begin{objectifs}
\begin{itemize}
\item By the end of this chapter I can define economic growth and distinguish it clearly from economic development.
\item By the end of this chapter I can calculate and interpret the growth rate of GDP/GNP, in nominal and real terms, and per capita.
\item By the end of this chapter I can explain the sources and factors of economic growth (labour, capital, natural resources, technology, institutions).
\item By the end of this chapter I can identify the benefits and costs (limits) of economic growth.
\item By the end of this chapter I can present the Harrod–Domar model, state its assumptions and equation, and use it to compute a warranted growth rate.
\item By the end of this chapter I can describe Rostow's five stages of growth and apply them critically to Cameroon's situation.
\item By the end of this chapter I can evaluate the relevance and limitations of simple growth models for developing countries.
\item By the end of this chapter I can define economic development and explain its multidimensional character (economic, social, political, cultural).
\item By the end of this chapter I can interpret development indicators such as GDP per capita, HDI and other composite indices, and discuss their limitations.
\item By the end of this chapter I can discuss strategies and obstacles to development in Cameroon and Less Developed Countries.
\end{itemize}
\end{objectifs}

\begin{prerequis}
\begin{itemize}
\item National income accounting: GDP, GNP, NNP and their measurement (Lower Sixth).
\item The concepts of saving, investment and capital formation.
\item Population growth and demographic transition (Lower Sixth).
\item Basic macroeconomic aggregates and the circular flow of income.
\item Simple percentage and rate-of-change calculations.
\end{itemize}
\end{prerequis}

\newpage

\coursec{Economic Growth: Meaning, Measurement and Sources}

In 1986, a bag of cement in Douala cost a fraction of what it costs today, yet many Cameroonians feel poorer than their parents did. Cameroon produces more goods and services than it did thirty years ago --- its GDP has multiplied several times --- but traffic jams in Yaound\'e, overcrowded classrooms in Buea and persistent youth unemployment raise a troubling question: if the economy is ``growing'', why is life not visibly improving for everyone? Is growth the same thing as development? And what do economists like Harrod, Domar and Rostow say a country like Cameroon must do to take off? This chapter equips you to answer these questions with rigour. We begin with the most fundamental concept of all: \cle{economic growth} itself --- what it means, how we measure it, and where it comes from.

\courssub{The Meaning of Economic Growth}

\textbf{Intuition first.}\par
Imagine a small bakery in Bamenda. In 2010 it produced 1\,000 loaves of bread per week. In 2020, with a second oven and two more workers, it produces 2\,500 loaves per week. The bakery has \emph{grown}: its capacity to produce has expanded in a lasting way. Now imagine instead that in one particular week of 2015 the bakery sold 1\,200 loaves because of a funeral celebration in the neighbourhood, then fell back to 1\,000 the following week. That temporary jump is \emph{not} growth --- it is a passing fluctuation. Economic growth, at the level of a whole country, works exactly the same way: it is about a \textbf{lasting expansion of productive capacity}, not a one-off good year.

\begin{definition}[Economic growth]
\cle{Economic growth} is a \textbf{sustained increase in a country's real national output} (real GDP or real GNP) \textbf{over a long period of time}. Three words carry the whole weight of the definition:
\begin{itemize}
\item \emph{Sustained}: the increase must persist year after year, not occur once;
\item \emph{Real}: output must be measured at constant prices, stripping out inflation;
\item \emph{Long period}: growth is a long-run phenomenon, distinct from the short-run ups and downs of the business cycle.
\end{itemize}
\end{definition}

\begin{attention}[Growth is not recovery, and one good year is not growth]
A frequent GCE question asks candidates to \emph{distinguish between economic growth and an economic recovery}. A \cle{recovery} is the upward phase of the business cycle in which output returns towards its previous level after a recession --- the economy is simply re-using idle capacity (workers and machines that were unemployed). \cle{Growth}, by contrast, means that \textbf{productive capacity itself has expanded}: the economy can now produce more than it ever could before. In terms of the production possibility frontier (PPF), a recovery is a movement from a point \emph{inside} the frontier towards the frontier; growth is an \emph{outward shift} of the frontier itself. A single good harvest of cocoa after a bad year is recovery, not growth.
\end{attention}

\begin{popfigure}
\begin{center}
\begin{tikzpicture}
\begin{axis}[
width=12cm, height=6.5cm,
xlabel={Year}, ylabel={Real GDP (billion FCFA, constant prices)},
xmin=0, xmax=20, ymin=0, ymax=12,
xtick={0,5,10,15,20}, ytick={0,3,6,9,12},
legend style={at={(0.02,0.98)}, anchor=north west, font=\small},
grid=both, grid style={popline, dashed, very thin},
]
\addplot[popmuted, thick, dashed] coordinates {(0,2.5) (20,10.5)};
\addlegendentry{Long-run trend (growth)}
\addplot[popblue, very thick, smooth] coordinates {(0,2.5) (1,2.9) (2,3.4) (3,3.2) (4,3.6) (5,4.3) (6,5.0) (7,4.7) (8,5.2) (9,5.9) (10,6.6) (11,6.3) (12,6.9) (13,7.6) (14,8.4) (15,8.0) (16,8.6) (17,9.3) (18,10.0) (19,9.7) (20,10.5)};
\addlegendentry{Actual real GDP (fluctuates around trend)}
\end{axis}
\end{tikzpicture}
\end{center}
\legende{Real GDP of a stylised economy over 20 years: actual output fluctuates around a rising long-run trend. The rising trend is growth; the wiggles are the business cycle.}
\end{popfigure}

\textbf{Extensive versus intensive growth.}\par
Growth can be achieved in two fundamentally different ways. \cle{Extensive growth} occurs when output rises simply because \emph{more inputs} are used: more workers, more hectares of farmland, more machines of the same kind. A village that doubles its cocoa output by clearing twice as much forest is experiencing extensive growth. \cle{Intensive growth} occurs when output rises because \emph{each unit of input becomes more productive}: better fertiliser, better-trained farmers, improved varieties, better organisation. Intensive growth is the more sustainable path, because labour and land are ultimately limited, whereas productivity can in principle rise without limit. Modern growth in advanced economies is overwhelmingly intensive; much of Africa's historical growth has been extensive, which is why the question of productivity is so central for Cameroon.

\begin{aretenir}[Key distinction]
\begin{itemize}
\item \textbf{Extensive growth} = more output from \emph{more inputs} (quantity of labour, land, capital).
\item \textbf{Intensive growth} = more output from \emph{higher productivity} of each input (quality, technology, efficiency).
\end{itemize}
\end{aretenir}

\courssub{The Measurement of Economic Growth}

\textbf{Nominal GDP versus real GDP.}\par
Gross Domestic Product (GDP) is the market value of all final goods and services produced within a country in a year. But GDP measured at \emph{current prices} mixes together two things: genuine increases in the volume of production, and mere increases in prices. If the price of a bag of cement doubles and the number of bags produced stays the same, nominal GDP from cement doubles --- yet not a single extra house can be built. To isolate genuine growth, economists \emph{deflate} nominal GDP by a price index (usually the GDP deflator):

\[ \text{Real GDP} = \frac{\text{Nominal GDP}}{\text{GDP deflator}} \times 100 \]

\begin{propriete}[Growth rate formulas]
The rate of economic growth between year $t-1$ and year $t$ is:
\[ g = \frac{\text{Real GDP}_{t} - \text{Real GDP}_{t-1}}{\text{Real GDP}_{t-1}} \times 100 \]
The \textbf{per capita} growth rate, which tells us whether output is growing faster than the population, is approximated by:
\[ g_{\text{per capita}} \approx g_{\text{GDP}} - g_{\text{population}} \]
(The exact formula is $g_{pc} = \left(\frac{1+g_{GDP}}{1+g_{pop}} - 1\right)\times 100$, but the approximation is accepted at GCE level.) These formulas are directly examinable and must be known with their units (\%).
\end{propriete}

\begin{attention}[The classic trap: confusing nominal growth with real growth]
A rise in \textbf{nominal} GDP may be \emph{pure inflation}. If nominal GDP rises by 12\% while the price level rises by 10\%, real growth is only about 2\%. In an examination, \textbf{always deflate by the price index} before concluding that the economy has grown. Candidates who compute growth from current-price figures lose marks.
\end{attention}

\begin{exoresolu}[Computing real and per capita growth with stylised Cameroonian data]
\textbf{Statement.} A stylised economy (call it ``Kamerunia'', inspired by Cameroon) reports the following data:

\begin{center}
\begin{tabular}{|l|c|c|}
\hline
\rowcolor{popblueL} & \textbf{Year 1} & \textbf{Year 2} \\
\hline
Nominal GDP (billion FCFA) & 20\,000 & 22\,880 \\
\hline
Price index (base 100, Year 1) & 100 & 104 \\
\hline
Population (millions) & 25.0 & 25.7 \\
\hline
\end{tabular}
\end{center}

Calculate: (a) real GDP in each year; (b) the rate of economic growth; (c) the per capita growth rate. Comment.
\tcblower
\textbf{Solution.}
\textbf{(a) Real GDP (Year 1 prices):}
\begin{align*}
\text{Real GDP}_1 &= \frac{20\,000}{100}\times 100 = 20\,000 \text{ billion FCFA}\\
\text{Real GDP}_2 &= \frac{22\,880}{104}\times 100 = 22\,000 \text{ billion FCFA}
\end{align*}
\textbf{(b) Growth rate:}
\[ g = \frac{22\,000 - 20\,000}{20\,000}\times 100 = 10\% \]
Note the trap avoided: nominal GDP rose by $\frac{22\,880-20\,000}{20\,000}\times 100 = 14.4\%$, but 4 percentage points of that were pure inflation (the price index rose by 4\%).
\textbf{(c) Population growth:} $g_{pop} = \frac{25.7 - 25.0}{25.0}\times 100 = 2.8\%$. Hence:
\[ g_{\text{per capita}} \approx 10\% - 2.8\% = 7.2\% \]
\textbf{Comment.} Output per person rose by about 7.2\%: on average, citizens became materially better off. Had population grown at 10\%, per capita income would have stagnated despite impressive headline growth --- a situation not far from the experience of several fast-growing African economies.
\end{exoresolu}

\begin{methode}[Answering a ``calculate the growth rate'' question]
\begin{enumerate}
\item Check whether the GDP figures are nominal or real. If nominal, deflate each year by the price index first.
\item Apply $g = \dfrac{\text{GDP}_t - \text{GDP}_{t-1}}{\text{GDP}_{t-1}}\times 100$ using \textbf{real} figures.
\item If asked for per capita growth, compute population growth and subtract.
\item Conclude in words: state whether the average citizen is better off, and mention any caveat (inflation, population, distribution).
\end{enumerate}
\end{methode}

\courssub{The Sources of Economic Growth}

Where does growth come from? Economists decompose the growth of output into the contributions of the factors of production and of efficiency (Total Factor Productivity). The main sources are the following.

\begin{itemize}
\item \textbf{Quantity of labour.} A growing, working-age population increases the number of producers. Cameroon, with its young population, possesses this potential --- but only if jobs exist.
\item \textbf{Quality of labour (human capital).} Education, health, training and experience make each worker more productive. A trained technician at SONARA produces far more value than an untrained one. Investment in \cle{human capital} is often the most profitable investment a poor country can make.
\item \textbf{Capital accumulation.} Machines, roads, ports, electricity and factories --- \cle{physical capital} --- raise output per worker. This requires saving and investment, the central variable in the Harrod--Domar model studied in the next section.
\item \textbf{Natural resources.} Fertile land, forests, minerals and oil can fuel growth (Cameroon's oil, timber and cocoa), but resources alone guarantee nothing --- the ``resource curse'' shows that some resource-rich countries grow slowly.
\item \textbf{Technological progress.} New knowledge --- better seed varieties, mobile banking, improved refining techniques --- allows more output from the same inputs. It is the engine of \emph{intensive} growth and appears as the \cle{Total Factor Productivity (TFP)} residual in growth accounting.
\item \textbf{Entrepreneurship.} Individuals who organise the other factors, take risks and innovate --- from the Douala importer to the Silicon Mountain start-up founder in Buea --- turn potential into production.
\item \textbf{Institutions and political stability.} Secure property rights, the rule of law, low corruption, sound money (such as the stability provided by the CFA franc's peg) and peace encourage investment. Instability destroys capital and frightens investors away.
\end{itemize}

\begin{popfigure}
\begin{center}
\begin{tikzpicture}
\begin{axis}[
ybar, bar width=0.9cm,
width=11cm, height=6cm,
ylabel={Contribution to growth (percentage points)},
symbolic x coords={Labour, Capital, Technology \& efficiency},
xtick=data,
ymin=0, ymax=4,
ytick={0,1,2,3,4},
nodes near coords, nodes near coords style={font=\small},
grid=both, grid style={popline, dashed, very thin},
]
\addplot[fill=popblue, draw=popdark] coordinates {(Labour,1.2) (Capital,2.1) (Technology \& efficiency,1.7)};
\end{axis}
\end{tikzpicture}
\end{center}
\legende{Sources-of-growth decomposition for a stylised economy growing at 5\%: labour contributes 1.2 points, capital 2.1 points, and technology/efficiency 1.7 points.}
\end{popfigure}

\begin{popfigure}
\begin{center}
\begin{tikzpicture}[
box/.style={draw=popdark, fill=popblueL, rounded corners, align=center, text width=3.6cm, minimum height=1.2cm, font=\small},
arr/.style={->, very thick, poporange}
]
\node[box, align=center] (f) at (0,0){\textbf{Factors of growth}\\ labour, human capital, capital, resources, technology, institutions};
\node[box, align=center] (c) at (6.2,0){\textbf{Increase in productive capacity}\\ PPF shifts outwards};
\node[box, align=center] (o) at (12.4,0){\textbf{Rise in real output}\\ sustained growth of real GDP};
\draw[arr] (f) -- (c);
\draw[arr] (c) -- (o);
\end{tikzpicture}
\end{center}
\legende{From factors to growth: inputs and efficiency expand productive capacity, which raises real output.}
\end{popfigure}

\courssub{The Benefits and the Limits of Economic Growth}

\textbf{The benefits.}\par
Growth is ardently pursued because of what it makes possible:
\begin{itemize}
\item \textbf{Higher average incomes} and consumption of goods and services;
\item \textbf{Employment creation}, as expanding firms hire more workers;
\item \textbf{Higher fiscal revenue} for the state without raising tax rates, financing schools, hospitals and roads;
\item \textbf{Poverty reduction}: sustained growth is historically the most powerful anti-poverty force;
\item \textbf{Improved public services and infrastructure}, which in turn feed back into further growth.
\end{itemize}

\textbf{The costs and limits.}\par
Growth is not an unmixed blessing, and a good GCE essay always evaluates:
\begin{itemize}
\item \textbf{Environmental degradation}: deforestation, pollution of the Wouri estuary, greenhouse gas emissions;
\item \textbf{Depletion of non-renewable resources}: oil reserves are finite; growth based on exhausting them cannot last;
\item \textbf{Inequality}: growth may enrich a small urban elite while rural areas stagnate --- average income rises, but the median citizen sees little change;
\item \textbf{Inflationary pressure}: if demand grows faster than supply capacity, growth overheats into inflation;
\item \textbf{Jobless growth}: output may rise through capital-intensive sectors (oil extraction, automated factories) that create few jobs, leaving youth unemployment untouched;
\item \textbf{Social costs}: congestion, urban squalor and the erosion of traditional ways of life.
\end{itemize}

This is precisely why the opening question of this chapter matters: \emph{growth is not the same thing as development}. Growth is a quantitative increase in output; development, as we shall see in Section 4, is a qualitative transformation of living standards, health, education and freedoms. Growth is usually necessary for development, but it is not sufficient.

\begin{savaistu}[Did you know?]
Cameroon's experience illustrates the difference between growth and fluctuations vividly. The oil boom of the late 1970s and early 1980s produced spectacular GDP increases, but when oil prices collapsed in the mid-1980s, the economy fell into a deep crisis, showing that a commodity windfall is not the same as sustained, broad-based growth. Economists now speak of the need to ``grow with stability and equity'' rather than to grow at any price.
\end{savaistu}

\begin{exercice}[Testing your understanding]
An economy's nominal GDP rises from 8\,000 billion FCFA to 9\,240 billion FCFA, while its price index rises from 100 to 110 and its population grows by 3\%.
\begin{enumerate}
\item Compute real GDP in each year and the rate of economic growth.
\item Compute the per capita growth rate.
\item A classmate claims: ``GDP rose by 15.5\%, so living standards rose by 15.5\%.'' Identify and correct the two mistakes in this claim.
\item Give two reasons why even a genuine rise in per capita GDP might not improve the welfare of the majority of the population.
\end{enumerate}
\end{exercice}

\begin{corrige}[Exercise 1]
\textbf{1.} Real GDP$_1 = \frac{8\,000}{100}\times 100 = 8\,000$; Real GDP$_2 = \frac{9\,240}{110}\times 100 = 8\,400$ billion FCFA. Growth: $g = \frac{8\,400-8\,000}{8\,000}\times 100 = 5\%$.
\textbf{2.} Per capita growth $\approx 5\% - 3\% = 2\%$.
\textbf{3.} First mistake: 15.5\% is the \emph{nominal} increase; 10 points of it are inflation, so real growth is only 5\%. Second mistake: even real GDP growth is not per capita growth; after population growth of 3\%, income per person rose by only about 2\%.
\textbf{4.} (i) The gains may be unequally distributed (rising inequality); (ii) growth may be jobless or may impose environmental and social costs that offset higher incomes.
\end{corrige}

\begin{aretenir}[What you must remember]
\begin{itemize}
\item Economic growth = a \textbf{sustained, long-run} increase in \textbf{real} national output; a recovery or a single good year is not growth.
\item Always compute growth from \textbf{real} (deflated) GDP: $g = \frac{\text{GDP}_t-\text{GDP}_{t-1}}{\text{GDP}_{t-1}}\times 100$; per capita growth $\approx g_{GDP} - g_{pop}$.
\item Extensive growth uses more inputs; intensive growth uses inputs more productively.
\item Sources of growth: labour (quantity and human capital), capital accumulation, natural resources, technology (TFP), entrepreneurship, institutions and stability.
\item Growth brings incomes, jobs and fiscal revenue, but also risks: inequality, inflation, environmental damage, resource depletion and jobless growth.
\item Growth $\neq$ development: growth is quantitative; development is qualitative. This distinction opens the rest of the chapter.
\end{itemize}
\end{aretenir}

\coursec{The Harrod–Domar Growth Model}

In the previous section we defined economic growth, saw how it is measured through real GDP and GDP per capita, and identified its sources: the accumulation of physical capital, the growth of the labour force and technical progress. A natural question now arises: \emph{can we predict, and therefore plan, the rate at which an economy will grow?} This is precisely the question that two economists working independently — Sir Roy Harrod in England (1939) and Evsey Domar in the United States (1946) — tried to answer. Their common framework, known as the \cle{Harrod--Domar model}, became the first formal growth model of the post-Keynesian era and, for decades, the favourite planning tool of developing countries, including the early planners of Cameroon.

\courssub{Context and Purpose of the Model}

After the Great Depression of the 1930s, Keynes had shown how aggregate demand determines the level of national income \emph{in the short run}, taking the capital stock as given. But Keynes's analysis was static: it did not explain how income grows over time. Harrod and Domar asked a dynamic question: \textbf{what conditions must be satisfied for an economy to grow at a steady, balanced rate, year after year, without falling into unemployment or inflation?}

The question was especially urgent for the newly independent countries of Africa and Asia in the 1950s and 1960s. Their governments needed a simple rule of thumb to answer a very practical question: \emph{``If we want national income to grow by 6\% per year, how much must we invest?''} The Harrod--Domar model gave a strikingly simple answer, and this is why it was built into the first Five-Year Development Plans of many countries, including the early development plans of Cameroon after independence.

\courssub{Assumptions of the Model}

Like every model in economics, the Harrod--Domar model simplifies reality. Its main assumptions are:

\begin{itemize}
\item A \textbf{closed economy} with no foreign trade and no government sector (these can be added later).
\item \textbf{Saving is proportional to income}: $S = sY$, where $s$ is the constant average (and marginal) propensity to save.
\item A \textbf{fixed capital--output ratio}: technology is such that a fixed amount of capital is always needed to produce one unit of output. There is no substitution between capital and labour.
\item \textbf{No technical progress}: the productivity of capital does not change over time.
\item \textbf{Labour is abundant}: labour supply does not constrain production (a reasonable assumption for labour-surplus developing economies).
\item \textbf{Saving equals investment}: in a closed economy with no government, all saving is automatically channelled into investment, $I = S$.
\item \textbf{No depreciation}: capital stock does not wear out (or gross investment equals net investment plus replacement, keeping the ratio simple).
\end{itemize}

\courssub{Key Variables: The Saving Ratio and the Capital--Output Ratio}

Two variables drive the whole model.

\begin{definition}[The saving ratio]
The \cle{saving ratio} $s$ is the proportion of national income that is saved:
\[
s = \frac{S}{Y}.
\]
If Cameroonian households, firms and the government together save 150 billion FCFA out of a national income of 1\,000 billion FCFA, then $s = 150/1000 = 0.15$, i.e. 15\%.
\end{definition}

\begin{definition}[The capital--output ratio (ICOR)]
The \cle{capital--output ratio}, also called the \cle{incremental capital--output ratio} (ICOR) when expressed in marginal terms, measures the amount of capital needed to produce one additional unit of output:
\[
v = \frac{K}{Y} = \frac{\Delta K}{\Delta Y}.
\]
If an economy needs 3 FCFA of extra capital (machines, roads, factories) to produce 1 extra FCFA of output per year, then $v = 3$.
\end{definition}

The ICOR captures the \emph{efficiency} of investment: the lower $v$, the more productive capital is. A capital-intensive project such as a deep-sea port (e.g. Kribi Deep Seaport) or an aluminium smelter (e.g. Alucam) has a high $v$; investment in small-scale agriculture or light manufacturing usually has a lower $v$.

\courssub{Derivation of the Fundamental Equation}

The derivation is short and \textbf{examinable} — you must be able to reproduce it step by step.

Start from three building blocks:
\begin{enumerate}
\item Saving behaviour: $S = sY$.
\item Equilibrium condition (closed economy, no government): $I = S$.
\item Fixed capital--output ratio: since $v = \dfrac{\Delta K}{\Delta Y}$ and investment is the addition to the capital stock ($I = \Delta K$, ignoring depreciation), we have $I = v\,\Delta Y$.
\end{enumerate}

Combining (1), (2) and (3):
\begin{align*}
I &= S \\
v\,\Delta Y &= sY \\
\frac{\Delta Y}{Y} &= \frac{s}{v}.
\end{align*}

\begin{propriete}[The Harrod--Domar fundamental equation]
The \cle{warranted growth rate} of output — the rate at which the economy must grow for planned investment to equal planned saving — is:
\[
g = \frac{s}{v}.
\]
Growth is faster the higher the saving ratio $s$ and the lower the capital--output ratio $v$.
\end{propriete}

The intuition is beautiful in its simplicity: saving finances investment; investment becomes new capital; new capital produces new output. Income therefore grows at the rate at which saving, transformed through the productivity of capital, expands productive capacity.

\begin{popfigure}
\begin{center}
\begin{tikzpicture}[node distance=1.1cm, every node/.style={font=\small}]
\node[draw=popblue, fill=popblueL, rounded corners, minimum width=2.6cm, minimum height=0.9cm] (y) {National income $Y$};
\node[draw=poporange, fill=popgold!30, rounded corners, minimum width=2.6cm, minimum height=0.9cm, right=2.2cm of y] (s) {Saving $S = sY$};
\node[draw=poppurple, fill=poppurple!15, rounded corners, minimum width=2.6cm, minimum height=0.9cm, below=1.2cm of s] (i) {Investment $I = S$};
\node[draw=popgreen, fill=popgreenL, rounded corners, minimum width=2.6cm, minimum height=0.9cm, left=2.2cm of i] (k) {Capital stock $K$};
\draw[-latex, thick, popdark] (y) -- (s);
\draw[-latex, thick, popdark] (s) -- (i);
\draw[-latex, thick, popdark] (i) -- node[above, font=\footnotesize]{$\Delta K = I$} (k);
\draw[-latex, thick, popdark] (k) -- node[left, font=\footnotesize, align=center]{output $\Delta Y = \Delta K / v$} (y);
\end{tikzpicture}
\end{center}
\legende{The circular logic of the Harrod--Domar model: income generates saving, saving finances investment, investment enlarges the capital stock, and capital produces more output.}
\end{popfigure}

\begin{exemplebox}[Two worked numerical examples]
\textbf{Example 1.} Suppose the saving ratio is $s = 15\%$ and the capital--output ratio is $v = 3$. Then:
\[
g = \frac{s}{v} = \frac{0.15}{3} = 0.05 = 5\% \text{ per year.}
\]
The economy can grow steadily at 5\% per year.

\textbf{Example 2.} Suppose a country saves 20\% of its income and needs 4 units of capital per extra unit of output ($s = 0.20$, $v = 4$). Then:
\[
g = \frac{0.20}{4} = 0.05 = 5\% \text{ per year.}
\]
Notice that a higher saving ratio does \emph{not} guarantee faster growth if capital is used inefficiently (high $v$).
\end{exemplebox}

\begin{attention}[Common mistakes with $g = s/v$]
Three errors appear again and again in examination scripts:
\begin{itemize}
\item \textbf{Inverting the formula}: writing $g = v/s$ instead of $g = s/v$. Remember: more saving means more growth, so $s$ must be in the numerator.
\item \textbf{Mixing percentages and decimals}: $g = 15/3 = 5\%$ is correct only because 15 is a percentage; if you write $g = 0.15/3 = 0.05$, convert back to 5\%. Never compute $g = 15\%/3$ and report ``5'' without the \% sign.
\item \textbf{Using nominal instead of real magnitudes}: growth must be measured in \emph{real} (constant-price) terms. If prices rise by 8\% while nominal income rises by 10\%, real growth is only about 2\%.
\end{itemize}
\end{attention}

\courssub{Harrod's Three Growth Rates and the Knife-Edge Problem}

Harrod distinguished three growth rates, and the relationship between them reveals a deep instability at the heart of capitalist growth.

\begin{definition}[Harrod's three growth rates]
\begin{itemize}
\item The \cle{actual growth rate} $G$: the growth rate actually achieved by the economy in a given period.
\item The \cle{warranted growth rate} $G_w = s/v$: the rate at which firms' expectations are fulfilled — the rate at which the investment firms undertake is exactly validated by the growth of demand, so they are willing to keep investing.
\item The \cle{natural growth rate} $G_n$: the maximum rate allowed by the growth of the labour force and technical progress (the growth of ``full employment'' output).
\end{itemize}
\end{definition}

Steady, full-employment growth requires $G = G_w = G_n$. But Harrod showed this equality is \textbf{unstable} — the famous \cle{knife-edge} problem.

\begin{aretenir}[The knife-edge instability]
If the actual rate $G$ rises even slightly above the warranted rate $G_w$, demand outruns capacity, firms feel they have invested \emph{too little}, so they invest \emph{more} — pushing $G$ still further above $G_w$ (boom and inflation). If $G$ falls slightly below $G_w$, firms find excess capacity, conclude they invested \emph{too much}, and cut investment — pushing $G$ still lower (recession and unemployment). Instead of returning to equilibrium, any small deviation is \emph{amplified}. Moreover, even if $G = G_w$, there is no guarantee that $G_w = G_n$: the economy may grow steadily while leaving part of its growing labour force unemployed — a situation familiar to many African economies where jobs grow more slowly than the population of young job-seekers.
\end{aretenir}

\courssub{Domar's Contribution: The Dual Role of Investment}

Domar reached the same fundamental equation by a slightly different route, emphasising that investment has \textbf{two faces}:

\begin{itemize}
\item The \textbf{demand effect}: investment is a component of aggregate demand. Through the Keynesian multiplier, an increase in investment raises income: $\Delta Y_d = \frac{1}{s}\Delta I$ (since $I = sY \Rightarrow \Delta I = s\Delta Y$).
\item The \textbf{capacity effect}: investment adds to the capital stock and therefore to productive capacity: $\Delta Y_s = \frac{\Delta K}{v} = \frac{I}{v}$.
\end{itemize}

For balanced growth, the growth rate of demand must equal the growth rate of capacity. Equating the two growth rates:
\[
\frac{\Delta Y_d}{Y} = \frac{\Delta Y_s}{Y} \quad\Longrightarrow\quad \frac{\Delta I / s}{Y} = \frac{I / v}{Y}.
\]
Since $I = sY$, the right-hand side simplifies to $s/v$. The left-hand side is $\frac{\Delta I}{sY} = \frac{\Delta I}{I} \cdot \frac{I}{sY} = \frac{\Delta I}{I} \cdot \frac{sY}{sY} = \frac{\Delta I}{I}$.
Thus, for balance we require:
\[
\frac{\Delta I}{I} = \frac{s}{v} = g.
\]
Domar thus showed that investment must \emph{itself grow} at the rate $s/v$ for demand and capacity to expand together. Building the Lom Pangar dam or a road between Yaoundé and Douala creates income for workers and suppliers immediately (demand effect) and raises the economy's ability to produce electricity and move goods for decades (capacity effect); both effects must grow in step.

\courssub{Policy Implications for Developing Countries: The Financing Gap}

The model's message to planners in less developed countries (LDCs) is direct: \textbf{to grow faster, either raise the saving/investment rate $s$, or lower the capital--output ratio $v$} (use capital more efficiently).

\begin{methode}[Computing the required saving rate for a growth target]
To find the investment (and saving) effort needed to reach a target growth rate $g^*$:
\begin{enumerate}
\item Fix the target growth rate $g^*$ (e.g. 6\% per year).
\item Estimate the ICOR $v$ from past data: $v = \dfrac{I}{\Delta Y}$ over recent years (incremental capital-output ratio).
\item Compute the required saving rate: $s^* = g^* \times v$.
\item Compare $s^*$ with the actual domestic saving rate $s$.
\item The \cle{financing gap} is the difference: $(s^* - s) \times Y$. It must be filled by foreign aid, foreign direct investment or borrowing.
\end{enumerate}
\end{methode}

\begin{exoresolu}[The financing gap: a worked example]
A developing economy has a national income of $Y = 10\,000$ billion FCFA, an estimated ICOR of $v = 4$, and a domestic saving rate of $s = 12\%$. The government targets growth of $g^* = 6\%$ per year. (a) What growth rate can domestic saving alone sustain? (b) What saving rate is required? (c) What is the financing gap?
\tcblower
\textbf{(a) Growth sustainable with domestic saving:}
\[
g = \frac{s}{v} = \frac{0.12}{4} = 0.03 = 3\% \text{ per year.}
\]
Domestic resources alone deliver only 3\% growth — half the target.

\textbf{(b) Required saving rate:}
\[
s^* = g^* \times v = 0.06 \times 4 = 0.24 = 24\%.
\]

\textbf{(c) Financing gap:}
\[
(s^* - s) \times Y = (0.24 - 0.12) \times 10\,000 = 0.12 \times 10\,000 = 1\,200 \text{ billion FCFA per year.}
\]
The country must find 1\,200 billion FCFA per year from abroad (aid, FDI, loans) or raise domestic saving, otherwise the 6\% target is unattainable. This is exactly the logic used in the \emph{two-gap models} (Chenery-Strout) that justified foreign aid flows to African countries in the 1960s--1980s.
\end{exoresolu}

\begin{popfigure}
\begin{center}
\begin{tikzpicture}
\begin{axis}[
width=11cm, height=7.5cm,
xlabel={Saving rate $s$ (\%)},
ylabel={Growth rate $g$ (\%)},
xmin=0, xmax=30, ymin=0, ymax=10,
xtick={0,5,10,15,20,25,30},
ytick={0,2,4,6,8,10},
grid=both, grid style={popline!40},
legend style={at={(0.03,0.97)}, anchor=north west, font=\small},
axis lines=left, thick
]
\addplot[domain=0:30, samples=50, very thick, popblue] {x/3};
\addlegendentry{$v = 3$}
\addplot[domain=0:30, samples=50, very thick, poporange] {x/4};
\addlegendentry{$v = 4$}
\addplot[domain=0:30, samples=50, very thick, popgreen] {x/6};
\addlegendentry{$v = 6$}
\addplot[mark=*, popdark, only marks] coordinates {(15,5)};
\node[font=\footnotesize, popdark, anchor=west] at (axis cs:15.5,5) {$s{=}15\%,\, v{=}3 \Rightarrow g{=}5\%$};
\end{axis}
\end{tikzpicture}
\end{center}
\legende{The warranted growth rate $g = s/v$ as a function of the saving rate, for three values of the capital--output ratio. A lower $v$ (more efficient capital) rotates the line upward: the same saving effort yields faster growth.}
\end{popfigure}

\begin{center}
\begin{tabularx}{\linewidth}{|l|X|X|X|X|}
\hline
\rowcolor{popblueL}
\textbf{Target $g^*$} & \textbf{ICOR $v$} & \textbf{Required $s^* = g^* v$} & \textbf{Actual $s$} & \textbf{Financing gap (\% of $Y$)} \\
\hline
5\% & 4 & 20\% & 12\% & 8\% \\
\hline
6\% & 4 & 24\% & 12\% & 12\% \\
\hline
6\% & 3 & 18\% & 12\% & 6\% \\
\hline
7\% & 4 & 28\% & 12\% & 16\% \\
\hline
\end{tabularx}
\end{center}

The table shows two lessons. First, the financing gap grows quickly with the ambition of the target. Second, improving the efficiency of investment (reducing $v$ from 4 to 3) shrinks the gap dramatically — from 12\% to 6\% of national income for the same 6\% target. This is why the quality of public investment (avoiding ``white elephant'' projects) matters as much as its quantity.

\courssub{Critical Evaluation of the Harrod--Domar Model}

\begin{propriete}[Strengths and weaknesses of the Harrod--Domar model]
\textbf{Strengths:}
\begin{enumerate}
\item \textbf{Simplicity and operational power}: with only two statistics ($s$ and $v$), a planner can set consistent growth and investment targets — invaluable for countries with limited data, such as Cameroon in the 1960s.
\item \textbf{It highlights the central role of capital accumulation and saving} in the early stages of development, and it provides a rigorous justification for policies that mobilise domestic savings and for foreign aid (the financing gap).
\item \textbf{It founded modern growth theory}: the knife-edge problem forced later economists (Solow, endogenous growth theorists) to ask how flexibility of technology and factor prices could stabilise growth.
\end{enumerate}
\textbf{Weaknesses:}
\begin{enumerate}
\item \textbf{Unrealistic assumptions}: the capital--output ratio is not fixed (firms can substitute labour for capital), technical progress does occur, and saving is not automatically transformed into productive investment — in many LDCs savings leak into consumption imports, capital flight or poorly managed projects.
\item \textbf{Saving is necessary but not sufficient}: the model ignores human capital (education, health), institutions (property rights, good governance, fight against corruption), infrastructure quality and entrepreneurship, all of which determine whether investment actually raises output.
\item \textbf{Limited applicability to Cameroon today}: Cameroon is an open economy dependent on cocoa, oil and timber exports, foreign capital and the CFA franc framework; growth has often diverged from what saving rates alone would predict, because of terms-of-trade shocks, governance problems and the efficiency of public investment — factors the model simply does not contain.
\end{enumerate}
\end{propriete}

\begin{savaistu}[The model that built the Five-Year Plans]
The Harrod--Domar logic travelled the world in the 1950s and 1960s. India's Second Five-Year Plan (1956) was explicitly built on a Harrod--Domar-type calculation, and the first development plans of many African states, including Cameroon's early plans after independence, used ICOR estimates to decide how much to invest in roads, energy and agriculture. Even today, when institutions such as the World Bank estimate the ``investment needs'' of a country to reach a growth target, they are — knowingly or not — applying the formula $s^* = g \times v$.
\end{savaistu}

\begin{exercice}[Applying the model]
An economy has an ICOR of 5 and a domestic saving rate of 10\%. (a) Calculate its warranted growth rate. (b) The government targets 6\% growth: what saving rate is required, and what is the financing gap as a percentage of income? (c) If investment efficiency improves so that the ICOR falls to 3, what growth rate does the original 10\% saving rate now sustain? Comment.
\end{exercice}

\begin{corrige}[Exercise 1]
(a) $g = s/v = 0.10/5 = 2\%$ per year. (b) $s^* = g^* \times v = 0.06 \times 5 = 30\%$; financing gap $= 30\% - 10\% = 20\%$ of national income — an enormous gap, showing the target is unrealistic with $v = 5$. (c) $g = 0.10/3 \approx 3.3\%$: the same saving effort now yields growth two-thirds higher. Improving the efficiency of investment (better project selection, maintenance, governance) can be as powerful as raising saving — a key lesson for essay answers on Cameroonian growth policy.
\end{corrige}

In summary, the Harrod--Domar model reduces the mystery of growth to a powerful arithmetic: growth equals the saving rate divided by the capital--output ratio. Its simplicity is both its strength — it made planning possible — and its weakness, because real economies like Cameroon's grow through technology, human capital and institutions as much as through capital accumulation. The next section examines a more descriptive account of how economies transform over time: Rostow's stages of growth.

\coursec{Rostow's Stages of Growth and Other Perspectives}

\courssub{From Harrod–Domar to Rostow: A Historical Perspective}

In the previous section we saw that the \cle{Harrod–Domar model} focuses on a single macroeconomic ratio — the savings rate relative to the capital–output ratio — to determine whether an economy can sustain a warranted growth path. While elegant, it treats the economy as a homogeneous aggregate and says little about the \emph{structural transformation} that accompanies long-run development. 

Walt Whitman Rostow, an American economic historian and political adviser, published \emph{The Stages of Economic Growth: A Non-Communist Manifesto} in 1960. Written at the height of the Cold War, the book was explicitly intended as a liberal-democratic alternative to Marx’s historical materialism. Rostow argued that all societies, regardless of their political system, pass through a \cle{linear sequence of five stages} on their way to modern economic maturity. His framework shifted attention from a single differential equation to the \emph{historical narrative} of how leading sectors, social overhead capital, and entrepreneurial elites interact to produce a self-sustaining expansion.

\begin{savaistu}[Rostow and the Cold War]
Rostow served as National Security Adviser to President Lyndon B. Johnson (1966–1969). His “non-communist manifesto” was meant to show that developing countries could achieve prosperity through capitalist industrialisation without passing through a socialist phase. The 10\% investment threshold he identified became a benchmark for US foreign aid programmes in the 1960s.
\end{savaistu}

\courssub{The Five Stages of Economic Growth}

Rostow’s schema is best visualised as a staircase where the \cle{rate of productive investment} (gross domestic capital formation as a share of national income) rises stepwise over time.

\begin{popfigure}
\begin{tikzpicture}[scale=0.9]
  % Axes
  \draw[->, thick] (0,0) -- (13,0) node[right] {Time / Historical progression};
  \draw[->, thick] (0,0) -- (0,5.5) node[above] {Investment rate (\% of GDP)};
  
  % Stage 1: Traditional Society
  \draw[popblue, very thick] (0.5,0.8) -- (2.5,0.8);
  \draw[popblue, dashed] (2.5,0.8) -- (2.5,1.5);
  \node[popblue, anchor=south] at (1.5,1.0) {Stage 1: Traditional Society};
  \node[popblue, anchor=north, font=\small] at (1.5,0.6) {Investment $<$ 5\%};
  
  % Stage 2: Preconditions
  \draw[poporange, very thick] (2.5,1.5) -- (4.5,1.5);
  \draw[poporange, dashed] (4.5,1.5) -- (4.5,2.8);
  \node[poporange, anchor=south] at (3.5,1.7) {Stage 2: Preconditions};
  \node[poporange, anchor=north, font=\small] at (3.5,1.3) {Investment 5–10\%};
  
  % Stage 3: Take-off
  \draw[popgreen, very thick] (4.5,2.8) -- (6.5,2.8);
  \draw[popgreen, dashed] (6.5,2.8) -- (6.5,3.8);
  \node[popgreen, anchor=south] at (5.5,3.0) {Stage 3: Take-off};
  \node[popgreen, anchor=north, font=\small] at (5.5,2.6) {Investment $\approx$ 10–15\%};
  
  % Stage 4: Drive to Maturity
  \draw[poppurple, very thick] (6.5,3.8) -- (9.5,3.8);
  \draw[poppurple, dashed] (9.5,3.8) -- (9.5,4.5);
  \node[poppurple, anchor=south] at (8.0,4.0) {Stage 4: Drive to Maturity};
  \node[poppurple, anchor=north, font=\small] at (8.0,3.6) {Investment 15–20\%};
  
  % Stage 5: High Mass Consumption
  \draw[poppink, very thick] (9.5,4.5) -- (12,4.5);
  \node[poppink, anchor=south] at (10.75,4.7) {Stage 5: High Mass Consumption};
  \node[poppink, anchor=north, font=\small] at (10.75,4.3) {Investment $>$ 20\% (services dominate)};
  
  % Vertical markers
  \foreach \x/\lbl in {2.5/1$\to$2, 4.5/2$\to$3, 6.5/3$\to$4, 9.5/4$\to$5} {
    \draw[gray, thin] (\x,0) -- (\x,5);
    \node[font=\tiny, rotate=90, anchor=south] at (\x,5.1) {\lbl};
  }
  
  % Tick marks on y-axis
  \foreach \y/\lab in {0.8/5\%, 1.5/5\%, 2.8/10\%, 3.8/15\%, 4.5/20\%} {
    \draw[thin] (-0.1,\y) -- (0.1,\y);
    \node[left, font=\small] at (-0.1,\y) {\lab};
  }
\end{tikzpicture}
\legende{Rostow's five stages as a rising investment-rate staircase. The vertical jumps correspond to structural breakthroughs (e.g., take-off).}
\end{popfigure}

\begin{definition}[Stage 1: Traditional Society]
A \cle{traditional society} is one in which output per head is low and largely determined by the limits of pre-Newtonian technology. Agriculture dominates (often $>$ 75\% of employment), social organisation is hierarchical and family-based, and the investment rate is typically below 5\% of national income. Surplus, if any, is absorbed by landlords or the state for non-productive purposes (palaces, wars, rituals).
\end{definition}

\begin{definition}[Stage 2: Preconditions for Take-off]
The \cle{preconditions for take-off} emerge when external shocks (trade, colonisation, resource discovery) or internal reforms (property rights, banking, education) create a modern sector alongside the traditional one. Investment rises to 5–10\% of GDP. Infrastructure (ports, railways, power) — what Rostow calls \cle{social overhead capital} — expands. A new entrepreneurial class appears, often in a \cle{leading sector} (e.g., cotton textiles in Britain, oil in the Gulf, cocoa in Ghana).
\end{definition}

\begin{definition}[Stage 3: The Take-off]
The \cle{take-off} is the decisive interval (20–30 years) during which growth becomes self-sustaining. Three conditions must be met simultaneously:
\begin{enumerate}
  \item The rate of productive investment rises from about 5\% to \textbf{over 10\% of national income}.
  \item One or more \cle{leading sectors} grow rapidly, pulling the rest of the economy forward (e.g., railways in the USA 1843–1860, textiles in Japan 1878–1900).
  \item An institutional framework (banking, corporate law, fiscal system) emerges to mobilise domestic savings and direct them into the modern sector.
\end{enumerate}
\end{definition}

\begin{aretenir}[Rostow's 10\% Threshold]
The \textbf{10\% investment-to-GDP ratio} is the most famous quantitative benchmark in Rostow's work. It signals that the economy generates enough surplus to finance its own capital widening and deepening without relying on foreign aid or primary-commodity windfalls. In GCE essays, quote this figure precisely.
\end{aretenir}

\begin{definition}[Stage 4: Drive to Maturity]
After take-off, the economy enters the \cle{drive to maturity} (roughly 40–60 years). Technology spreads beyond the leading sectors; heavy industry, chemicals, and machinery replace the original leading sector. The investment rate stabilises at 15–20\%. The economy demonstrates the capacity to produce \emph{any} good it chooses, not merely those dictated by comparative advantage.
\end{definition}

\begin{definition}[Stage 5: Age of High Mass Consumption]
In the \cle{age of high mass consumption}, the leading sectors shift to durable consumer goods (automobiles, appliances) and services (health, education, finance). Real income per head rises sufficiently for the majority to afford more than basic subsistence. The workforce moves massively into services; the investment rate may exceed 20\% but is increasingly directed toward human capital and R\&D.
\end{definition}

\courssub{Historical Timeline of Take-off Dates}

Rostow provided estimated take-off periods for several now-developed economies. The figure below places them alongside plausible African reference points.

\begin{popfigure}
\begin{tikzpicture}[scale=0.85]
  % Timeline axis: 1700 to 1960 (260 years), 1 unit = 20 years
  \draw[thick, popdark] (0,0) -- (13,0);
  \foreach \x/\year in {0/1700, 1/1720, 2/1740, 3/1760, 4/1780, 5/1800, 6/1820, 7/1840, 8/1860, 9/1880, 10/1900, 11/1920, 12/1940, 13/1960} {
    \draw (\x,-0.1) -- (\x,0.1);
    \node[below, font=\small] at (\x,-0.3) {\year};
  }
  
  % Developed countries (blue bars)
  % Britain 1783-1802 -> x = (1783-1700)/20 = 4.15 to (1802-1700)/20 = 5.1
  \draw[popblue, line width=4pt] (4.15,0.5) -- (5.1,0.5);
  \node[popblue, left, font=\small] at (4.15,0.5) {Britain (1783–1802)};
  
  % USA 1843-1860 -> (1843-1700)/20 = 7.15 to (1860-1700)/20 = 8.0
  \draw[popblue, line width=4pt] (7.15,1.0) -- (8.0,1.0);
  \node[popblue, left, font=\small] at (7.15,1.0) {USA (1843–1860)};
  
  % Germany 1850-1873 -> 7.5 to 8.65
  \draw[popblue, line width=4pt] (7.5,1.5) -- (8.65,1.5);
  \node[popblue, left, font=\small] at (7.5,1.5) {Germany (1850–1873)};
  
  % Japan 1878-1900 -> 8.9 to 10.0
  \draw[popblue, line width=4pt] (8.9,2.0) -- (10.0,2.0);
  \node[popblue, left, font=\small] at (8.9,2.0) {Japan (1878–1900)};
  
  % Russia 1890-1914 -> 9.5 to 10.7
  \draw[popblue, line width=4pt] (9.5,2.5) -- (10.7,2.5);
  \node[popblue, left, font=\small] at (9.5,2.5) {Russia (1890–1914)};
  
  % Asian Tigers (green bars)
  % South Korea 1953-1965 -> 12.65 to 13.25 (but 1965 > 1960, so cap at 13)
  \draw[popgreen, line width=4pt] (12.65,3.0) -- (13,3.0);
  \node[popgreen, left, font=\small] at (12.65,3.0) {S. Korea (1953–1965)};
  
  % Taiwan 1952-1964 -> 12.6 to 13.2
  \draw[popgreen, line width=4pt] (12.6,3.5) -- (13,3.5);
  \node[popgreen, left, font=\small] at (12.6,3.5) {Taiwan (1952–1964)};
  
  % African reference points (orange bars - tentative)
  % Mauritius ~1970s (beyond 1960, but indicate with arrow)
  \draw[poporange, line width=4pt, dashed] (13,4.0) -- (13.5,4.0);
  \node[poporange, left, font=\small] at (13,4.0) {Mauritius (1970s$\to$)};
  
  % Botswana ~1970s
  \draw[poporange, line width=4pt, dashed] (13,4.5) -- (13.5,4.5);
  \node[poporange, left, font=\small] at (13,4.5) {Botswana (1970s$\to$)};
  
  % Cameroon ? (ongoing preconditions)
  \draw[poporange, line width=4pt, dashed] (10,5.0) -- (13,5.0);
  \node[poporange, left, font=\small] at (10,5.0) {Cameroon (pre-conditions?)};
\end{tikzpicture}

\begin{center}
\begin{tabular}{ll}
\textcolor{popblue}{\rule{1em}{0.5em}} & Early industrialisers (Rostow's dates) \\
\textcolor{popgreen}{\rule{1em}{0.5em}} & Asian Tigers (post-WWII) \\
\textcolor{poporange}{\rule{1em}{0.5em}} & African economies (tentative, debated)
\end{tabular}
\end{center}
\legende{Historical take-off dates per Rostow (blue), Asian Tigers (green), and tentative African cases (orange, dashed). Cameroon is widely considered still in the preconditions stage.}
\end{popfigure}

\courssub{Applying Rostow's Model to Cameroon and African Economies}

\begin{experience}[Classifying Cameroon in Rostow's Schema]
\textbf{Step 1: Structure of output (World Bank, 2022 estimates).} Agriculture $\approx$ 16\% of GDP but $\approx$ 43\% of employment; industry (including oil, mining, construction) $\approx$ 27\%; services $\approx$ 57\%. The large employment share in low-productivity agriculture is a hallmark of Stage 2.

\textbf{Step 2: Investment rate.} Gross capital formation averaged 20–22\% of GDP in the 2010s (World Bank), but \emph{productive} investment (excluding oil exploration and public works with low multiplier) is estimated closer to 10–12\%. This hovers around the take-off threshold.

\textbf{Step 3: Leading sectors.} 
\begin{itemize}
  \item \textbf{Oil} (since 1977): capital-intensive, enclave, limited linkages; subject to Dutch disease effects.
  \item \textbf{Agriculture} (cocoa, coffee, cotton, bananas): traditional exports, vulnerable to terms-of-trade shocks.
  \item \textbf{Timber}: raw log exports dominate; local processing rising slowly.
  \item \textbf{Emerging}: telecommunications, mobile money (e.g., MTN MoMo, Orange Money), construction.
\end{itemize}
No single manufacturing sector has yet achieved the \emph{rapid, self-propelling growth} Rostow requires for take-off.

\textbf{Step 4: Institutional framework.} The OHADA legal system, BEAC monetary stability (CFA franc), and recent improvements in port infrastructure (Kribi deep-sea port) are preconditions being put in place.

\textbf{Conclusion:} Cameroon sits in \textbf{late Stage 2 (preconditions)} with pockets of Stage 3 dynamics in urban services and telecoms, but has not yet achieved a broad-based industrial take-off.
\end{experience}

\begin{methode}[How to Classify a Country in Rostow's Stages (Exam Technique)]
\begin{enumerate}
  \item \textbf{Check the investment/GDP ratio} (World Bank data). $<$ 5\% $\to$ Stage 1; 5–10\% $\to$ Stage 2; $>$ 10\% sustained $\to$ Stage 3+.
  \item \textbf{Identify the leading sector(s).} Is there a manufacturing or modern service sector growing at $>$ 10\% annually and pulling up productivity economy-wide?
  \item \textbf{Examine employment structure.} Declining agricultural share ($<$ 40\%) and rising manufacturing/services share signal Stage 3/4.
  \item \textbf{Assess institutional maturity.} Domestic savings mobilisation, financial depth (M2/GDP $>$ 30\%), and rule of law.
  \item \textbf{State the stage clearly} and justify with \emph{two or three} quantitative indicators. Avoid vague statements like “it is developing”.
\end{enumerate}
\end{methode}

\courssub{Critical Evaluation of Rostow's Stages}

\begin{attention}[Common Mistake: The Automatic Ladder]
Many candidates write: “Rostow says all countries must pass through the five stages in order.” \textbf{This is a misreading.} Rostow described a \emph{historical pattern} observed in the now-developed world; he did not prove it is a universal law. The GCE rewards answers that distinguish between \emph{description} and \emph{prescription}, and that highlight where African experience diverges.
\end{attention}

\begin{propriete}[Main Critiques of Rostow's Model]
\begin{enumerate}
  \item \textbf{Linear and Eurocentric.} The model is based on the history of Western Europe and North America. It assumes a single path, ignoring that late industrialisers (Japan, South Korea, China) used state-led strategies, export orientation, and technology transfer — shortcuts not available to 18th-century Britain.
  \item \textbf{Neglect of colonial history and external dependence.} Rostow treats the “preconditions” as internally generated. For Africa, colonial extraction (railways built for export, not integration), the imposition of cash-crop monocultures, and the terms-of-trade deterioration (Prebisch–Singer thesis) created structural distortions that persist.
  \item \textbf{Saving is not the only constraint.} Harrod–Domar and Rostow both emphasise the savings gap. In reality, developing countries often face \emph{foreign exchange gaps}, \emph{fiscal gaps}, \emph{absorptive capacity limits} (lack of skilled labour, management), and \emph{governance failures}.
  \item \textbf{The “leading sector” concept is vague.} In Cameroon, oil raised the investment rate but created few linkages (Dutch disease). A leading sector must have \emph{strong backward and forward linkages} (Hirschman) to trigger take-off.
  \item \textbf{Stage 5 is not the end of history.} Deindustrialisation, inequality, climate change, and the digital revolution challenge the idea that high mass consumption is a stable equilibrium.
\end{enumerate}
\end{propriete}

\begin{exoresolu}[Assess the relevance of Rostow's stages of growth model for explaining the development experience of Cameroon.]
\textbf{Statement:} Assess the relevance of Rostow's stages of growth model for explaining the development experience of Cameroon.
\tcblower
\textbf{Plan:}
\begin{enumerate}
  \item \textbf{Introduction:} Define Rostow's model (5 stages, 10\% threshold, leading sector). State thesis: useful as a heuristic checklist but limited as a predictive theory for Cameroon.
  \item \textbf{Relevance (Where it fits):}
  \begin{itemize}
    \item Cameroon's investment rate (20–22\% GDP) meets the quantitative take-off threshold.
    \item Infrastructure push (Kribi port, highways, dams) mirrors “social overhead capital” in Stage 2.
    \item Emerging sectors (telecoms, fintech) show Stage 3 dynamism in services.
  \end{itemize}
  \item \textbf{Limits (Where it fails):}
  \begin{itemize}
    \item \textbf{No manufacturing take-off:} Oil is an enclave; agriculture remains low-productivity. The “leading sector” lacks Hirschman's linkages.
    \item \textbf{Structural dualism persists:} 43\% employment in agriculture vs. 16\% GDP — not the smooth transition Rostow depicts.
    \item \textbf{External dependence:} CFA franc zone, commodity exports, debt service constrain policy autonomy — absent from Rostow's closed-economy logic.
    \item \textbf{Institutional gaps:} Corruption, weak property rights, electricity shortages reduce absorptive capacity.
  \end{itemize}
  \item \textbf{Alternative perspectives:} Mention Hirschman's \emph{unbalanced growth} (deliberate imbalances to induce investment) and the \emph{big push} (coordinated industrialisation) as more policy-relevant for Cameroon.
  \item \textbf{Conclusion:} Rostow provides a vocabulary (take-off, preconditions) and a benchmark (10\% investment), but Cameroon's path is better understood through structuralist and institutional lenses.
\end{enumerate}
\end{exoresolu}

\courssub{Other Perspectives: Balanced vs. Unbalanced Growth and the Big Push}

\begin{savaistu}[Beyond Rostow: Three Classic Debates]
While Rostow focuses on \emph{stages}, three contemporaneous theories focus on \emph{how} to industrialise:
\end{savaistu}

\begin{definition}[Balanced Growth (Ragnar Nurkse, 1953)]
\cle{Balanced growth} argues that the main obstacle in poor countries is the \emph{small size of the market}. Since demand is limited, no single firm can profitably invest at a large scale. The solution: simultaneous investment across a wide range of industries so that workers in each new factory become consumers for the others, expanding the market \emph{endogenously}. \textbf{Policy message:} broad-based industrial planning, import substitution across the board.
\end{definition}

\begin{definition}[Unbalanced Growth (Albert Hirschman, 1958)]
\cle{Unbalanced growth} counters that poor countries lack not only capital but also \emph{decision-making inputs} (entrepreneurs, managers, planners). Deliberately creating \emph{imbalances} (e.g., building a steel mill without enough coal) generates pressure to fill the gaps, inducing private investment where it has the highest \cle{linkage effects} (backward: demand for inputs; forward: supply for downstream industries). \textbf{Policy message:} target strategic sectors with strong linkages (the “leading sectors” Rostow mentions, but chosen deliberately).
\end{definition}

\begin{definition}[The Big Push (Paul Rosenstein-Rodan, 1943; formalised by Murphy, Shleifer, Vishny 1989)]
The \cle{big push} theory shows that when there are \emph{complementarities} across sectors (e.g., a shoe factory needs a leather industry, which needs cattle ranching), no single firm will invest alone because its profitability depends on others investing simultaneously. The state must coordinate a \emph{large-scale, simultaneous investment programme} to jump the economy to a high-level equilibrium. \textbf{Policy message:} industrial parks, special economic zones, coordinated infrastructure–industry packages.
\end{definition}

\begin{popfigure}
\begin{center}
\begin{tabularx}{\linewidth}{|l|X|X|X|}
\hline
\rowcolor{popblueL}
\textbf{Dimension} & \textbf{Harrod–Domar} & \textbf{Rostow} & \textbf{Balanced / Unbalanced / Big Push} \\
\hline
\textbf{Nature} & Formal growth model (differential equation) & Historical-stylised stages & Development strategy prescriptions \\
\hline
\textbf{Key variable} & Savings rate ($s$) and capital–output ratio ($v$) & Investment rate crossing 10\%; leading sector & Market size (Nurkse), linkages (Hirschman), complementarities (Big Push) \\
\hline
\textbf{Policy message} & Raise savings (domestic + aid) to reach warranted growth & Create preconditions; identify leading sector; wait for take-off & Nurkse: invest everywhere. Hirschman: invest in high-linkage sectors. Big Push: coordinate massive simultaneous investment. \\
\hline
\textbf{Main limit} & Knife-edge instability; ignores structure & Linear, Eurocentric; ignores colonial legacy & Nurkse: unrealistic resource needs. Hirschman: requires strong state capacity. Big Push: coordination failures, corruption risk. \\
\hline
\end{tabularx}
\end{center}
\legende{Comparison of Harrod–Domar, Rostow, and classic development strategy debates.}
\end{popfigure}

\begin{exercice}[Essay Practice]
\textbf{“Rostow’s stages of growth model is more relevant to the historical experience of Western Europe than to the contemporary challenges of African economies.” Discuss. (25 marks)}
\end{exercice}

\begin{corrige}[Essay Practice]
\textbf{Introduction:} Define Rostow's five stages, the 10\% investment threshold, and the leading-sector concept. Acknowledge the model's origin (1960, Cold War context).

\textbf{Body — Relevance to Western Europe:}
\begin{itemize}
  \item Fits the historical sequence: Britain (textiles $\to$ railways $\to$ steel $\to$ chemicals $\to$ services).
  \item Captures the role of social overhead capital (canals, railways) and institutional evolution (Bank of England, joint-stock companies).
  \item The 10\% benchmark aligns with historical national accounts (e.g., UK investment/GDP crossed 10\% c. 1780s).
\end{itemize}

\textbf{Body — Limits for Africa (Cameroon case):}
\begin{itemize}
  \item \textbf{Colonial distortion:} Infrastructure built for extraction (Douala–Nkongsamba railway for cocoa), not integrated development.
  \item \textbf{Enclave leading sectors:} Oil in Cameroon, diamonds in Botswana, copper in Zambia — capital-intensive, few linkages, volatile revenues.
  \item \textbf{Premature deindustrialisation:} Many African economies see manufacturing share peak at 10–15\% (vs. 30\%+ in East Asia), skipping Stage 4.
  \item \textbf{Institutional fragility:} Rostow assumes a functional state; weak property rights and corruption raise the effective investment threshold well above 10\%.
  \item \textbf{Demographic pressure:} Population growth 2.5–3\% requires much higher investment just to keep capital per worker constant (Harrod–Domar: $g = s/v - n$).
\end{itemize}

\textbf{Body — Alternative frameworks more suited to Africa:}
\begin{itemize}
  \item \textbf{Structuralist (Prebisch, Singer):} Focus on terms of trade, commodity dependence, need for industrial policy.
  \item \textbf{Institutionalist (Acemoglu, Robinson):} Inclusive vs. extractive institutions as the fundamental cause.
  \item \textbf{New Structural Economics (Lin):} Comparative advantage following (CAF) strategy — upgrade step by step, not big push.
\end{itemize}

\textbf{Conclusion:} Rostow remains a useful \emph{organising framework} (vocabulary, benchmarks) but a poor \emph{predictive theory} for Africa. Policy should draw on Hirschman's linkage analysis and the big push's coordination insight, adapted to local institutional reality.
\end{corrige}

\begin{aretenir}[Key Takeaways for the GCE Exam]
\begin{itemize}
  \item Know the \textbf{five stages} and the \textbf{10\% investment threshold} by heart.
  \item Be able to \textbf{classify a country} (Cameroon, Nigeria, South Korea) using \emph{three observable criteria}: investment rate, leading sector, employment structure.
  \item \textbf{Critique} must cover: linearity, Eurocentrism, neglect of colonialism, enclave sectors, institutional prerequisites.
  \item Contrast with \textbf{Nurkse (balanced), Hirschman (unbalanced), Rosenstein-Rodan (big push)} — know the \emph{core idea} and \emph{policy implication} of each.
  \item Essay phrase bank: “heuristic device”, “stylised facts”, “linkage effects”, “absorptive capacity”, “Dutch disease”, “premature deindustrialisation”.
\end{itemize}
\end{aretenir}

\coursec{Economic Development: Meaning, Indicators and Strategies}

In the previous section we studied Rostow's stages of growth and other models which try to explain how economies grow over time. Yet a moment's reflection on the Cameroonian experience raises a troubling question: Cameroon has recorded positive GDP growth in most recent decades, oil has been extracted and exported since the late 1970s, and yet millions of Cameroonians still lack clean water, good schools and reliable health care. If the economy is ``growing'', why are so many people not visibly better off? The answer is that \cle{economic growth} and \cle{economic development} are not the same thing. Growth is only one ingredient of development. This section defines economic development, shows how it is measured, describes the typical features of Less Developed Countries (LDCs), analyses the obstacles Cameroon faces, and reviews the main strategies proposed to achieve development.

\courssub{The Meaning of Economic Development}

\textbf{Intuition first.} Imagine two villages. In village A, a foreign company opens a large plantation; output doubles, GDP rises, but the profits are repatriated abroad, the workers remain on subsistence wages, and the village still has no school, no clinic and no paved road. In village B, output rises more modestly, but the extra income is taxed and used to build a health centre, train teachers, electrify homes and diversify activities away from a single crop. Which village is ``developing''? Clearly village B. Development is about \emph{people's lives changing for the better}, not merely about statistics rising.

\begin{definition}[Economic development]
\cle{Economic development} is a sustained improvement in the material and non-material well-being of the population, accompanied by a structural transformation of the economy. It combines economic growth with improvements in health, education, equity, political freedoms and institutional quality, and with a shift of labour and output from low-productivity traditional activities towards higher-productivity modern ones.
\end{definition}

Development is therefore \cle{multidimensional}. It has at least four components:
\begin{itemize}
\item \textbf{Economic growth:} a sustained rise in real output and income per head (the quantitative base).
\item \textbf{Structural change:} transformation of the economy — industrialisation, urbanisation, modernisation of agriculture, formalisation of activities.
\item \textbf{Social progress:} better health, longer life expectancy, higher literacy and schooling, reduced poverty and inequality.
\item \textbf{Institutional and human freedoms:} rule of law, security of property, participation, and the expansion of people's real choices (what the economist Amartya Sen famously called ``development as freedom'').
\end{itemize}

\begin{aretenir}[Growth versus development: the key distinction]
\cle{Growth is quantitative}: it is measured by the percentage change in real GDP. \cle{Development is qualitative and multidimensional}: it includes growth \emph{plus} structural, social and institutional change. Growth is therefore a \emph{necessary but not sufficient} condition for development. Development is the larger concept; growth is a subset of it.
\end{aretenir}

The relationship can be shown visually: development is the big circle, growth only one part of it.

\begin{popfigure}
\begin{center}
\begin{tikzpicture}
\draw[fill=popblueL, draw=popblue, very thick] (0,0) circle (3.1);
\node[text=popdark, font=\bfseries] at (0,2.55) {ECONOMIC DEVELOPMENT};
\draw[fill=popgreenL, draw=popgreen, very thick] (-1.35,-0.7) circle (1.15);
\node[text=popdark, align=center, font=\small] at (-1.35,-0.7) {Economic\\growth};
\node[text=popdark, align=center, font=\small] at (1.35,0.15) {Structural\\change};
\node[text=popdark, align=center, font=\small] at (1.35,-1.35) {Health \&\\education};
\node[text=popdark, align=center, font=\small] at (-0.1,-2.35) {Equity \& freedoms};
\end{tikzpicture}
\end{center}
\legende{Development as the wider concept: growth is a subset of development, alongside structural change, social progress and expanded freedoms.}
\end{popfigure}

\textbf{Growth without development.} History offers many cases where GDP rose while poverty deepened. Oil- and mineral-producing economies illustrate this well: extraction enclaves generate large export revenues, but if these are captured by a small elite, spent on imports of luxury goods, or repatriated by foreign firms, the majority of the population sees no improvement in schools, clinics or jobs. Economists call this pattern \cle{growth without development}, and the broader paradox of resource-rich countries remaining poor is known as the \cle{resource curse}. The lesson for exam essays is clear: a rising GDP figure, by itself, proves nothing about development.

\courssub{Characteristics of Less Developed Countries}

Although LDCs differ greatly, they share a cluster of common features. The table below summarises them with Cameroonian illustrations.

\begin{center}
\begin{tabularx}{\linewidth}{|l|X|X|}
\hline
\rowcolor{popblueL} \textbf{Characteristic} & \textbf{Explanation} & \textbf{Cameroonian illustration} \\
\hline
Low income per capita & GDP per head is low, limiting consumption, saving and investment. & Cameroon is classified as a lower-middle-income country; average incomes remain far below those of industrialised economies. \\
\hline
Dualism & A modern, monetised sector coexists with a large traditional subsistence sector. & Modern firms in Douala and Yaound\'e coexist with subsistence farming in rural areas. \\
\hline
Large agricultural and informal sectors & A big share of the labour force works in low-productivity agriculture and unregistered activities. & A large share of Cameroonians earn their living from smallholder farming (cocoa, coffee, food crops) and informal trade, artisanship and transport. \\
\hline
Rapid population growth & High birth rates create a young, dependent population. & A very large proportion of the population is under 25, pressing on schools and the labour market. \\
\hline
Unemployment and underemployment & Many workers are idle or work few hours at low productivity. & Visible underemployment among urban youth: hawking, motorcycle taxis (\emph{bendskins}), casual labour. \\
\hline
External dependence & Reliance on exports of a few raw materials, on imported manufactures, foreign capital and aid. & Exports dominated by unprocessed cocoa, oil, timber and cotton; imports of machinery, fuel products and manufactured goods. \\
\hline
\end{tabularx}
\end{center}

\begin{savaistu}[Did you know?]
Economists speak of the \emph{informal sector} as the set of small, unregistered, untaxed activities. In many African cities it provides the majority of jobs. It is a survival strategy for millions, but because it escapes taxation and regulation, it also limits the state's capacity to finance public services — one of the vicious circles of underdevelopment.
\end{savaistu}

\courssub{Indicators of Development}

\textbf{GDP/GNP per capita.} The oldest and simplest indicator is \cle{GDP per capita}: GDP divided by population. It is easy to compute and compare, and it correlates broadly with living standards. But it suffers from serious limitations:
\begin{itemize}
\item \textbf{Distribution:} it is an average; a high mean can hide extreme inequality (a few rich oil exporters, a poor majority).
\item \textbf{Informal and subsistence output:} much production in LDCs never passes through markets and is poorly measured, so GDP understates true activity.
\item \textbf{Non-market welfare:} it ignores health, education, leisure, political freedom and environmental quality.
\item \textbf{Composition of output:} GDP counts weapons and pollution-clean-up the same as schools and hospitals.
\item \textbf{Exchange-rate problems:} converting GDP into a common currency (US dollars) can distort comparisons; purchasing-power-parity (PPP) conversions partly correct this.
\end{itemize}

\begin{attention}[Common mistake]
Never use GDP per capita \emph{alone} as proof that a country is ``developed''. In an essay, always mention at least two limitations — income distribution, the informal sector, non-market output — and complement GDP with a broader indicator such as the HDI.
\end{attention}

\textbf{The Human Development Index (HDI).} To overcome these weaknesses, the United Nations Development Programme (UNDP) created the \cle{Human Development Index} in 1990.

\begin{definition}[Human Development Index]
The \cle{HDI} is a composite index, published annually by the UNDP, combining three dimensions of human development: (1) a long and healthy life, measured by \textbf{life expectancy at birth}; (2) knowledge, measured by \textbf{mean and expected years of schooling}; and (3) a decent standard of living, measured by \textbf{GNI per capita at PPP}. Each dimension is converted into an index between 0 and 1, and the HDI is their geometric mean.
\end{definition}

The calculation principle is simple: for each dimension, a country's achievement is placed between a minimum and a maximum goalpost to obtain an index between 0 and 1; the three indices are then combined (geometric mean). The HDI therefore ranges from 0 (worst) to 1 (best). Countries are grouped into very high, high, medium and low human development. Cameroon falls in the \emph{medium} human development group, ranked in the lower half of the world table — and, importantly, its HDI rank is typically \emph{worse} than its income rank, a sign that income is not being fully converted into health and education.

\textbf{Other indicators.} Development analysts also use:
\begin{itemize}
\item \textbf{Poverty incidence:} the percentage of the population living below a poverty line (e.g.\ the international extreme-poverty line).
\item \textbf{Adult literacy rate} and school enrolment ratios.
\item \textbf{Life expectancy at birth} and infant/maternal mortality rates.
\item \textbf{Inequality:} the \cle{Gini coefficient}, ranging from 0 (perfect equality) to 1 (one person receives all income); high inequality signals that growth is not inclusive.
\end{itemize}

\begin{popfigure}
\begin{center}
\begin{tikzpicture}
\begin{axis}[
ybar, bar width=0.55cm,
width=12.5cm, height=6.2cm,
ymin=0, ymax=100,
ylabel={Index (GDP per capita scaled; HDI $\times$ 100)},
symbolic x coords={Norway, S.\ Africa, Cameroon, Chad},
xtick=data,
nodes near coords, every node near coord/.append style={font=\scriptsize},
legend style={at={(0.5,-0.18)}, anchor=north, legend columns=2},
]
\addplot[fill=popblue] coordinates {(Norway,100) (S.\ Africa,22) (Cameroon,6) (Chad,3)};
\addplot[fill=poporange] coordinates {(Norway,96) (S.\ Africa,72) (Cameroon,59) (Chad,40)};
\legend{GDP per capita (Norway $=100$), HDI $\times$ 100}
\end{axis}
\end{tikzpicture}
\end{center}
\legende{Illustrative comparison (orders of magnitude): income differences between countries are far wider than HDI differences, because health and education partly compensate for low income. Cameroon's HDI exceeds its relative income position — yet remains in the medium-development band.}
\end{popfigure}

\courssub{Obstacles to Development in Cameroon and LDCs}

Why do LDCs find development so difficult? The obstacles can be classified into four groups — a classification you should reuse in essays.

\begin{itemize}
\item \textbf{Economic obstacles:} low incomes $\Rightarrow$ low saving $\Rightarrow$ low capital formation $\Rightarrow$ low productivity $\Rightarrow$ low incomes (the \cle{vicious circle of poverty}); poor infrastructure (roads, electricity, ports) raising production costs; narrow industrial base; dependence on a few primary exports whose prices fluctuate.
\item \textbf{Social obstacles:} rapid population growth diluting capital accumulation and overloading schools and hospitals; low levels of education and skills; poor health reducing labour productivity; rapid, unplanned urbanisation.
\item \textbf{Political and institutional obstacles:} corruption and mismanagement of public funds; weak rule of law and insecure property rights discouraging investment; administrative inefficiency; \cle{brain drain} — the emigration of trained doctors, nurses, engineers and lecturers to richer countries.
\item \textbf{External obstacles:} heavy debt service absorbing scarce public revenue; unfavourable \cle{terms of trade} (prices of exported raw materials tend to stagnate or fluctuate while prices of imported manufactures rise); dependence on foreign aid and capital with attached conditions.
\end{itemize}

\begin{methode}[Structuring an essay on ``obstacles to development in Cameroon'']
\begin{enumerate}
\item \textbf{Introduction:} define development; state that Cameroon, despite growth, faces persistent obstacles; announce the plan.
\item \textbf{Classify the obstacles} into four families: economic (low saving, poor infrastructure), social (population pressure, weak human capital), political/institutional (corruption, brain drain), external (debt, terms of trade).
\item \textbf{Illustrate each obstacle} with one concrete Cameroonian example (e.g.\ electricity shortages raising firms' costs; emigration of health workers).
\item \textbf{Show the interconnections:} obstacles reinforce one another (vicious circles).
\item \textbf{Conclusion:} evaluation — which obstacles are most binding, and why overcoming them requires combined strategies.
\end{enumerate}
\end{methode}

\courssub{Strategies of Development}

There is no single recipe; the debate concerns which engine should drive development.

\textbf{1. Industrialisation versus agricultural development.} The classical view (Lewis, and the models seen earlier) sees \cle{industrialisation} as the engine: industry offers higher productivity, economies of scale and dynamic demand. But critics argue that neglecting \cle{agriculture} — where most poor people work — is a mistake: raising farm productivity releases labour, raises rural incomes, supplies food and raw materials, and creates demand for industrial goods. The modern consensus favours a \emph{balanced} approach: agricultural modernisation first or alongside industry.

\textbf{2. Import substitution versus export promotion.} \cle{Import-substitution industrialisation (ISI)} protects infant industries behind tariffs to replace imports; it was tried widely in Africa and Latin America but often produced inefficient, high-cost industries dependent on continued protection. \cle{Export promotion} encourages industries to compete on world markets (the path of the Asian ``tigers''); it imposes efficiency but requires competitiveness, infrastructure and stable macroeconomic policies. Many countries now combine elements of both, and regional initiatives such as the African Continental Free Trade Area (AfCFTA) aim to widen markets for African exporters.

\textbf{3. Human capital investment.} Spending on \cle{education and health} is not consumption but investment: healthier, better-educated workers are more productive, innovate more and adapt faster. Human capital is now seen as a precondition for all other strategies — and as the best antidote to brain drain when combined with opportunities at home.

\textbf{4. The state versus the market.} Should development be state-led (planning, public enterprises, industrial policy) or market-led (liberalisation, privatisation, price signals)? Experience suggests a pragmatic mix: markets are powerful allocators, but the state must provide public goods (infrastructure, education, health, justice), correct market failures and protect the vulnerable. Weak or predatory states fail; absent states fail too.

\textbf{5. Foreign aid and foreign direct investment.} \cle{Aid} can finance infrastructure and social programmes but may create dependency, be mismanaged, or come with conditions. \cle{Foreign direct investment (FDI)} brings capital, technology and market access, but may repatriate profits, create enclaves and exploit resources without local linkages. The evaluation in an essay: aid and FDI help only if institutions channel them into productive, employment-creating uses.

\courssub{Sustainable Development and Inclusive Growth}

Recent thinking — increasingly present in GCE themes — adds two requirements. \cle{Sustainable development} is development that meets the needs of the present without compromising the ability of future generations to meet theirs: growth must not destroy forests, soils, fisheries and the climate. \cle{Inclusive growth} is growth whose benefits are widely shared, reducing poverty and inequality rather than concentrating gains. For Cameroon, this means valuing its forests and biodiversity, investing in renewable energy, and ensuring that growth creates jobs for its huge youth population.

\textbf{Cameroon focus.} Cameroon's long-term vision (often summarised as the ambition of ``emergence'') sets objectives of industrialisation, infrastructure development, and poverty reduction, monitored through indicators such as the HDI. Reading Cameroon's position in the HDI table — medium human development, ranked below what its income alone would predict — is a perfect application for the exam: it shows concretely the gap between growth and development, and why structural and social policies matter.

\begin{exoresolu}[Worked exercise: growth or development?]
\textbf{Statement.} Country X's real GDP grows by $6\%$ per year for ten years thanks to oil exports. Over the same period, life expectancy stagnates, the Gini coefficient rises from 0.42 to 0.55, and the share of manufacturing in GDP falls. (a) Has country X experienced economic growth? (b) Has it experienced economic development? Justify. (c) Which indicators would you consult to complete the diagnosis?
\tcblower
\textbf{Solution.}
(a) \textbf{Yes}, country X has experienced economic growth: by definition, growth is a sustained rise in real GDP, and here real GDP rises by $6\%$ per year for a decade — a strong performance.

(b) \textbf{No — or at most very partially.} Development requires growth \emph{plus} structural and social progress. Here the three accompanying signals are negative: life expectancy stagnates (no health progress), the Gini coefficient rises sharply from 0.42 to 0.55 (growth is accompanied by increasing inequality, so its benefits are concentrated), and the falling share of manufacturing shows \emph{no} structural transformation — rather a reinforcement of the primary-export enclave. This is a textbook case of \cle{growth without development}, typical of the resource-curse pattern.

(c) To complete the diagnosis one would consult: the \textbf{HDI} (to combine income, health and education), the \textbf{poverty incidence} (has the number of poor fallen?), the \textbf{adult literacy rate} and school enrolment, \textbf{infant mortality}, and the \textbf{unemployment/underemployment rate}. Only if these improve alongside GDP could one speak of development.
\end{exoresolu}

\begin{exercice}[Essay practice]
``Economic growth is a necessary but not a sufficient condition for economic development.'' Discuss this statement with reference to Cameroon or any African country you have studied. (Structure: definitions; why growth helps development; why it is not enough — limitations of GDP, distribution, structural change; indicators you would use; conclusion with evaluation.)
\end{exercice}

\begin{exercice}[Data response]
The table below gives hypothetical data for two countries.\par
\begin{center}
\begin{tabular}{|l|c|c|}
\hline
\rowcolor{popblueL} \textbf{Indicator} & \textbf{Country A} & \textbf{Country B} \\
\hline
GDP per capita (USD, PPP) & 4\,000 & 3\,200 \\
\hline
Life expectancy (years) & 55 & 68 \\
\hline
Mean years of schooling & 5 & 9 \\
\hline
Gini coefficient & 0.58 & 0.36 \\
\hline
\end{tabular}
\end{center}
Which country is more ``developed''? Justify your answer carefully, explaining why GDP per capita alone is misleading here.
\end{exercice}

\begin{corrige}[Exercise 2]
Although country A has the higher GDP per capita (4\,000 against 3\,200 USD), country B is arguably more developed. Its citizens live much longer (68 against 55 years) and are far better educated (9 against 5 mean years of schooling), so its HDI would almost certainly be higher. Moreover, income in B is distributed much more equally (Gini 0.36 against 0.58), so average income in A overstates the typical person's standard of living: A's higher mean may reflect a rich minority alongside a poor majority. This illustrates precisely why GDP per capita alone is a misleading indicator of development: it ignores distribution, health and education. A complete judgement must use a composite indicator such as the HDI plus inequality and poverty measures.
\end{corrige}

\begin{aretenir}[Exam tip]
For \emph{any} essay on development, master: (1) one table of indicators — GDP per capita, HDI and its three components, life expectancy, literacy, poverty incidence, Gini; and (2) one Cameroonian example — e.g.\ growth driven by oil and primary exports coexisting with a medium HDI rank, a large informal sector and youth underemployment. With these two tools you can illustrate and evaluate any development question.
\end{aretenir}

\newpage

\coursec{Integration activity}

\coursec{Economic Growth and Development}

\courssub{Chapter Overview and Learning Objectives}
\begin{aretenir}[Official Syllabus Coverage (Upper Sixth -- Second Term)]
This chapter addresses the following official syllabus items:
\begin{itemize}
    \item \textbf{TU 61 -- Lesson 112:} Economic growth -- meaning, measurement (real GDP, real GDP per capita), sources (factor accumulation, productivity) and limits.
    \item \textbf{TU 62 -- Further Studies 8:} Simple growth models -- Harrod-Domar model ($g = s/v$), Rostow's stages of growth, and other perspectives (Solow, endogenous growth).
    \item \textbf{TU 63 -- Lesson 113:} Economic development -- meaning, indicators (HDI, HPI, MDGs/SDGs), strategies, and the distinction between growth and development.
\end{itemize}
\end{aretenir}

\begin{aretenir}[Key Competences Targeted]
\begin{itemize}
    \item Compute and interpret real GDP growth and real GDP per capita growth.
    \item Apply the Harrod-Domar model to diagnose a financing gap and propose policy responses.
    \item Classify an economy within Rostow's stages using structural and social evidence.
    \item Evaluate whether growth translates into development using the Human Development Index (HDI) and Amartya Sen's capability approach.
    \item Synthesise quantitative analysis into a concise policy memo for decision-makers.
\end{itemize}
\end{aretenir}

\courssub{Case Study: Camerland Take-off Diagnostic and Strategy Memo (2025--2029)}

\begin{experience}[Context and Data]
\textbf{Context.} You are a team of junior economic advisers attached to the Ministry of Economy, Planning and Regional Development (MINEPAT) of \textbf{Camerland}, a lower-middle-income CEMAC country whose structure closely mirrors Cameroon's. The Minister has asked each advisory group to prepare a \textbf{``Take-off Diagnostic and Strategy Memo''} for the next five-year development plan (2025--2029). The objective is to assess whether Camerland can achieve the \textbf{7\% annual real GDP growth target} set in the National Development Strategy (SND30) and, crucially, whether that growth will improve the well-being of the population.

\textbf{Stylised Data Sheet for Camerland (2023 estimates).}
\begin{center}
\begin{tabularx}{\linewidth}{|l|X|}\hline
\rowcolor{popblueL}
\textbf{Indicator} & \textbf{Value} \\ \hline
Nominal GDP (2023) & 30\,000 billion CFA francs \\ \hline
Real GDP (2015 base year, 2023) & 25\,000 billion CFA francs \\ \hline
Real GDP (2015 base year, 2022) & 24\,000 billion CFA francs \\ \hline
Population (2023) & 28.0 million \\ \hline
Population growth rate & 2.5\% per year \\ \hline
Gross National Saving Rate ($s$) & 14\% of GDP \\ \hline
Incremental Capital-Output Ratio ($v$) & 3.5 \\ \hline
Sectoral share of real GDP (2023) & Primary: 15\% \quad Secondary: 20\% \quad Services: 65\% \\ \hline
Adult literacy rate (15+ years) & 78\% \\ \hline
Life expectancy at birth & 60 years \\ \hline
Mean years of schooling (adults 25+) & 6.2 years \\ \hline
Expected years of schooling (children) & 10.5 years \\ \hline
GNI per capita (PPP, current international \$) & 3\,800 \$ \\ \hline
\end{tabularx}
\end{center}

\textbf{Additional qualitative notes.}
\begin{itemize}
    \item The primary sector is dominated by cocoa, coffee, cotton, timber and subsistence food crops; productivity growth is slow.
    \item The secondary sector includes agro-processing (SONARA-type refinery, aluminium smelter), light manufacturing and construction.
    \item The service sector is led by wholesale/retail trade, transport, telecommunications (mobile money penetration $>$ 80\%) and public administration.
    \item Tax revenue is 13.5\% of GDP; the budget deficit is 3.2\% of GDP, financed partly by regional bond issuance on the BEAC market.
    \item The financial system is bank-dominated; credit to the private sector is 14\% of GDP.
\end{itemize}
\end{experience}

\courssub{Guided Activities}
\textbf{Work in groups of 4--5.} Use the data sheet and your course notes. Show all calculations and reasoning. Prepare a one-page policy memo (Task 5) to be defended in a 10-minute plenary debate.

\begin{enumerate}
    \item \textbf{Growth measurement.}
          \begin{enumerate}
              \item Compute Camerland's \textbf{real GDP growth rate} for 2023.
              \item Compute the \textbf{real GDP per capita} for 2022 and 2023, then derive the \textbf{per capita real growth rate}.
              \item Comment on the gap between the two rates and its implication for living standards.
          \end{enumerate}
    \item \textbf{Harrod-Domar analysis.}
          \begin{enumerate}
              \item Using $g = s/v$, calculate the \textbf{warranted growth rate} implied by the current saving rate and capital-output ratio.
              \item Determine the \textbf{required saving rate} to achieve the 7\% growth target, assuming $v$ remains constant.
              \item Compute the \textbf{financing gap} (in percentage points of GDP and in absolute CFA francs) between the required and actual saving rates.
              \item Discuss two realistic policies to close the financing gap (domestic and external).
          \end{enumerate}
    \item \textbf{Rostow's stages of growth.}
          \begin{enumerate}
              \item Place Camerland in one of Rostow's five stages. Justify your choice using \textbf{at least three} pieces of evidence from the data sheet (sectoral structure, social indicators, savings/investment behaviour, etc.).
              \item Identify the key \textbf{``leading sectors''} that could drive the take-off, linking them to Camerland's comparative advantages.
          \end{enumerate}
    \item \textbf{Growth vs. Development.}
          \begin{enumerate}
              \item Using the HDI components (health, education, income), assess Camerland's current level of human development. (You may approximate the HDI value using the standard formula or describe the likely category.)
              \item If the 7\% growth target is met \textbf{without} structural changes in income distribution, health or education access, would you consider that \textbf{development} has occurred? Explain using the distinction between growth and development.
              \item Propose one specific policy that would make growth more \textbf{inclusive} and \textbf{sustainable}.
          \end{enumerate}
    \item \textbf{Policy Memo (one A4 page max).} Write a memo addressed to the Minister of Economy with:
          \begin{itemize}
              \item A clear \textbf{subject line} and \textbf{executive summary} (3--4 lines).
              \item \textbf{Two prioritized strategies} (e.g., ``Boost domestic resource mobilisation'' and ``Reduce the capital-output ratio through infrastructure and technology''), each with:
                    \begin{itemize}
                        \item A measurable target (e.g., ``Raise tax revenue to 16\% of GDP by 2027'').
                        \item Two concrete actions.
                        \item A risk/mitigation note.
                    \end{itemize}
              \item A closing sentence linking the strategies to both the 7\% growth target and HDI improvement.
          \end{itemize}
\end{enumerate}

\courssub{Model Answers with Explanatory Notes}

\courssub{Task 1 -- Growth Measurement}

\begin{definition}[Real GDP Growth Rate]
The real GDP growth rate measures the percentage change in the volume of goods and services produced in an economy from one period to the next, using constant prices (a base year) to eliminate inflation effects.
\[
g_Y = \frac{Y_t - Y_{t-1}}{Y_{t-1}} \times 100
\]
\end{definition}

\begin{exemplebox}[Real GDP Growth Calculation for Camerland (2023)]
\begin{align*}
g_Y &= \frac{\text{Real GDP}_{2023} - \text{Real GDP}_{2022}}{\text{Real GDP}_{2022}} \times 100 \\
    &= \frac{25\,000 - 24\,000}{24\,000} \times 100 \\
    &= \frac{1\,000}{24\,000} \times 100 \\
    &= \mathbf{4.17\%}
\end{align*}
\end{exemplebox}

\begin{definition}[Real GDP per Capita]
Real GDP per capita = Real GDP / Population. It approximates the average income available per person and is a better indicator of living standards than aggregate GDP.
\end{definition}

\begin{exemplebox}[Real GDP per Capita and Per Capita Growth]
\textbf{Step 1: Estimate 2022 population.}
Given 2023 population = 28.0 million and growth rate = 2.5\%:
\[
\text{Pop}_{2022} = \frac{28.0}{1.025} \approx 27.317 \text{ million}
\]

\textbf{Step 2: Compute per capita real GDP.}
\begin{align*}
\text{Per capita}_{2022} &= \frac{24\,000 \times 10^9}{27.317 \times 10^6} \approx 878\,600 \text{ CFA francs} \\
\text{Per capita}_{2023} &= \frac{25\,000 \times 10^9}{28.0 \times 10^6} \approx 892\,857 \text{ CFA francs}
\end{align*}

\textbf{Step 3: Per capita growth rate.}
\begin{align*}
g_{\text{pc}} &= \frac{892\,857 - 878\,600}{878\,600} \times 100 \approx \mathbf{1.62\%}
\end{align*}

\textbf{Shortcut (approximation):} $g_{\text{pc}} \approx g_Y - g_{\text{pop}} = 4.17\% - 2.5\% = 1.67\%$ (slight difference due to rounding).
\end{exemplebox}

\begin{attention}[Interpretation Trap]
Do not confuse aggregate growth with per capita growth. A high population growth rate dilutes the gains from aggregate growth. Here, 4.17\% aggregate growth translates into only 1.62\% per capita growth. At this rate, the Rule of 70 gives a doubling time of $\approx 70/1.62 \approx 43$ years -- too slow for rapid poverty reduction.
\end{attention}

\begin{aretenir}[Key Point]
Camerland's 2023 real GDP growth (4.17\%) is below the SND30 target of 7\%. Per capita growth (1.62\%) is modest, implying slow improvement in average living standards unless structural changes accelerate productivity.
\end{aretenir}

\courssub{Task 2 -- Harrod-Domar Analysis}

\begin{propriete}[Harrod-Domar Growth Equation (Examinable Derivation)]
The model assumes a fixed capital-output ratio $v = \Delta K / \Delta Y$ and that investment $I$ equals saving $S = sY$. Since $\Delta K = I$ (ignoring depreciation for simplicity), we have:
\[
\Delta Y = \frac{\Delta K}{v} = \frac{sY}{v} \quad \Rightarrow \quad g = \frac{\Delta Y}{Y} = \frac{s}{v}
\]
This is the \textbf{warranted growth rate} ($g_w$) -- the growth rate that maintains full employment of capital.
\end{propriete}

\begin{exemplebox}[Warranted Growth Rate and Financing Gap]
\textbf{1. Current warranted growth rate:}
\[
g_w = \frac{s}{v} = \frac{0.14}{3.5} = 0.04 = \mathbf{4\%}
\]
This closely matches the observed 4.17\%, suggesting the economy is near its warranted path.

\textbf{2. Required saving rate for 7\% target:}
\[
s_{\text{req}} = g_{\text{target}} \times v = 0.07 \times 3.5 = 0.245 = \mathbf{24.5\% \text{ of GDP}}
\]

\textbf{3. Financing gap:}
\begin{align*}
\text{Gap (percentage points)} &= 24.5\% - 14\% = \mathbf{10.5 \text{ pp of GDP}} \\
\text{Gap (absolute)} &= 0.105 \times 30\,000 \text{ billion} = \mathbf{3\,150 \text{ billion CFA francs}}
\end{align*}
\end{exemplebox}

\begin{methode}[Policies to Close the Financing Gap]
\begin{enumerate}
    \item \textbf{Domestic Resource Mobilisation (DRM):} Broaden the tax base and improve collection efficiency.
          \begin{itemize}
              \item \textit{Action:} Digitalise land registry and property taxation in major cities (Yaoundé, Douala, Bafoussam) by end-2025; introduce simplified turnover tax regime for informal SMEs (threshold 50M CFAF) with mobile-money filing.
              \item \textit{Target:} Raise tax revenue from 13.5\% to 16\% of GDP by 2027 (+2.5 pp $\approx$ 750 bn CFAF/yr).
              \item \textit{Risk/Mitigation:} Political resistance from urban elites $\rightarrow$ earmark 50\% of new property tax revenue for visible neighbourhood infrastructure (roads, street lighting, clinics).
          \end{itemize}
    \item \textbf{External Concessional Financing \& Efficiency Gains:} Lower the effective $v$ while accessing cheap capital.
          \begin{itemize}
              \item \textit{Action:} Negotiate IDA/AfDB concessional loans and issue diaspora bonds for high-return infrastructure (energy, transport); launch ``Energy-Transport Corridor'' PPP (Lom Pangar $\rightarrow$ industrial zones $\rightarrow$ Kribi deep-sea port) to cut logistics costs by 20\% and power outages by 50\%.
              \item \textit{Target:} Reduce incremental capital-output ratio $v$ from 3.5 to 3.0 by 2028 $\rightarrow$ required $s$ for 7\% growth falls from 24.5\% to 21\% (closes 3.5 pp of gap).
              \item \textit{Risk/Mitigation:} PPP contingent liabilities $\rightarrow$ use IMF-approved PPP fiscal risk framework; cap sovereign guarantees at 2\% of GDP/yr.
          \end{itemize}
\end{enumerate}
\end{methode}

\begin{savaistu}[Cameroon Context]
Cameroon's actual tax revenue has hovered around 13--14\% of GDP for years (well below the CEMAC convergence criterion of 17\%). The ``Programme d'Urgence Triennal'' (PLANUT) and SND30 both emphasise DRM. The Harrod-Domar gap of 10.5 pp illustrates why Cameroon relies on BEAC advances and Eurobonds -- but debt service now exceeds 20\% of revenue, crowding out social spending.
\end{savaistu}

\courssub{Task 3 -- Rostow's Stages of Growth}

\begin{definition}[Rostow's Five Stages of Growth (1960)]
\begin{enumerate}
    \item \textbf{Traditional Society:} Subsistence agriculture, low savings ($<$5\%), limited technology, hierarchical social structure.
    \item \textbf{Preconditions for Take-off:} Investment rises (5--10\% of GDP), infrastructure expands, commercial agriculture emerges, social overhead capital builds.
    \item \textbf{Take-off:} Investment $>$10--12\% of GDP, manufacturing ``leading sectors'' drive growth, self-sustaining growth begins, institutional transformation.
    \item \textbf{Drive to Maturity:} Investment 15--20\%+, technology spreads across sectors, economy diversifies, heavy industry replaces light manufacturing.
    \item \textbf{Age of High Mass Consumption:} Services dominate, durable goods consumption rises, welfare state expands.
\end{enumerate}
\end{definition}

\begin{exemplebox}[Classifying Camerland: Take-off Stage (Stage 3)]
\textbf{Justification with three+ evidence points:}
\begin{enumerate}
    \item \textbf{Sectoral structure:} Primary sector only 15\%, Secondary 20\%, Services 65\%. The declining primary share and rising secondary/services share signal structural transformation typical of take-off, where manufacturing and modern services become leading sectors.
    \item \textbf{Savings/Investment behaviour:} Gross saving rate of 14\% exceeds Rostow's 10--12\% take-off threshold but is below the 20\%+ of drive to maturity.
    \item \textbf{Social/technological indicators:} Adult literacy 78\% and mobile money penetration $>$80\% show a workforce capable of absorbing new technologies. However, life expectancy (60) and mean schooling (6.2 years) indicate human capital is still incomplete -- consistent with a society in transition.
    \item \textbf{Institutional context:} Existence of a national development plan (SND30), active regional integration (CEMAC, AfCFTA), and a banking system providing 14\% of GDP credit to private sector suggest institutional preconditions are met.
\end{enumerate}
\end{exemplebox}

\begin{attention}[Common Mistake]
Do not confuse ``Take-off'' with ``Preconditions'' just because indicators are not perfect. Rostow's stages are ideal types; real economies show mixed features. The key is whether manufacturing/industrial growth has become \textbf{self-sustaining} and investment exceeds the critical threshold. Camerland's 14\% savings and expanding secondary sector meet this.
\end{attention}

\begin{exemplebox}[Leading Sectors for Camerland's Take-off]
\begin{itemize}
    \item \textbf{Agro-processing (cocoa grinding, cotton ginning, fruit juicing):} Leverages primary comparative advantage, creates manufacturing jobs, raises export unit values, reduces vulnerability to raw commodity price swings.
    \item \textbf{Digital services \& fintech:} Builds on mobile money ubiquity ($>$80\%), low marginal cost, high scalability; can export services regionally (CEMAC, AfCFTA).
    \item \textbf{Renewable energy (hydro, solar):} Reduces energy costs for industry (binding constraint), attracts FDI, directly lowers incremental capital-output ratio ($v$).
\end{itemize}
\end{exemplebox}

\begin{popfigure}
\begin{tikzpicture}[scale=0.9, every node/.style={font=\small}]
% Rostow stages diagram
\node[draw, rounded corners, fill=popblueL, minimum width=3.5cm, minimum height=1cm] (s1) at (0,0) {1. Traditional Society};
\node[draw, rounded corners, fill=popblueL, minimum width=3.5cm, minimum height=1cm] (s2) at (4,0) {2. Preconditions};
\node[draw, rounded corners, fill=poporange, minimum width=3.5cm, minimum height=1cm] (s3) at (8,0) {\textbf{3. Take-off}};
\node[draw, rounded corners, fill=popblueL, minimum width=3.5cm, minimum height=1cm] (s4) at (12,0) {4. Drive to Maturity};
\node[draw, rounded corners, fill=popblueL, minimum width=3.5cm, minimum height=1cm] (s5) at (16,0) {5. High Mass Consumption};

\draw[->, thick, popdark] (s1.east) -- (s2.west);
\draw[->, thick, popdark] (s2.east) -- (s3.west);
\draw[->, thick, poporange] (s3.east) -- (s4.west);
\draw[->, thick, popdark] (s4.east) -- (s5.west);

% Annotations for Camerland
\node[draw, rounded corners, fill=popgreenL, minimum width=5cm, minimum height=1.2cm, align=center, anchor=north] (cam) at (8,-1.8) {
\textbf{Camerland (2023)} \\
Savings 14\% > 10\% threshold \\
Secondary sector 20\%, Services 65\% \\
Literacy 78\%, Mobile money >80\%
};
\draw[->, thick, popgreen] (s3.south) -- (cam.north);
\end{tikzpicture}
\legende{Rostow's Stages of Growth -- Camerland positioned at Take-off (Stage 3).}
\end{popfigure}

\courssub{Task 4 -- Growth vs. Development}

\begin{definition}[Human Development Index (HDI) -- 2010 Methodology]
The HDI is the geometric mean of three normalised indices:
\[
\text{HDI} = \sqrt[3]{I_{\text{Health}} \times I_{\text{Education}} \times I_{\text{Income}}}
\]
\begin{itemize}
    \item \textbf{Health index:} $I_{\text{Health}} = \frac{\text{LE} - 20}{85 - 20}$ (Life Expectancy at birth).
    \item \textbf{Education index:} $I_{\text{Education}} = \frac{1}{2}\left(\frac{\text{MYS}}{15} + \frac{\text{EYS}}{18}\right)$ (Mean Years of Schooling, Expected Years of Schooling).
    \item \textbf{Income index:} $I_{\text{Income}} = \frac{\ln(\text{GNIpc}) - \ln(100)}{\ln(75\,000) - \ln(100)}$ (GNI per capita PPP \$).
\end{itemize}
\end{definition}

\begin{exemplebox}[HDI Approximation for Camerland (2023)]
\textbf{Health index:} $\frac{60 - 20}{85 - 20} = \frac{40}{65} = \mathbf{0.615}$

\textbf{Education index:}
\begin{align*}
I_{\text{MYS}} &= \frac{6.2}{15} = 0.413 \\
I_{\text{EYS}} &= \frac{10.5}{18} = 0.583 \\
I_{\text{Education}} &= \frac{0.413 + 0.583}{2} = \mathbf{0.498}
\end{align*}

\textbf{Income index:}
\begin{align*}
I_{\text{Income}} &= \frac{\ln(3800) - \ln(100)}{\ln(75000) - \ln(100)} \\
&= \frac{8.242 - 4.605}{11.225 - 4.605} = \frac{3.637}{6.620} = \mathbf{0.550}
\end{align*}

\textbf{HDI:}
\[
\text{HDI} = \sqrt[3]{0.615 \times 0.498 \times 0.550} = \sqrt[3]{0.168} \approx \mathbf{0.552}
\]
\textbf{Classification: Medium Human Development (0.550--0.699).}
\end{exemplebox}

\begin{propriete}[Growth vs. Development -- Core Distinction]
\begin{itemize}
    \item \textbf{Economic Growth:} Quantitative increase in real GDP (or real GDP per capita). Necessary but \textbf{not sufficient} for development.
    \item \textbf{Economic Development:} Qualitative improvement in living standards, capabilities, and freedom (Sen, 1999). Requires structural transformation, equitable distribution, and progress in health, education, and institutions.
\end{itemize}
\textbf{Kuznets' Inverted-U Hypothesis:} Inequality may rise in early growth stages then fall -- but this is not automatic; policy matters.
\end{propriete}

\begin{exemplebox}[Would 7\% Growth Alone Constitute Development?]
\textbf{No.} If the 7\% growth is achieved without structural changes in income distribution, health or education access:
\begin{itemize}
    \item The extra income may accrue mainly to capital owners and urban elites (high inequality persists).
    \item Public revenues may not be invested in health, education, social protection $\rightarrow$ HDI components stagnate.
    \item Growth without development is ``jobless growth'' or ``voiceless growth'' -- it expands the pie but does not share it or build capabilities.
\end{itemize}
\textbf{Development requires:} (1) Structural transformation (shift to higher-productivity activities), (2) Inclusive institutions that translate growth into capabilities (Sen's approach), (3) Sustainability (environmental and fiscal).
\end{exemplebox}

\begin{methode}[Inclusive and Sustainable Policy Proposal]
\textbf{Policy: Universal Health Coverage (UHC) financed by a progressive health levy on mining/telecom profits.}
\begin{itemize}
    \item \textbf{Inclusiveness:} Directly raises health index (reduces out-of-pocket impoverishment), covers informal sector workers.
    \item \textbf{Sustainability:} Improves labour productivity (lower $v$), builds social cohesion, reduces inequality.
    \item \textbf{Financing:} Levy of 2--3\% on super-profits of extractive/telecom sectors (low elasticity), complemented by BEAC-supported health bonds.
    \item \textbf{HDI impact:} Life expectancy could rise from 60 to 65+ within 5 years, pushing HDI toward 0.60 (high medium).
\end{itemize}
\end{methode}

\courssub{Task 5 -- Policy Memo (Model Answer)}

\begin{exemplebox}[Policy Memo to the Minister of Economy]
\begin{center}
\fbox{\begin{minipage}{0.92\linewidth}
\textbf{REPUBLIC OF CAMERLAND}\\
\textbf{MINISTRY OF ECONOMY, PLANNING AND REGIONAL DEVELOPMENT}\\
\textbf{INTERNAL MEMO -- CONFIDENTIAL}
\end{minipage}}
\end{center}

\vspace{0.5cm}
\noindent\textbf{TO:} Minister of Economy, Planning and Regional Development\\
\textbf{FROM:} Advisory Group \#3 (Macroeconomic Modelling Unit)\\
\textbf{DATE:} 15 March 2025\\
\textbf{SUBJECT:} \textbf{Closing the Financing Gap for 7\% Growth \& Ensuring Inclusive Development -- Two Priority Strategies for SND30 (2025--2029)}

\vspace{0.3cm}
\noindent\textbf{EXECUTIVE SUMMARY.} Camerland's current warranted growth is 4\% ($s=14\%, v=3.5$). Achieving the 7\% SND30 target requires raising the saving rate to 24.5\% of GDP -- a financing gap of 10.5 pp (3\,150 bn CFAF). Without structural reforms, 7\% growth will not lift HDI from 0.55 (Medium). We recommend two complementary strategies: (1) Boost domestic resource mobilisation to 16\% of GDP by 2027; (2) Reduce the incremental capital-output ratio to 3.0 via targeted infrastructure and digitalisation. Together they close 60\% of the gap and directly improve health/education outcomes.

\vspace{0.3cm}
\noindent\textbf{STRATEGY 1: BOOST DOMESTIC RESOURCE MOBILISATION (DRM)}
\begin{itemize}
    \item \textbf{Target:} Tax revenue $\uparrow$ from 13.5\% to 16\% of GDP by 2027 (+2.5 pp $\approx$ 750 bn CFAF/yr).
    \item \textbf{Action 1:} Deploy a \textbf{digitalised land registry \& property tax} in 10 major cities (Yaoundé, Douala, Bafoussam, etc.) by end-2025; broadens base, reduces evasion.
    \item \textbf{Action 2:} Introduce a \textbf{simplified turnover tax regime} for informal SMEs (threshold 50M CFAF) with mobile-money filing; brings 200\,000+ firms into net.
    \item \textbf{Risk/Mitigation:} Political resistance from urban elites $\rightarrow$ earmark 50\% of new property tax revenue for visible neighbourhood infrastructure (roads, street lighting, clinics).
\end{itemize}

\vspace{0.2cm}
\noindent\textbf{STRATEGY 2: REDUCE THE INCREMENTAL CAPITAL-OUTPUT RATIO ($v$) FROM 3.5 TO 3.0}
\begin{itemize}
    \item \textbf{Target:} Lower $v$ to 3.0 by 2028 $\rightarrow$ reduces required saving rate for 7\% growth from 24.5\% to 21\% (closes 3.5 pp of gap).
    \item \textbf{Action 1:} Launch a \textbf{``Energy-Transport Corridor'' PPP} (Lom Pangar $\rightarrow$ industrial zones $\rightarrow$ Kribi deep-sea port) to cut logistics costs by 20\% and power outages by 50\%.
    \item \textbf{Action 2:} Scale up \textbf{digital public infrastructure} (interoperable ID, e-procurement, tax API) to reduce transaction costs for firms; pilot in agro-processing clusters.
    \item \textbf{Risk/Mitigation:} PPP contingent liabilities $\rightarrow$ use IMF-approved PPP fiscal risk framework; cap sovereign guarantees at 2\% of GDP/yr.
\end{itemize}

\vspace{0.3cm}
\noindent\textbf{CONCLUSION.} Implementing these two strategies simultaneously attacks the financing gap from both the numerator (savings) and denominator (capital efficiency) of the Harrod-Domar equation, while the earmarked social spending and digital inclusion directly raise the HDI components. This is the path from \textbf{growth} to \textbf{development} -- the essence of Camerland's take-off.

\vspace{0.2cm}
\noindent\textit{Prepared for Ministerial Arbitration -- 20 March 2025.}
\end{exemplebox}

\begin{aretenir}[Examination Tips for This Chapter]
\begin{itemize}
    \item \textbf{Growth measurement:} Always distinguish nominal vs. real, aggregate vs. per capita. Show the formula and units.
    \item \textbf{Harrod-Domar:} Know the derivation $g = s/v$. Be able to compute warranted growth, required savings, and financing gap. Discuss limitations (fixed $v$, no technology, knife-edge instability).
    \item \textbf{Rostow:} Memorise the five stages and key thresholds (savings 10--12\% for take-off). Apply to a country case using structural evidence.
    \item \textbf{HDI:} Know the three components, the min/max values, and the geometric mean formula. Be able to approximate.
    \item \textbf{Growth vs. Development:} Use Sen's capability approach, Kuznets curve, and the distinction between necessary/sufficient conditions. Always link policies to both GDP and HDI.
\end{itemize}
\end{aretenir}

\newpage

\coursec{Methods and worked examples}

\coursec{Economic Growth: Meaning, Measurement and Sources}

\courssub{Skill 1: Computing and Interpreting the Growth Rate of Real GDP}

\begin{definition}[Economic Growth]
Economic growth is the sustained increase in the \cle{real gross domestic product (real GDP)} of an economy over time. It measures the expansion of the volume of goods and services produced, valued at constant prices to eliminate the effect of inflation.
\end{definition}

\begin{propriete}[Growth Rate Formula]
The annual rate of economic growth $g$ (in percent) is computed as:
\[
g = \frac{Y_t - Y_{t-1}}{Y_{t-1}} \times 100
\]
where $Y_t$ is real GDP in year $t$ and $Y_{t-1}$ is real GDP in the previous year, both measured in the \emph{same base year} prices.
\end{propriete}

\begin{methode}[Computing the rate of economic growth]
\begin{enumerate}
\item \textbf{Identify the correct variable.} Growth must be measured with \cle{real GDP} (GDP at constant prices), never nominal GDP, because nominal GDP mixes price changes and output changes.
\item \textbf{Apply the growth-rate formula:}
\[ g = \frac{Y_t - Y_{t-1}}{Y_{t-1}} \times 100 \]
where $Y_t$ is real GDP in year $t$ and $Y_{t-1}$ real GDP in the previous year.
\item \textbf{If nominal GDP and a price index are given}, first deflate: $\text{Real GDP} = \dfrac{\text{Nominal GDP}}{\text{Price index}} \times 100$. The price index is usually the GDP deflator (base year = 100).
\item \textbf{Interpret the result.} A positive $g$ means expansion of output; compare it with population growth to judge whether living standards are really rising.
\item \textbf{Check units and rounding.} Express the answer as a percentage, usually to one decimal place, and state the base year of the constant prices used.
\end{enumerate}
\end{methode}

\begin{exemplebox}[Chain-linking and base years]
National accounts offices (e.g. INS in Cameroon) periodically update the base year (e.g. from 2000 to 2016) and \cle{chain-link} the series to avoid substitution bias. In examinations, you will be given data already expressed in a single base year; do not mix series with different base years.
\end{exemplebox}

\begin{exoresolu}[Growth rate of the Cameroonian economy (fictional data)]
The real GDP of a CEMAC country, measured at constant 2016 prices, was 22\,000 billion FCFA in 2022 and 23\,100 billion FCFA in 2023. Over the same period, the population grew from 27.2 million to 28.0 million inhabitants.
\begin{enumerate}
\item Calculate the rate of economic growth in 2023.
\item Calculate real GDP per head in 2022 and 2023 and comment.
\end{enumerate}
\tcblower
\textbf{1. Rate of economic growth.}
\begin{align*}
g &= \frac{Y_{2023} - Y_{2022}}{Y_{2022}} \times 100 \\
&= \frac{23\,100 - 22\,000}{22\,000} \times 100 \\
&= \frac{1\,100}{22\,000} \times 100 = 5.0\ \%
\end{align*}
The economy grew by \textbf{5.0\,\%} in 2023: total output of goods and services expanded.

\textbf{2. Real GDP per head.}
\[ \text{GDP per head}_{2022} = \frac{22\,000 \times 10^9}{27.2 \times 10^6} \approx 808\,824 \ \text{FCFA} \]
\[ \text{GDP per head}_{2023} = \frac{23\,100 \times 10^9}{28.0 \times 10^6} = 825\,000 \ \text{FCFA} \]
The growth of GDP per head is:
\[ \frac{825\,000 - 808\,824}{808\,824} \times 100 \approx 2.0\ \% \]

\textbf{Comment.} Population grew by $\frac{28.0-27.2}{27.2}\times 100 \approx 2.9\ \%$. Because output (5.0\,\%) grew faster than population (2.9\,\%), income per head rose by about 2.0\,\%. This illustrates the key rule: \emph{living standards improve only when output growth exceeds population growth}. Had population grown at 5\,\% too, GDP per head would have stagnated despite respectable aggregate growth.
\end{exoresolu}

\courssub{Skill 2: Distinguishing Nominal GDP, Real GDP and Per Capita Measures}

\begin{definition}[Nominal vs Real GDP]
\begin{itemize}
\item \cle{Nominal GDP} (GDP at current prices) values output at the prices of the year of production.
\item \cle{Real GDP} (GDP at constant prices) values output at the prices of a chosen base year, using the GDP deflator.
\item \cle{GDP deflator} = $\dfrac{\text{Nominal GDP}}{\text{Real GDP}} \times 100$. It is a Paasche index covering all domestically produced goods and services.
\end{itemize}
\end{definition}

\begin{methode}[Deflating nominal GDP]
\begin{enumerate}
\item \textbf{Recall the relationship:} $\text{Nominal GDP} = \text{Real GDP} \times \dfrac{\text{GDP deflator}}{100}$.
\item \textbf{Rearrange to deflate:} $\text{Real GDP} = \text{Nominal GDP} \times \dfrac{100}{\text{GDP deflator}}$.
\item \textbf{Compare real GDP across years}, not nominal GDP, to isolate the change in the volume of production.
\item \textbf{Divide by population} to obtain per capita figures when the question concerns welfare or living standards.
\item \textbf{State the limitation} in interpretation questions: per capita GDP is an average and hides inequality, the informal sector and non-market production.
\end{enumerate}
\end{methode}

\begin{exoresolu}[Nominal versus real growth]
In a Central African economy, nominal GDP rose from 10\,000 billion FCFA in 2021 to 11\,220 billion FCFA in 2022. The GDP price index (deflator) rose from 100 to 110 over the same period.
\begin{enumerate}
\item Calculate the nominal growth rate.
\item Calculate real GDP for each year (base 2021) and the real growth rate.
\item Explain the difference between the two rates.
\end{enumerate}
\tcblower
\textbf{1. Nominal growth rate.}
\[ g_{\text{nominal}} = \frac{11\,220 - 10\,000}{10\,000} \times 100 = 12.2\ \% \]

\textbf{2. Real GDP and real growth.}
\[ \text{Real GDP}_{2021} = 10\,000 \times \frac{100}{100} = 10\,000 \ \text{billion FCFA} \]
\[ \text{Real GDP}_{2022} = 11\,220 \times \frac{100}{110} = 10\,200 \ \text{billion FCFA} \]
\[ g_{\text{real}} = \frac{10\,200 - 10\,000}{10\,000} \times 100 = 2.0\ \% \]

\textbf{3. Explanation.} The price level rose by $\frac{110-100}{100}\times 100 = 10\ \%$. Of the 12.2\,\% nominal increase, about 10 percentage points simply reflect \cle{inflation}; only 2\,\% corresponds to a genuine increase in the volume of goods and services produced. This is why examiners insist that economic growth be measured in \emph{real} terms: using nominal figures would grossly overstate the performance of the economy. A government reporting 12.2\,\% "growth" from nominal data would be misleading the public.
\end{exoresolu}

\begin{aretenir}[Sources of Economic Growth]
Growth accounting (Solow decomposition) identifies three proximate sources:
\begin{enumerate}
\item \textbf{Growth of the labour force} (quantity and quality / human capital).
\item \textbf{Growth of the capital stock} (physical capital accumulation).
\item \textbf{Total Factor Productivity (TFP) growth} / technical progress — the residual not explained by factor accumulation.
\end{enumerate}
In the long run, TFP growth is the only sustainable source of rising living standards.
\end{aretenir}

\coursec{The Harrod--Domar Growth Model}

\courssub{Skill 3: Applying the Harrod--Domar Growth Model}

\begin{definition}[Harrod--Domar Model (Simplified)]
The model links the \cle{warranted growth rate} $g_w$ to the savings ratio $s$ and the capital--output ratio $c$ (or ICOR):
\[
g_w = \frac{s}{c}
\]
where $s = S/Y$ (savings as a share of national income) and $c = \Delta K / \Delta Y$ (incremental capital--output ratio, ICOR), i.e. the extra capital needed to produce one extra unit of output.
\end{definition}

\begin{propriete}[Key Assumptions]
\begin{enumerate}
\item Fixed coefficients production function (Leontief): capital and labour are used in fixed proportions; no substitution.
\item Constant returns to scale.
\item Savings are a constant fraction of income ($S = sY$) and are automatically invested ($I = S$).
\item No technical progress, no government, no foreign trade (closed economy).
\item The capital--output ratio $c$ is constant.
\item Labour supply grows at a constant natural rate $g_n$ (Harrod) or is unlimited (Domar).
\end{enumerate}
\end{propriete}

\begin{methode}[Using the Harrod--Domar formula $g = s/c$]
\begin{enumerate}
\item \textbf{Write down the model's identity:} the warranted growth rate is
\[ g = \frac{s}{c} \]
where $s$ is the national savings ratio ($S/Y$) and $c$ is the incremental capital--output ratio (ICOR).
\item \textbf{Identify the two given variables} and solve for the third: $s = g \times c$ or $c = s/g$.
\item \textbf{Remember the logic:} growth requires investment; investment must be financed by savings; the ICOR tells how much capital is needed to produce one extra unit of output.
\item \textbf{Check plausibility:} a high $c$ (e.g.\ 4--5) means capital is used inefficiently or the economy lacks infrastructure; a low $s$ (typical of low-income countries) constrains growth — this is the \emph{savings gap}.
\item \textbf{Conclude with evaluation} if asked: the model assumes a fixed $c$, a closed economy, no technical progress, and that savings automatically become investment — all questionable in developing economies.
\end{enumerate}
\end{methode}

\begin{popfigure}
\begin{tikzpicture}[scale=0.9, every node/.style={font=\small}]
% Harrod-Domar flow
\node[draw, rounded corners, fill=popblueL] (sav) at (0,0) {Savings $S = sY$};
\node[draw, rounded corners, fill=popgreenL] (inv) at (4,0) {Investment $I = S$};
\node[draw, rounded corners, fill=poporangeL] (cap) at (8,0) {Capital Stock $\Delta K = I$};
\node[draw, rounded corners, fill=poppinkL] (out) at (12,0) {Output $\Delta Y = \Delta K / c$};
\draw[->, thick] (sav) -- (inv) node[midway, above] {$I=S$};
\draw[->, thick] (inv) -- (cap) node[midway, above] {$\Delta K = I$};
\draw[->, thick] (cap) -- (out) node[midway, above] {$\Delta Y = \Delta K/c$};
% Growth rate formula
\node[draw, rounded corners, fill=popgold, align=center] at (6,-2) {Warranted Growth Rate:\\ $g_w = \frac{\Delta Y}{Y} = \frac{s}{c}$};
\draw[->, dashed] (inv.south) -- ++(0,-0.5) -| (6,-1.5);
\draw[->, dashed] (cap.south) -- ++(0,-0.5) -| (6,-1.5);
\end{tikzpicture}
\legende{The Harrod--Domar causal chain: savings finance investment, which adds to capital, which generates output growth.}
\end{popfigure}

\begin{exoresolu}[Financing a growth target]
A developing economy has a capital--output ratio (ICOR) of 4 and a savings ratio of 12\,\%.
\begin{enumerate}
\item Calculate the growth rate predicted by the Harrod--Domar model.
\item The government targets a growth rate of 6\,\%. What savings ratio is required?
\item If domestic savings cannot rise above 15\,\%, suggest two ways the country could still finance the required investment.
\end{enumerate}
\tcblower
\textbf{1. Predicted growth rate.}
\[ g = \frac{s}{c} = \frac{0.12}{4} = 0.03 = 3\ \% \]
With current behaviour, the economy can grow at only \textbf{3\,\%} per year.

\textbf{2. Required savings ratio for 6\,\% growth.}
\[ s = g \times c = 0.06 \times 4 = 0.24 = 24\ \% \]
The country must save and invest \textbf{24\,\% of national income} — double its current ratio. The shortfall, $24\% - 12\% = 12$ percentage points of GDP, is the \cle{savings gap} (or financing gap).

\textbf{3. Bridging the gap.}
\begin{itemize}
\item \textbf{Foreign aid and concessional loans} (e.g.\ from the World Bank or the African Development Bank) can supplement domestic savings, though they create dependency and debt-servicing burdens.
\item \textbf{Foreign direct investment (FDI)}, for instance in Cameroon's oil, timber or agro-industrial sectors, brings capital without requiring domestic savings, but profits may be repatriated.
\end{itemize}
One could also mention improving the \emph{efficiency} of capital: if better management lowered $c$ from 4 to 3, the required savings ratio would fall to $0.06 \times 3 = 18\ \%$, much closer to the feasible 15\,\%. This shows that growth policy is not only about saving more but also about investing \emph{better} (lowering the ICOR).
\end{exoresolu}

\courssub{Skill 4: Critical Evaluation of the Harrod--Domar Model}

\begin{methode}[Evaluating a growth model in essay form]
\begin{enumerate}
\item \textbf{State the contribution first:} what does the model explain well? (Here: the central role of saving, investment and capital accumulation — relevant to capital-scarce African economies.)
\item \textbf{List the assumptions} and test each against reality: fixed capital--output ratio; savings automatically invested; closed economy; no technical progress; labour unlimited.
\item \textbf{Draw the implications for developing countries:} the model justifies aid ("financing gap" approach) but ignores human capital, institutions, corruption, and the productivity of investment.
\item \textbf{Conclude with a balanced judgement:} useful as a planning benchmark (many five-year plans used it), insufficient as a complete theory of growth.
\end{enumerate}
\end{methode}

\begin{exoresolu}[Essay: "The Harrod--Domar model adequately explains growth in developing countries." Discuss.]
Write a structured essay plan with developed arguments.
\tcblower
\textbf{Introduction.} Define the model ($g = s/c$) and its purpose: it shows that growth depends on the share of income saved and invested and on the productivity of capital. Thesis: it illuminates one condition of growth — capital accumulation — but cannot by itself explain the growth experience of developing countries.

\textbf{Paragraph 1 — Strengths.}
\begin{itemize}
\item Correctly stresses that low-income countries are trapped by low savings: low income $\Rightarrow$ low saving $\Rightarrow$ low investment $\Rightarrow$ low growth (the \emph{vicious circle of poverty}).
\item Provides a practical planning tool: target growth rates can be translated into required investment and aid needs; post-war reconstruction plans and many African development plans used exactly this logic (the "financing gap" approach of the World Bank and IMF).
\item Highlights the \cle{knife-edge instability}: if actual growth $g$ deviates from warranted growth $g_w$, the economy moves into cumulative inflation or deflation (Harrod) — a warning against unbalanced expansion.
\end{itemize}

\textbf{Paragraph 2 — Weaknesses.}
\begin{itemize}
\item The capital--output ratio is \emph{not fixed}: it varies with technology, management and the type of investment (a road versus a steel mill). The ICOR often falls as an economy develops.
\item Savings do not automatically become productive investment: financial systems are underdeveloped, funds may be misallocated, embezzled or held abroad (capital flight). The \cle{investment gap} may be larger than the savings gap.
\item It ignores \cle{human capital}, education and health, which modern growth theory (and the Asian experience) shows to be decisive.
\item It ignores technical progress, institutions, political stability and trade — factors that explain why countries with similar savings ratios grow very differently.
\item The \cle{natural growth rate} $g_n$ (labour force growth) may differ from $g_w$, leading to chronic unemployment or labour shortages (Harrod's instability).
\end{itemize}

\textbf{Paragraph 3 — Evidence and alternative views.} Some economies grew fast with modest investment thanks to efficiency gains (TFP); others invested heavily (oil boom projects in Nigeria or Gabon) with disappointing results because of poor governance. Rostow and later endogenous growth theories add stages, technology and knowledge to the picture.

\textbf{Conclusion.} The Harrod--Domar model is a valuable \emph{starting point} — capital accumulation is necessary — but it is not sufficient. Growth in developing countries also requires human capital, sound institutions and productive use of investment.
\end{exoresolu}

\coursec{Rostow's Stages of Growth and Other Perspectives}

\courssub{Skill 5: Identifying Rostow's Stages from a Description}

\begin{definition}[Rostow's Five Stages of Economic Growth (1960)]
\begin{enumerate}
\item \textbf{Traditional Society:} Subsistence agriculture, low technology, hierarchical social structure, limited surplus, production for own consumption.
\item \textbf{Preconditions for Take-off:} Investment in social overhead capital (transport, power, ports), emergence of a commercial/entrepreneurial class, rise in savings/investment to $\approx 5-10\%$ of national income, one or two "leading sectors" appear.
\item \textbf{Take-off:} Self-sustaining growth; investment rises to $>10\%$ of national income; leading sectors drive rapid expansion; new industries appear; institutional changes support modern economic activity.
\item \textbf{Drive to Maturity:} Technology spreads to all sectors; economy diversifies; investment $10-20\%$; reliance on imports falls; professional/technical workforce grows.
\item \textbf{Age of High Mass Consumption:} Shift to durable consumer goods and services; real income per head rises sharply; welfare state emerges; sectoral structure dominated by services.
\end{enumerate}
\end{definition}

\begin{popfigure}
\begin{tikzpicture}[scale=0.85, every node/.style={font=\small}]
% Rostow stages flowchart
\node[draw, rounded corners, fill=popblueL, align=center] (s1) at (0,0) {Stage 1:\\Traditional Society};
\node[draw, rounded corners, fill=popgreenL, align=center] (s2) at (3.5,0) {Stage 2:\\Preconditions for\\Take-off};
\node[draw, rounded corners, fill=poporangeL, align=center] (s3) at (7,0) {Stage 3:\\Take-off};
\node[draw, rounded corners, fill=poppinkL, align=center] (s4) at (10.5,0) {Stage 4:\\Drive to Maturity};
\node[draw, rounded corners, fill=popgold, align=center] (s5) at (14,0) {Stage 5:\\High Mass\\Consumption};
\draw[->, thick] (s1) -- (s2);
\draw[->, thick] (s2) -- (s3);
\draw[->, thick] (s3) -- (s4);
\draw[->, thick] (s4) -- (s5);
% Key features below
\node[align=center, text width=2.5cm] at (0,-1.2) {Subsistence agri.\\(70-80\% labour)};
\node[align=center, text width=2.5cm] at (3.5,-1.2) {Infrastructure,\\rising savings,\\leading sector};
\node[align=center, text width=2.5cm] at (7,-1.2) {Investment $>10\%$,\\self-sustaining\\growth};
\node[align=center, text width=2.5cm] at (10.5,-1.2) {Diversification,\\tech diffusion,\\import substitution};
\node[align=center, text width=2.5cm] at (14,-1.2) {Services dominant,\\durables,\\welfare state};
\end{tikzpicture}
\legende{Rostow's linear stages of growth.}
\end{popfigure}

\begin{methode}[Applying Rostow's five stages]
\begin{enumerate}
\item \textbf{Memorise the five stages in order:} (1) traditional society; (2) preconditions for take-off; (3) take-off; (4) drive to maturity; (5) age of high mass consumption.
\item \textbf{Look for diagnostic clues} in the text: subsistence agriculture and custom $\Rightarrow$ stage 1; infrastructure investment, rising savings, a leading sector emerging $\Rightarrow$ stage 2; self-sustaining industrial growth, savings ratio jumping to around 10\,\% or more $\Rightarrow$ stage 3; diversification into many industries $\Rightarrow$ stage 4; dominance of consumer durables and services $\Rightarrow$ stage 5.
\item \textbf{Justify the classification} by quoting the clue and matching it to the stage's definition.
\item \textbf{Be ready to criticise:} the stages are based on Western history, assume linearity, and fit poorly economies dependent on a single export commodity (e.g.\ oil or cocoa).
\end{enumerate}
\end{methode}

\begin{exoresolu}[Classifying economies according to Rostow]
Classify each of the following hypothetical economies into one of Rostow's stages, justifying your answer:
\begin{enumerate}
\item Economy A: 70\,\% of the labour force works in subsistence agriculture; output per head has been stagnant for generations; land is held communally.
\item Economy B: the state is building railways, ports and power dams; a commercial class is emerging; savings are rising but industry remains embryonic.
\item Economy C: textile and food-processing industries are expanding rapidly; investment has reached 12\,\% of national income; growth has become self-sustaining.
\end{enumerate}
\tcblower
\textbf{1. Economy A — Stage 1: the traditional society.} The dominance of subsistence agriculture, stagnant productivity and production organised by custom rather than markets are the defining features of a traditional society. There is little surplus to invest, so growth is impossible.

\textbf{2. Economy B — Stage 2: the preconditions for take-off.} Investment in \emph{social overhead capital} (railways, ports, power) and the appearance of a commercial, enterprising class show that the society is preparing the conditions for growth. Savings are rising but have not yet reached the critical threshold; industry is still embryonic. This resembles the phase of colonial and early post-independence infrastructure building in many African countries.

\textbf{3. Economy C — Stage 3: take-off.} The investment ratio crossing roughly 10\,\% of national income, the emergence of \cle{leading sectors} (textiles, food processing) and growth becoming self-sustaining are precisely Rostow's criteria for take-off — the decisive break with stagnation.

\textbf{Critical note.} Rostow's schema assumes every country follows the same linear path. A country whose "take-off" rests on a single commodity (cocoa, oil) may see growth collapse when world prices fall, showing that the stages are not automatic or irreversible. Many African economies remain stuck in Stage 2 or experience "aborted take-off" due to commodity dependence, weak institutions, and lack of structural transformation.
\end{exoresolu}

\begin{savaistu}[Rostow and African Economies]
Rostow wrote with the US and Western Europe in mind. In Cameroon, the "preconditions" phase saw massive infrastructure investment (e.g. the Kribi deep-sea port, Lom Pangar dam, railway extensions). However, the "take-off" has been elusive because leading sectors (oil, cocoa, timber, aluminium) are enclave industries with weak linkages to the rest of the economy. This illustrates the \cle{structuralist critique}: growth requires not just investment but \emph{structural transformation} and \emph{linkages} (Hirschman).
\end{savaistu}

\coursec{Economic Development: Meaning, Indicators and Strategies}

\courssub{Skill 6: Distinguishing Growth from Development and Interpreting the HDI}

\begin{definition}[Economic Growth vs Economic Development]
\begin{itemize}
\item \cle{Economic Growth}: A quantitative increase in real GDP (or real GDP per capita) over time. It is a \emph{necessary} condition for development.
\item \cle{Economic Development}: Growth \emph{plus} qualitative improvements in living standards — health, education, reduced poverty and inequality, structural transformation of the economy, and greater freedoms. It is a \emph{multidimensional} process.
\end{itemize}
\end{definition}

\begin{propriete}[Human Development Index (HDI) — UNDP]
Since 2010, the HDI is the geometric mean of three normalized dimension indices:
\[
\text{HDI} = \sqrt[3]{I_{\text{Health}} \times I_{\text{Education}} \times I_{\text{Income}}}
\]
where
\begin{itemize}
\item $I_{\text{Health}}$: Life expectancy at birth (min 20, max 85).
\item $I_{\text{Education}}$: Mean of (Mean years of schooling, max 15) and (Expected years of schooling, max 18).
\item $I_{\text{Income}}$: GNI per capita (PPP \$), log scale (min 100, max 75\,000).
\end{itemize}
Value ranges: 0--1. Categories: Very High ($\ge 0.800$), High ($0.700-0.799$), Medium ($0.550-0.699$), Low ($<0.550$).
\end{propriete}

\begin{popfigure}
\begin{tikzpicture}[scale=1.1, every node/.style={font=\small}]
% HDI triangle
\coordinate (A) at (0,2.5);
\coordinate (B) at (-2.165,-1.25);
\coordinate (C) at (2.165,-1.25);
\draw[thick, popblue] (A) -- (B) -- (C) -- cycle;
\node[popblue] at (0,2.9) {HDI};
\node[align=center] at (0,1.2) {Health\\Life Expectancy};
\node[align=center] at (-2.5,-1.8) {Education\\Mean \& Expected\\Years of Schooling};
\node[align=center] at (2.5,-1.8) {Income\\GNI per capita (PPP)};
% Dimensions labels on edges
\node[rotate=60, popdark] at (-1.2,0.5) {Long \& Healthy Life};
\node[rotate=-60, popdark] at (1.2,0.5) {Decent Standard of Living};
\node[popdark] at (0,-1.5) {Knowledge};
\end{tikzpicture}
\legende{The three dimensions of the Human Development Index (geometric mean since 2010).}
\end{popfigure}

\begin{methode}[Analysing development indicators]
\begin{enumerate}
\item \textbf{Define the distinction:} \cle{economic growth} is a quantitative rise in real output; \cle{economic development} is growth \emph{plus} qualitative change — better health, education, reduced poverty and inequality, structural transformation.
\item \textbf{Recall the HDI's three dimensions:} a long and healthy life (life expectancy at birth), knowledge (mean and expected years of schooling), and a decent standard of living (GNI per head, PPP).
\item \textbf{Interpret an HDI value} between 0 and 1: closer to 1 means higher human development; compare countries or years, never treat it as a complete measure.
\item \textbf{Cross-check with GDP per head:} if GDP per head is high but HDI is low, growth has not been translated into development (typical of some oil economies).
\item \textbf{Mention limitations:} the HDI ignores inequality within the country (use \cle{IHDI} — Inequality-adjusted HDI), environmental sustainability, and political freedoms. Complement with \cle{Multidimensional Poverty Index (MPI)}, Gini coefficient, and poverty headcount ratios.
\end{enumerate}
\end{methode}

\begin{exoresolu}[Growth without development?]
Countries X and Y have the following data:
\begin{center}
\begin{tabular}{|l|c|c|c|}
\hline
\rowcolor{popblueL} Country & GDP per head (USD) & Life expectancy (years) & Mean schooling (years) \\
\hline
X & 4\,000 & 55 & 5 \\
\hline
Y & 2\,500 & 68 & 9 \\
\hline
\end{tabular}
\end{center}
\begin{enumerate}
\item Which country is more "developed"? Justify.
\item What does this comparison reveal about the limits of GDP per head as a development indicator?
\end{enumerate}
\tcblower
\textbf{1. Which country is more developed?} By GDP per head alone, X appears richer (4\,000 USD against 2\,500 USD). But development is multidimensional. Y's citizens live 13 years longer and receive on average 4 more years of schooling, indicating better health services, nutrition and education. Y would therefore record a higher \cle{HDI} despite lower income: it has converted modest resources into human well-being more effectively. Country Y is arguably \emph{more developed}, while X is merely \emph{richer on average}. This pattern is observed in reality: some oil-exporting economies (e.g. Equatorial Guinea, Gabon historically) display relatively high income per head alongside mediocre health and education outcomes, because resource rents are concentrated and weakly redistributed.

\textbf{2. Limits of GDP per head.}
\begin{itemize}
\item It is an \textbf{average}: it conceals inequality. If oil rents accrue to a small elite, most citizens may be poor despite high per capita GDP.
\item It ignores the \textbf{non-monetary dimensions} of welfare: health, education, security, freedom, a clean environment.
\item It omits the \textbf{informal sector} and subsistence production, which are large in African economies (estimated 30-50\% of GDP in many CEMAC countries), so it mismeasures even the income dimension.
\item It says nothing about the \textbf{sustainability} of growth (depletion of forests, oil reserves, soil degradation).
\end{itemize}
\textbf{Conclusion.} GDP per head measures \emph{growth}; the HDI and complementary indicators (poverty rates, inequality, literacy, MPI) are needed to assess \emph{development}. The examiner's favourite formulation: \emph{growth is a necessary but not sufficient condition for development}.
\end{exoresolu}

\courssub{Skill 7: Evaluating Development Strategies in Essay Form}

\begin{definition}[Major Development Strategies]
\begin{itemize}
\item \cle{Import-Substitution Industrialisation (ISI)}: Producing domestically goods previously imported, protected by tariffs/quotas. Inward-looking.
\item \cle{Export Promotion (EP) / Export-Oriented Industrialisation}: Orienting production towards world markets. Outward-looking.
\item \cle{Agricultural Development / Rural Development}: Raising productivity in agriculture (Green Revolution, land reform, extension services) to release labour and surplus for industry.
\item \cle{Human Capital Investment}: Education, health, nutrition, skills training.
\item \cle{Regional Integration}: Enlarging markets (e.g. \cle{AfCFTA} — African Continental Free Trade Area, CEMAC, ECOWAS).
\item \cle{Structural Transformation}: Deliberate policy to shift labour from low- to high-productivity sectors (industrial policy, special economic zones).
\end{itemize}
\end{definition}

\begin{methode}[Structuring an essay on development strategies]
\begin{enumerate}
\item \textbf{Define the strategy} precisely (import-substitution industrialisation, export promotion, agricultural development, human capital investment, aid, regional integration such as AfCFTA).
\item \textbf{Explain the mechanism:} how is the strategy supposed to raise growth and development?
\item \textbf{Give advantages} with a concrete illustration (Cameroonian, CEMAC or African example).
\item \textbf{Give limitations and risks}, again illustrated.
\item \textbf{Conclude with a nuanced judgement:} strategies are complements rather than substitutes; success depends on governance, institutions and sequencing.
\end{enumerate}
\end{methode}

\begin{exoresolu}[Essay: Import substitution versus export promotion]
"For a country like Cameroon, export promotion is a better development strategy than import-substitution industrialisation." Discuss.
\tcblower
\textbf{Introduction.} Define the two strategies. \cle{Import-substitution industrialisation (ISI)} consists of producing domestically goods previously imported, protected by tariffs and quotas. \cle{Export promotion} consists of orienting production towards world markets. Thesis: each has merits, but their success depends on how they are implemented and on the country's specific context.

\textbf{Paragraph 1 — The case for ISI.} It reduces dependence on imports, saves foreign exchange, creates urban jobs and allows infant industries to learn behind protection. After independence, many African states, including Cameroon, established processing industries (sugar at SOSUCAM, textiles at CICAM, aluminium at Edea/ALUCAM) on this logic. The "infant industry argument" (List, Mill) justifies temporary protection until firms achieve economies of scale.

\textbf{Paragraph 2 — The limits of ISI.} Protected firms often become inefficient and rent-seeking; domestic markets are small (Cameroon ~28 million), so economies of scale are impossible; protection raises prices for consumers and downstream industries; and industries still import machinery and inputs, so the foreign-exchange saving is smaller than expected. Several state enterprises accumulated losses and had to be restructured or privatised in the 1990s SAPs.

\textbf{Paragraph 3 — The case for export promotion.} Exporting forces firms to be competitive, earns foreign exchange, exploits comparative advantage (cocoa, coffee, cotton, timber, oil for Cameroon) and permits economies of scale through access to world demand. The success of East Asian economies (South Korea, Taiwan) illustrates its potential when combined with disciplined industrial policy.

\textbf{Paragraph 4 — The limits of export promotion.} Dependence on a few primary commodities exposes the country to volatile world prices (a fall in cocoa or oil prices immediately hits export earnings and the budget); primary exports have low value added and declining terms of trade (Prebisch-Singer hypothesis); and world markets are protected by rich countries' subsidies (EU Common Agricultural Policy, US cotton subsidies). Export promotion of \emph{unprocessed} commodities can lock a country into underdevelopment — hence the importance of local processing (grinding cocoa locally rather than exporting raw beans) and of regional markets through the \cle{AfCFTA} to build larger markets for manufactured goods.

\textbf{Paragraph 5 — Synthesis for Cameroon.} Cameroon's "Vision 2035" aims at structural transformation: moving from a commodity exporter to an upper-middle-income industrialising economy. This requires \emph{both} selective protection of genuine infant industries (e.g. pharmaceuticals, agro-processing) \emph{and} export orientation with quality upgrading. The \cle{AfCFTA} offers a continental market of 1.3 billion consumers to overcome the small domestic market constraint. However, neither strategy works without \cle{good governance}, infrastructure (energy, transport), skilled labour, and a stable macroeconomic environment (CFA franc stability helps).

\textbf{Conclusion.} Neither strategy is superior in the abstract. A pragmatic path for Cameroon combines selective protection of genuine infant industries, diversification and local processing of exports, investment in human capital and infrastructure, and regional integration to widen markets — with good governance as the precondition for any strategy to work.
\end{exoresolu}

\begin{attention}[Frequent errors in this chapter]
\begin{itemize}
\item Computing growth with \textbf{nominal} GDP instead of real GDP.
\item Confusing the \textbf{level} of GDP with its \textbf{growth rate}: a rich country can grow slowly, a poor one quickly.
\item In Harrod--Domar exercises, forgetting to convert percentages into decimals ($s = 0.12$, not 12) before applying $g = s/c$.
\item Confusing the \cle{average capital-output ratio} ($K/Y$) with the \cle{incremental capital-output ratio (ICOR)} ($\Delta K/\Delta Y$); the model uses ICOR.
\item Treating GDP per head as a measure of \emph{development}: always add health, education and distribution.
\item Listing Rostow's stages without \textbf{justifying} the classification from the clues in the text.
\item One-sided essays: the GCE rewards balanced evaluation — advantages \emph{and} limits, with examples.
\item Forgetting that the HDI uses a \textbf{geometric mean} (since 2010), not an arithmetic mean — this penalises unbalanced development.
\end{itemize}
\end{attention}

\begin{exercice}[Practice Questions]
\begin{enumerate}
\item The real GDP of a country (base 2015) was 15\,000 billion FCFA in 2021 and 15\,750 billion FCFA in 2022. Population was 25.0 million and 25.8 million respectively. Calculate (a) the economic growth rate, (b) real GDP per capita for both years, (c) the growth rate of real GDP per capita. Comment.
\item An economy has an ICOR of 3.5 and a target growth rate of 7\,\%. If the current savings ratio is 18\,\%, calculate the savings gap. Suggest two policies to close it, other than raising domestic savings.
\item "Rostow's stages of growth are irrelevant to modern African economies." Discuss with reference to at least two African countries.
\item Explain why a country with a high GDP per capita may have a low HDI. Use the concepts of inequality, resource curse, and social spending priorities in your answer.
\item Evaluate the view that regional integration (AfCFTA) is the most promising development strategy for CEMAC countries.
\end{enumerate}
\end{exercice}

\begin{corrige}[Exercise 1]
\begin{enumerate}
\item[(a)] $g = \frac{15\,750 - 15\,000}{15\,000} \times 100 = 5.0\%$.
\item[(b)] GDP per capita 2021 = $\frac{15\,000 \times 10^9}{25.0 \times 10^6} = 600\,000$ FCFA. GDP per capita 2022 = $\frac{15\,750 \times 10^9}{25.8 \times 10^6} \approx 610\,465$ FCFA.
\item[(c)] Growth of GDP per capita = $\frac{610\,465 - 600\,000}{600\,000} \times 100 \approx 1.7\%$. Population growth $\approx 3.2\%$. Output growth (5.0\%) > population growth (3.2\%), so living standards rose, but modestly.
\end{enumerate}
\end{corrige}

\begin{corrige}[Exercise 2]
Required savings $s = g \times c = 0.07 \times 3.5 = 0.245 = 24.5\%$. Savings gap = $24.5\% - 18\% = 6.5$ percentage points of GDP.
Policies: (1) Attract FDI in priority sectors (e.g. special economic zones with tax incentives). (2) Negotiate concessional loans/grants from multilateral banks (AfDB, World Bank) for infrastructure that lowers the ICOR over time. (3) Improve public investment management to reduce waste and lower the effective ICOR.
\end{corrige}

\newpage

\coursec{Exercises}

\courssub{I check my knowledge}

\begin{exercice}[Multiple choice questions]
For each question, choose the single correct answer.
\begin{enumerate}
\item Economic growth is best defined as:
\begin{enumerate}
\item an improvement in the general well-being of the population;
\item a sustained increase in real national output over time;
\item a reduction in income inequality;
\item an increase in the money supply.
\end{enumerate}
\item If nominal GDP rises from $800$ billion FCFA to $920$ billion FCFA while prices rise by $10\%$, the real growth rate is approximately:
\begin{enumerate}
\item $15\%$; \item $10\%$; \item $4.5\%$; \item $25\%$.
\end{enumerate}
\item In the Harrod--Domar model, the warranted growth rate is given by:
\begin{enumerate}
\item $g_w = v/s$; \item $g_w = s \times v$; \item $g_w = s/v$; \item $g_w = s + v$.
\end{enumerate}
\item According to Rostow, the stage during which an economy becomes self-sustaining, with growth spreading to many sectors, is:
\begin{enumerate}
\item the traditional society; \item the preconditions for take-off; \item the take-off; \item the age of high mass consumption.
\end{enumerate}
\end{enumerate}
\end{exercice}

\begin{exercice}[True or false — justify your answer]
State whether each assertion is \textbf{true} or \textbf{false}, and justify briefly.
\begin{enumerate}
\item A country can experience economic growth without economic development.
\item GDP per capita is a perfect measure of the standard of living of a population.
\item In the Harrod--Domar model, a higher saving rate leads, \emph{ceteris paribus}, to a higher growth rate.
\item The Human Development Index (HDI) combines income, education and life expectancy.
\end{enumerate}
\end{exercice}

\begin{exercice}[Definitions]
Define precisely each of the following terms in two or three sentences:
\begin{enumerate}
\item economic growth;
\item economic development;
\item the capital--output ratio;
\item the take-off (in Rostow's theory).
\end{enumerate}
\end{exercice}

\begin{exercice}[Direct calculations]
\begin{enumerate}
\item A country's real GDP rises from $5\,000$ billion FCFA to $5\,300$ billion FCFA in one year. Calculate the rate of economic growth.
\item The population of this country grows from $25$ million to $25.75$ million over the same period. Calculate the growth rate of real GDP \emph{per capita} and comment on the difference with the answer in (a).
\item Given a saving rate $s = 15\%$ and a capital--output ratio $v = 3$, compute the warranted growth rate of the economy.
\end{enumerate}
\end{exercice}

\courssub{I practise}

\begin{exercice}[Real and per capita growth from a data table]
The table below gives data for a hypothetical Central African economy.
\begin{center}
\begin{tabular}{|l|c|c|}
\hline
\rowcolor{popblueL} & \textbf{Year 1} & \textbf{Year 2} \\ \hline
Nominal GDP (billion FCFA) & $12\,000$ & $13\,440$ \\ \hline
Price index (base 100, Year 1) & $100$ & $105$ \\ \hline
Population (million) & $20.0$ & $20.6$ \\ \hline
\end{tabular}
\end{center}
\begin{enumerate}
\item Calculate real GDP for each year.
\item Calculate the rate of economic growth between Year 1 and Year 2.
\item Calculate real GDP per capita for each year and its growth rate.
\item Comment on your results: is the population better off? Justify.
\end{enumerate}
\end{exercice}

\begin{exercice}[Applying the Harrod--Domar model]
In a developing economy, the saving rate is $s = 12\%$ and the capital--output ratio is $v = 4$.
\begin{enumerate}
\item Compute the warranted growth rate of this economy.
\item The government wishes to achieve a growth target of $6\%$ per year. Assuming $v$ remains constant, calculate the saving rate required.
\item Identify and explain \textbf{two} policies this government could use to raise the saving rate, and \textbf{one} difficulty a low-income country faces in doing so.
\item Explain what happens to the required saving rate if the capital--output ratio falls to $v = 3$ because of better technology. What lesson do you draw for development policy?
\end{enumerate}
\end{exercice}

\begin{exercice}[Ranking GDP per capita against HDI]
The table below shows data for five countries (hypothetical values).
\begin{center}
\begin{tabular}{|l|c|c|}
\hline
\rowcolor{popblueL} \textbf{Country} & \textbf{GDP per capita (USD)} & \textbf{HDI rank} \\ \hline
A & $9\,500$ & 2 \\ \hline
B & $8\,200$ & 1 \\ \hline
C & $6\,000$ & 4 \\ \hline
D & $5\,800$ & 3 \\ \hline
E & $3\,100$ & 5 \\ \hline
\end{tabular}
\end{center}
\begin{enumerate}
\item Rank the five countries by GDP per capita.
\item Compare this ranking with the HDI ranking and identify the countries whose positions change.
\item Explain \textbf{two} reasons why a country with higher GDP per capita may have a lower HDI rank.
\item What does this exercise reveal about the limitations of GDP per capita as a measure of development?
\end{enumerate}
\end{exercice}

\begin{exercice}[Distinguishing key concepts]
Write a short paragraph (five to eight lines) clearly distinguishing each pair of concepts, giving one example each time:
\begin{enumerate}
\item economic growth and economic development;
\item the warranted growth rate and the natural growth rate;
\item extensive growth and intensive growth;
\item GDP and GNP.
\end{enumerate}
\end{exercice}

\courssub{I prepare for the GCE}

\begin{exercice}[Structured question: measuring growth (20 marks)]
The table below relates to the economy of a CEMAC member state.
\begin{center}
\begin{tabular}{|l|c|c|}
\hline
\rowcolor{popblueL} & \textbf{2022} & \textbf{2023} \\ \hline
Nominal GDP (billion FCFA) & $25\,000$ & $27\,500$ \\ \hline
GDP deflator (base 100 in 2022) & $100$ & $104$ \\ \hline
Population (million) & $27.0$ & $27.8$ \\ \hline
\end{tabular}
\end{center}
\begin{enumerate}
\item Define \emph{economic growth} and \emph{economic development}. \textbf{(4 marks)}
\item Calculate real GDP for 2022 and 2023, and the rate of economic growth in 2023. \textbf{(5 marks)}
\item Calculate real GDP per capita for each year and its growth rate. \textbf{(4 marks)}
\item Explain why the growth rate of GDP per capita differs from the growth rate of GDP. \textbf{(3 marks)}
\item Give \textbf{two} reasons why the increase in GDP per capita calculated above may overstate the improvement in the population's welfare. \textbf{(4 marks)}
\end{enumerate}
\end{exercice}

\begin{exercice}[Essay: growth and development (20 marks)]
\textit{``A high rate of economic growth is a necessary but not sufficient condition for economic development.''}

Discuss this statement with reference to Cameroon. Your essay should:
\begin{enumerate}
\item clearly distinguish growth from development;
\item explain why growth is generally necessary for development;
\item show, using Cameroonian examples (oil, cocoa, regional inequalities, education and health), why growth alone does not guarantee development;
\item conclude with a balanced personal judgement.
\end{enumerate}
\end{exercice}

\begin{exercice}[Essay: growth models and African economies (20 marks)]
Answer \textbf{one} of the following two questions.

\textbf{Question A.} \textit{``Rostow's stages of growth are irrelevant to African economies.''} Assess this view critically. In your answer, present the five stages, apply them to an African economy of your choice, and evaluate at least three limitations of the model.

\textbf{Question B.} Examine the main obstacles to economic development in Cameroon (you may consider the saving gap, the capital--output ratio, dependence on primary exports, human capital and institutions), and suggest appropriate strategies to overcome them. Use the Harrod--Domar framework where relevant.
\end{exercice}

\newpage

\coursec{Solutions to the exercises}

\begin{corrige}[Exercise 1 — Multiple choice questions]
\begin{enumerate}
\item \textbf{Answer: (b)} — a sustained increase in real national output over time. Economic growth is a purely quantitative, long-run phenomenon measured by the rise of real GDP or real national income. Option (a) describes economic \emph{development}, (c) describes a distributional change, and (d) is a monetary phenomenon which, if unmatched by output, simply causes inflation.
\item \textbf{Answer: (c)} — approximately $4.5\%$. 
\begin{align*}
\text{Nominal growth} &= \frac{920-800}{800}\times 100 = 15\%. \\
\text{Real GDP in Year 2} &= \frac{920}{1.10} = 836.36 \text{ billion FCFA}. \\
\text{Real growth} &= \frac{836.36-800}{800}\times 100 \approx 4.55\% \approx 4.5\%.
\end{align*}
The approximation $g_{\text{real}} \approx g_{\text{nominal}} - \pi$ gives $15\% - 10\% = 5\%$, close but less precise.
\item \textbf{Answer: (c)} — $g_w = s/v$, where $s$ is the saving rate and $v$ the capital--output ratio. Growth is positively related to saving and negatively related to the amount of capital needed per unit of output.
\item \textbf{Answer: (c)} — the take-off. During take-off, growth becomes self-sustaining: investment rises to about $10\%$ or more of national income, leading sectors emerge, and expansion spreads through the economy. The age of high mass consumption comes later, when durable consumer goods dominate.
\end{enumerate}
\end{corrige}

\begin{corrige}[Exercise 2 — True or false]
\begin{enumerate}
\item \textbf{True.} Growth is a quantitative rise in real output; development includes qualitative changes (health, education, reduced poverty and inequality, structural transformation). A country exporting oil may record high GDP growth while most of the population sees no improvement in living standards — growth without development.
\item \textbf{False.} GDP per capita is only an average. It ignores the distribution of income, non-marketed (subsistence) production, the informal sector, working conditions, environmental damage, health and education. Two countries with identical GDP per capita may have very different standards of living.
\item \textbf{True.} Since $g_w = s/v$, with $v$ constant, any rise in $s$ raises $g_w$ proportionally. For example, with $v=4$, raising $s$ from $10\%$ to $16\%$ raises growth from $2.5\%$ to $4\%$.
\item \textbf{True.} The HDI, published by the UNDP, is a composite index of three dimensions: a long and healthy life (life expectancy at birth), knowledge (mean and expected years of schooling), and a decent standard of living (GNI per capita).
\end{enumerate}
\end{corrige}

\begin{aretenir}[Key distinction: Growth vs Development]
Economic \textbf{growth} is a \emph{quantitative} concept: sustained increase in real output (real GDP). Economic \textbf{development} is a \emph{qualitative and quantitative} process: rising real income per head \emph{plus} improvements in health, education, nutrition, reduced poverty/inequality, and structural transformation. Growth is \emph{necessary but not sufficient} for development.
\end{aretenir}

\begin{corrige}[Exercise 3 — Definitions]
\begin{enumerate}
\item \textbf{Economic growth} is the sustained increase, over a long period, in the real output of goods and services of an economy, usually measured by the percentage change in real GDP (or real GDP per capita). It is a quantitative phenomenon.
\item \textbf{Economic development} is a broader, qualitative and quantitative process involving rising real incomes per head together with improvements in health, education, nutrition, reduced poverty and inequality, and structural transformation of the economy (from subsistence agriculture towards industry and services). It implies an improvement in the general well-being of the population.
\item \textbf{The capital--output ratio} ($v$) is the amount of capital required to produce one additional unit of output: $v = \dfrac{\Delta K}{\Delta Y} = \dfrac{K}{Y}$. A ratio of $4$ means that $4$ FCFA of investment are needed to generate $1$ FCFA of extra annual output; the lower $v$, the more productive capital is.
\item \textbf{The take-off} is the third of Rostow's five stages: the decisive period (often $20$--$30$ years) during which the economy breaks the resistance of the traditional society. Investment rises to $10\%$ or more of national income, one or more leading manufacturing sectors expand rapidly, and growth becomes self-sustaining and automatic.
\end{enumerate}
\end{corrige}

\begin{corrige}[Exercise 4 — Direct calculations]
\begin{enumerate}
\item Rate of growth of real GDP:
\[ g = \frac{5\,300 - 5\,000}{5\,000}\times 100 = \frac{300}{5\,000}\times 100 = 6\%. \]
\item Population growth: $\dfrac{25.75-25}{25}\times 100 = 3\%$. Growth of GDP per capita:
\[ g_{pc} \approx g - g_{pop} = 6\% - 3\% = 3\%. \]
(Exact check: GDP per capita rises from $\frac{5\,000}{25}=200\,000$ to $\frac{5\,300}{25.75}\approx 205\,825$ FCFA, i.e. $+2.9\%$.) \textbf{Comment:} because population grows at $3\%$, half of the output growth is absorbed by the extra mouths to feed; average income per person rises by only about $3\%$. This illustrates why rapid population growth can dilute the welfare gains of growth.
\item Warranted growth rate (Harrod--Domar):
\[ g_w = \frac{s}{v} = \frac{15\%}{3} = 5\% \text{ per year.} \]
\end{enumerate}
\end{corrige}

\begin{methode}[Calculating real GDP and real growth]
\begin{enumerate}
\item Obtain nominal GDP and a price index (base year $=100$).
\item Real GDP $=$ Nominal GDP $\times \dfrac{100}{\text{Price Index}}$.
\item Real growth rate $= \dfrac{\text{Real GDP}_t - \text{Real GDP}_{t-1}}{\text{Real GDP}_{t-1}} \times 100$.
\item Real GDP per capita $= \dfrac{\text{Real GDP}}{\text{Population}}$.
\item Growth of real GDP per capita $\approx$ Real GDP growth rate $-$ Population growth rate.
\end{enumerate}
\end{methode}

\begin{corrige}[Exercise 5 — Real and per capita growth from a data table]
\begin{enumerate}
\item Real GDP $=$ nominal GDP $\times \dfrac{100}{\text{price index}}$:
\begin{itemize}
\item Year 1: $12\,000 \times \frac{100}{100} = 12\,000$ billion FCFA.
\item Year 2: $13\,440 \times \frac{100}{105} = 12\,800$ billion FCFA.
\end{itemize}
\item Rate of economic growth:
\[ g = \frac{12\,800 - 12\,000}{12\,000}\times 100 = \frac{800}{12\,000}\times 100 \approx 6.67\%. \]
(Note: nominal growth was $12\%$, but $5\%$ inflation eroded much of it.)
\item Real GDP per capita:
\begin{itemize}
\item Year 1: $\dfrac{12\,000\text{ billion}}{20.0\text{ million}} = 600\,000$ FCFA.
\item Year 2: $\dfrac{12\,800\text{ billion}}{20.6\text{ million}} \approx 621\,359$ FCFA.
\end{itemize}
Growth rate of GDP per capita: $\dfrac{621\,359 - 600\,000}{600\,000}\times 100 \approx 3.56\%$.
\item \textbf{Comment:} Yes, the average citizen is better off, but less than the headline figure suggests. Real output grew by $6.67\%$, yet because the population grew by $3\%$ (from $20.0$ to $20.6$ million), income per head rose by only about $3.6\%$. Moreover, this is an average: it says nothing about how the extra income is distributed, nor about health, education or the informal sector. The lesson is that real GDP per capita, not nominal GDP, is the meaningful measure of growth.
\end{enumerate}
\end{corrige}

\begin{attention}[Common mistake: confusing nominal and real growth]
Always deflate nominal GDP by the price index \emph{before} computing growth rates. A $12\%$ nominal increase with $5\%$ inflation does \emph{not} mean $7\%$ real growth; the correct formula is $\frac{1+g_{\text{nom}}}{1+\pi}-1$, which here gives $\frac{1.12}{1.05}-1 \approx 6.67\%$. The approximation $g_{\text{real}} \approx g_{\text{nom}} - \pi$ is only valid for small rates.
\end{attention}

\begin{corrige}[Exercise 6 — Applying the Harrod--Domar model]
\begin{enumerate}
\item Warranted growth rate:
\[ g_w = \dfrac{s}{v} = \dfrac{12\%}{4} = 3\% \text{ per year.} \]
\item Required saving rate for a target $g_w = 6\%$:
\[ s = g_w \times v = 6\% \times 4 = 24\%. \]
The country must double its saving rate from $12\%$ to $24\%$ — this is the \emph{saving gap} of $12$ percentage points.
\item \textbf{Two policies to raise saving:}
\begin{itemize}
\item \emph{Mobilising domestic financial savings:} developing banks, microfinance institutions and mobile-money savings products (e.g., MTN MoMo, Orange Money in Cameroon) so that rural households and informal-sector workers can save safely and earn interest; government can also run budget surpluses (public saving) through better tax collection.
\item \emph{Attracting external resources to fill the gap:} foreign direct investment (FDI), official development aid (ODA), and remittances from the diaspora (which are very significant for Cameroon, exceeding $300$ billion FCFA annually) supplement domestic saving.
\end{itemize}
\textbf{One difficulty:} in a low-income country, most income is consumed by basic needs (food, housing), so the marginal propensity to save is very low — the \emph{vicious circle of poverty}: low income $\rightarrow$ low saving $\rightarrow$ low investment $\rightarrow$ low growth $\rightarrow$ low income. Forced saving through taxation is also hard because the tax base is narrow and the informal sector large (about $90\%$ of employment in Cameroon).
\item If $v$ falls to $3$ (improved capital productivity): $s = g_w \times v = 6\% \times 3 = 18\%$. The required saving rate falls from $24\%$ to $18\%$. \textbf{Lesson:} growth depends not only on \emph{how much} is invested but on \emph{how efficiently} capital is used. Better technology, management, infrastructure (e.g., reliable electricity from EDC/SONATREL), and maintenance reduce $v$, so development policy should combine resource mobilisation with measures that raise the productivity of investment (education, technology transfer, good governance) rather than focusing on investment volume alone.
\end{enumerate}
\end{corrige}

\begin{savaistu}[Harrod--Domar in Cameroon]
Cameroon's average incremental capital-output ratio (ICOR) has historically been high ($4$--$5$) due to power shortages, poor transport infrastructure, and low project implementation capacity. The \emph{National Development Strategy 2020--2030 (SND30)} explicitly targets lowering the ICOR by improving public investment management and energy supply. Meanwhile, gross national savings remain around $15\%$--$18\%$ of GDP, implying a warranted growth rate of only $3\%$--$4.5\%$, barely above population growth ($\approx 2.6\%$). This explains the persistent gap between actual and desired per capita growth.
\end{savaistu}

\begin{corrige}[Exercise 7 — Ranking GDP per capita against HDI]
\begin{enumerate}
\item Ranking by GDP per capita (highest to lowest): \textbf{A} ($9\,500$) $>$ \textbf{B} ($8\,200$) $>$ \textbf{C} ($6\,000$) $>$ \textbf{D} ($5\,800$) $>$ \textbf{E} ($3\,100$).
\item HDI ranking: B (1st), A (2nd), D (3rd), C (4th), E (5th). \textbf{Changes:} A and B swap positions (B is poorer but better developed); C and D swap positions (D is poorer in income but ranks higher on HDI). Only E keeps the same position.
\item \textbf{Two reasons why higher GDP per capita may mean a lower HDI rank:}
\begin{itemize}
\item \emph{Unequal income distribution:} country A's wealth may be concentrated in a small elite (e.g., oil revenues in an enclave sector), so the average citizen is poor, uneducated and unhealthy, whereas B's income is more evenly spread, funding public services for all.
\item \emph{Different public priorities:} B may invest heavily in free education and primary health care (raising literacy and life expectancy), while A spends on prestige projects or the military; since HDI includes schooling and life expectancy, B scores higher despite lower income.
\end{itemize}
\item \textbf{Conclusion:} GDP per capita measures only the \emph{average quantity} of income. It ignores distribution, health, education and the quality of life. Development is multidimensional, so composite indicators such as the HDI give a fuller — though still imperfect — picture. A country can be ``richer'' yet less ``developed''.
\end{enumerate}
\end{corrige}

\begin{corrige}[Exercise 8 — Distinguishing key concepts]
\begin{enumerate}
\item \textbf{Growth vs development.} Economic growth is the sustained rise in real output (a quantitative change), while economic development adds qualitative improvements: better health, education, less poverty and inequality, and structural change. \emph{Example:} an oil-producing country like Chad or Equatorial Guinea may record $8\%$ GDP growth from crude exports while rural literacy and infant mortality barely improve — growth without development.
\item \textbf{Warranted vs natural growth rate.} In Harrod--Domar, the warranted rate $g_w = s/v$ is the rate at which producers' investment plans are exactly validated by demand, keeping capacity fully used. The natural rate $g_n$ is the maximum rate permitted by the growth of the labour force and productivity (labour-augmenting technical progress). \emph{Example:} if $s/v = 4\%$ but the labour force grows at $3\%$, $g_w > g_n$ and the economy faces a ceiling of labour shortage and inflationary pressure; if $g_w < g_n$, there is chronic unemployment.
\item \textbf{Extensive vs intensive growth.} Extensive growth comes from using \emph{more} factors (more workers, more land, more machines); intensive growth comes from using factors \emph{more efficiently} (higher productivity through technology and skills). \emph{Example:} doubling cocoa output by clearing twice as much forest in the Centre Region is extensive; doubling yields per hectare with improved varieties (e.g., IRAD's high-yielding cocoa seedlings) and fertiliser is intensive.
\item \textbf{GDP vs GNP.} GDP measures output produced \emph{within} the country's territory, whoever produces it; GNP measures income accruing to the country's \emph{nationals}, wherever they are: GNP $=$ GDP $+$ net property income from abroad. \emph{Example:} profits repatriated by a foreign oil company (e.g., Perenco, SNH partners) operating in Cameroon count in Cameroon's GDP but not in its GNP, so Cameroon's GNP is lower than its GDP. Conversely, remittances from Cameroonians abroad count in GNP but not GDP.
\end{enumerate}
\end{corrige}

\begin{corrige}[Exercise 9 — Structured question: measuring growth (20 marks)]
\textbf{(a) Definitions (4 marks — 2 each).} \emph{Economic growth} is the sustained increase over time in the real output of goods and services of an economy, measured by the percentage change in real GDP. \emph{Economic development} is a broader process combining rising real income per head with improvements in health, education, nutrition, reduced poverty and inequality, and structural transformation of the economy. \emph{Examiner's expectation:} the word ``real'' and the quantitative/qualitative distinction must appear for full marks.

\textbf{(b) Real GDP and growth rate (5 marks).}
\begin{itemize}
\item Real GDP 2022: $25\,000 \times \frac{100}{100} = 25\,000$ billion FCFA. \textbf{(1 mark)}
\item Real GDP 2023: $27\,500 \times \frac{100}{104} \approx 26\,442.3$ billion FCFA. \textbf{(2 marks: formula + result)}
\item Growth rate: $g = \dfrac{26\,442.3 - 25\,000}{25\,000}\times 100 \approx 5.77\%$. \textbf{(2 marks)}
\end{itemize}
Note: nominal growth was $10\%$, but $4\%$ inflation reduced real growth to about $5.8\%$.

\textbf{(c) Real GDP per capita (4 marks).}
\begin{itemize}
\item 2022: $\dfrac{25\,000\text{ billion}}{27.0\text{ million}} \approx 925\,926$ FCFA. \textbf{(1 mark)}
\item 2023: $\dfrac{26\,442.3\text{ billion}}{27.8\text{ million}} \approx 951\,162$ FCFA. \textbf{(1 mark)}
\item Growth rate: $\dfrac{951\,162 - 925\,926}{925\,926}\times 100 \approx 2.73\%$. \textbf{(2 marks)}
\end{itemize}

\textbf{(d) Why the rates differ (3 marks).} GDP per capita $=$ real GDP $\div$ population. Its growth rate equals approximately the growth rate of real GDP \emph{minus} the growth rate of population. Here population grew by $\frac{27.8-27.0}{27.0}\times 100 \approx 2.96\%$, so $g_{pc} \approx 5.77\% - 2.96\% \approx 2.81\%$. Part of the extra output is absorbed by the additional population; average income therefore rises more slowly than total output. \emph{Examiner's expectation:} the relationship $g_{pc} \approx g - g_{pop}$ and a numerical illustration.

\textbf{(e) Two reasons welfare may be overstated (4 marks — 2 each, any two):}
\begin{itemize}
\item \emph{Income distribution:} the average may rise while inequality widens — if the gains accrue to a small urban elite or to oil enclaves, the majority's welfare is unchanged.
\item \emph{Non-market and informal production:} subsistence food production, very important in CEMAC rural areas, is imperfectly captured, so GDP mismeasures true living standards.
\item \emph{Composition of output:} growth may come from enclave extractive industries (oil, mining) with few local jobs, or from activities with harmful side-effects (pollution, deforestation) whose costs are not deducted.
\item \emph{Non-income dimensions:} GDP per capita says nothing about health, education, security or leisure, which are essential to welfare.
\end{itemize}
\end{corrige}

\begin{corrige}[Exercise 10 — Essay: growth and development (20 marks)]
\textbf{Indicative mark scheme:} introduction and definitions (4); why growth is necessary (5); why it is not sufficient, with Cameroonian examples (7); balanced conclusion (2); quality of expression and structure (2).

\textbf{Introduction.} Economic growth is the sustained rise in real output; economic development is the wider improvement in human well-being — health, education, equity, structural transformation. The statement claims growth is \emph{necessary} (development cannot occur without it) but \emph{not sufficient} (it does not automatically produce development). This essay argues the statement is largely correct for Cameroon.

\textbf{Why growth is necessary.} Development requires resources that only a growing economy can generate. Growth raises tax revenues, allowing the state to finance schools, hospitals, roads and electricity; it creates employment and raises household incomes, reducing poverty; it permits investment in human capital without reducing current consumption. A stagnant economy faces a zero-sum choice: any redistribution to the poor means taking from someone else, which is socially explosive. Historically, no country has achieved high HDI without first achieving sustained growth.

\textbf{Why growth is not sufficient — Cameroonian evidence.}
\begin{itemize}
\item \emph{Oil enclave:} Cameroon's petroleum exports (via SONARA and SNH) have at times boosted GDP strongly, yet oil is capital-intensive, employs few Cameroonians, and its revenues accrue largely to the state and foreign companies; periods of high oil-driven growth (e.g., 1980s, 2000s) coexisted with persistent rural poverty.
\item \emph{Cocoa and cash crops:} export earnings from cocoa and coffee raise GDP, but smallholder farmers in the Centre, South-West and Littoral regions receive a small share of the world price; growth in the export sector does not automatically transform village living standards.
\item \emph{Regional inequalities:} growth is concentrated in Douala and Yaoundé, while the Far North, North, Adamawa and East regions lag in roads, water, electricity and schools — aggregate growth masks deep spatial disparities.
\item \emph{Education and health:} GDP per capita can rise while school completion, maternal health and access to clean water stagnate if public spending is misallocated or institutions are weak; growth does not guarantee that revenue becomes public services.
\end{itemize}
Growth must therefore be \emph{accompanied} by deliberate policies: progressive taxation and social spending, rural infrastructure, diversification, good governance and investment in human capital — in short, inclusive or pro-poor growth.

\textbf{Conclusion.} The statement is valid: without growth, development is impossible because the resources for schools, clinics and redistribution do not exist; but growth alone, especially enclave or commodity-led growth as in Cameroon, can bypass the majority. Development requires growth \emph{plus} equity-oriented institutions and policies that convert output into well-being.
\end{corrige}

\begin{corrige}[Exercise 11 — Essay: growth models and African economies (20 marks)]
\textbf{Question A — Rostow's stages: indicative mark scheme:} presentation of the five stages (6); application to an African economy (5); at least three limitations (6); conclusion (3).

\textbf{The five stages:} (1) the \emph{traditional society} — subsistence agriculture, low productivity, limited technology, hierarchical social structure; (2) the \emph{preconditions for take-off} — rise of transport, banking, a commercial class, investment in infrastructure, emergence of a manufacturing sector; (3) the \emph{take-off} — investment reaches about $10\%$ of national income, leading sectors emerge, growth becomes self-sustaining; (4) the \emph{drive to maturity} — diversification, technology spreads to the whole economy, import substitution deepens; (5) the \emph{age of high mass consumption} — durable consumer goods and services dominate, real wages rise, welfare state emerges.

\textbf{Application (example: Côte d'Ivoire or Cameroon):} one may argue that a country like Côte d'Ivoire displayed ``preconditions'' through investment in ports (Abidjan), roads and cocoa commercialisation, and a partial ``take-off'' during its post-independence boom (1960s--1970s), with agro-industry as a leading sector. Cameroon similarly built infrastructure (dams at Song Loulou, Edéa; roads) and saw manufacturing emerge around Douala (textiles, aluminium, food processing). This suggests the model retains some descriptive value for the early post-independence period.

\textbf{Limitations (at least three):}
\begin{enumerate}
\item \emph{Historical determinism:} the model generalises the experience of 19th-century Europe; African economies did not evolve in isolation but within colonial structures designed for extraction, so their ``traditional'' stage was distorted, not original.
\item \emph{Stages are not sequential or automatic:} African states can display features of several stages simultaneously (modern banking in cities beside subsistence farming), and take-off is not guaranteed — many have remained stuck despite decades of investment (e.g., the ``lost decades'' of the 1980s--1990s).
\item \emph{It ignores the international environment:} dependence on primary exports, deteriorating terms of trade (Prebisch-Singer hypothesis), debt crises and foreign domination — central to African underdevelopment — are absent from Rostow's closed-economy logic.
\item \emph{It neglects institutions, distribution and human capital:} growth is treated as a matter of investment rates alone, ignoring governance, education, inequality and the role of the state.
\end{enumerate}

\textbf{Conclusion:} Rostow's model is a useful pedagogical description of how investment and leading sectors can launch self-sustaining growth, and some African trajectories loosely echo it; but as a universal law it is inadequate for Africa, where colonial history, external dependence and institutional constraints dominate. It should be complemented, not applied mechanically.

\medskip
\textbf{Question B — Obstacles to development in Cameroon: indicative mark scheme:} identification and explanation of at least four obstacles (10); use of the Harrod--Domar framework (4); strategies (4); structure and expression (2).

\textbf{Obstacles:}
\begin{enumerate}
\item \emph{The saving gap:} with low per capita incomes ($\approx \$1\,500$), domestic saving is low. In Harrod--Domar terms, $g\_w = s/v$: if $s = 12\%$ and $v = 4$, growth is only $3\%$ — barely above population growth ($\approx 2.6\%$), so per capita income stagnates. Reaching $6\%$ growth requires $s = 24\%$, a gap that domestic resources cannot fill.
\item \emph{A high capital--output ratio:} poor infrastructure (only $15\%$ of roads paved), weak maintenance culture, power shortages (frequent load-shedding) and low skills mean each unit of investment yields little output ($v$ is high), so even available saving produces disappointing growth.
\item \emph{Dependence on primary exports:} oil, cocoa, coffee, cotton, timber and aluminium expose Cameroon to volatile world prices and deteriorating terms of trade; export earnings — and thus the capacity to import capital goods — are unstable.
\item \emph{Weak human capital:} shortages of skilled labour, brain drain (many engineers and doctors emigrate), and uneven education and health lower labour productivity and the productivity of capital.
\item \emph{Institutional weaknesses:} corruption (Cameroon ranks low on Transparency International's CPI), administrative inefficiency and insecurity (Boko Haram in Far North, Anglophone crisis in NW/SW) discourage private investment and waste public resources.
\end{enumerate}

\textbf{Strategies:}
\begin{itemize}
\item \emph{Raise saving and finance:} deepen the banking and microfinance system (MC2, UBA, Afriland First Bank), mobilise diaspora remittances via formal channels, broaden the tax base (reduce informality), attract FDI (e.g., in energy, agro-industry) and use concessional aid to fill the saving gap.
\item \emph{Lower the capital--output ratio:} invest in energy (Lom Pangar, Nachtigal, Memve'ele dams), roads (Yaoundé-Douala, Bamenda-Mamfe) and maintenance; improve project selection and management (PIM reforms) and technology transfer so that each FCFA invested yields more output — with $v = 3$, the same $12\%$ saving rate yields $4\%$ growth instead of $3\%$.
\item \emph{Diversify and industrialise:} process cocoa, timber and cotton locally (import-substitution and export-processing zones like Kribi) to escape raw-price dependence and create jobs; implement the \emph{Import-Substitution Policy (PSI)} and \emph{Local Transformation Policy}.
\item \emph{Invest in human capital:} education (especially vocational/technical), health (universal health coverage rollout) and nutrition raise productivity directly and reduce $v$.
\item \emph{Strengthen institutions:} transparent public finance (CONSUPE, Audit Bench of Supreme Court), fight against corruption (CONAC), decentralisation, and regional integration (CEMAC, AfCFTA) to widen markets and attract investment.
\end{itemize}

\textbf{Conclusion:} Cameroon's obstacles are mutually reinforcing — the vicious circle of poverty — but the Harrod--Domar framework shows the two levers of escape: mobilising more resources ($s$) and using them more efficiently ($v$). A credible strategy must act on both, within a supportive institutional framework that ensures inclusive growth and structural transformation.
\end{corrige}

\newpage

\coursec{Summary sheet}

\begin{popfigure}
\begin{tikzpicture}[mindmap, grow cyclic, every node/.style={concept, text=white, font=\small\bfseries, align=center},
root concept/.append style={concept color=popink, font=\large\bfseries},
level 1/.append style={level distance=3.6cm, sibling angle=51.4},
level 2/.append style={level distance=2.2cm, sibling angle=40, font=\scriptsize\bfseries}]
\node[root concept]{Growth \& Development}
  child[concept color=popblue]{ node{Economic Growth}
    child[concept color=popblue!75]{ node{GDP, GNP, per capita} }
    child[concept color=popblue!75]{ node{Sources: labour, capital, tech} }
  }
  child[concept color=poporange]{ node{Measuring Growth}
    child[concept color=poporange!75]{ node{Rate $=\frac{\Delta Y}{Y}\times100$} }
    child[concept color=poporange!75]{ node{Nominal vs real} }
  }
  child[concept color=popgreen]{ node{Harrod--Domar}
    child[concept color=popgreen!75]{ node{$g=s/c$} }
    child[concept color=popgreen!75]{ node{Savings \& ICOR} }
  }
  child[concept color=poppurple]{ node{Rostow's Stages}
    child[concept color=poppurple!75]{ node{5 stages} }
    child[concept color=poppurple!75]{ node{Take-off key} }
  }
  child[concept color=popgold]{ node{Development}
    child[concept color=popgold!75]{ node{HDI, poverty} }
    child[concept color=popgold!75]{ node{Growth $\neq$ development} }
  }
  child[concept color=poppink]{ node{Strategies}
    child[concept color=poppink!75]{ node{Industrialisation} }
    child[concept color=poppink!75]{ node{Human capital} }
  };
\end{tikzpicture}
\legende{Mind map of the chapter: Economic Growth and Development}
\end{popfigure}

\begin{bilan}
\textbf{Key definitions}\par
\begin{itemize}
\item \cle{Economic growth}: a sustained increase in a country's real output (real GDP or real GNP) over time. It is a \emph{quantitative} concept.
\item \cle{Economic development}: a wider \emph{qualitative} improvement in living standards — better health, education, income distribution, institutions and freedoms — alongside growth.
\item \cle{GDP}: value of final goods and services produced \emph{within} a country in a year; \cle{GNP} adds net factor income from abroad.
\item \cle{Per capita income}: $Y/N$ (output divided by population); a rough measure of average living standards.
\item \cle{HDI} (Human Development Index): composite index combining life expectancy, education and per capita income.
\item \cle{ICOR} (Incremental Capital-Output Ratio): extra capital needed to produce one extra unit of output, $c = \Delta K / \Delta Y$.
\end{itemize}

\textbf{Formulas to know}\par
\begin{itemize}
\item Growth rate: $g = \dfrac{Y_t - Y_{t-1}}{Y_{t-1}} \times 100$.
\item Real GDP $=$ Nominal GDP deflated by the price index: $Y_{real} = \dfrac{Y_{nominal}}{P}\times 100$.
\item Per capita growth $\approx$ output growth $-$ population growth.
\item \cle{Harrod--Domar}: $g = \dfrac{s}{c}$, where $s$ is the savings ratio and $c$ the capital-output ratio. Higher savings or a lower ICOR raise growth.
\item Harrod's three rates: actual $G$, warranted $G_w = s/c$, natural $G_n$ (set by labour-force and technical progress). Instability: if $G \neq G_w$, the economy diverges (``knife-edge'').
\end{itemize}

\textbf{Sources of growth}\par
\begin{itemize}
\item Increase in \textbf{labour} (quantity and quality — education, health).
\item Accumulation of \textbf{capital} (investment in machines, infrastructure).
\item \textbf{Natural resources} (e.g.\ Cameroon's oil, cocoa, timber) — necessary but not sufficient.
\item \textbf{Technical progress} and innovation — the key long-run source.
\item Institutional factors: political stability, property rights, good governance.
\end{itemize}

\textbf{Rostow's five stages}\par
\begin{enumerate}
\item Traditional society (subsistence agriculture).
\item Pre-conditions for take-off (investment in infrastructure, rising savings).
\item \cle{Take-off}: investment rises to about $10\%$ of national income; leading sectors emerge.
\item Drive to maturity (diversification, technology spreads).
\item Age of high mass consumption.
\end{enumerate}
Criticisms: linear and Eurocentric; ignores colonial history, dualism and external constraints of developing countries.

\textbf{Indicators of development}\par
\begin{itemize}
\item Economic: real GNP per capita, structure of output and employment.
\item Social: life expectancy, literacy, school enrolment, access to safe water, poverty and inequality (Gini coefficient).
\item Composite: HDI; also the \emph{physical quality of life index}.
\end{itemize}

\textbf{Development strategies}\par
\begin{itemize}
\item Import-substitution industrialisation vs export promotion.
\item Agricultural development and rural transformation.
\item Investment in \textbf{human capital} (education, health).
\item Infrastructure, regional integration (CEMAC, AfCFTA), good governance and the fight against corruption.
\end{itemize}

\textbf{Common mistakes to avoid}\par
\begin{itemize}
\item Confusing \cle{growth} (quantitative) with \cle{development} (qualitative): growth can occur without development (e.g.\ oil-led growth with rising inequality).
\item Using \emph{nominal} instead of \emph{real} GDP to measure growth.
\item Forgetting to subtract population growth when discussing living standards.
\item Inverting the Harrod--Domar formula: it is $g = s/c$, \emph{not} $g = c/s$.
\item Listing Rostow's stages in the wrong order or omitting the take-off condition.
\item Presenting high per capita income as proof of development while ignoring income distribution.
\end{itemize}

\textbf{GCE tips}\par
\begin{itemize}
\item In essays, always \textbf{define} growth and development first, then distinguish them with a Cameroonian example (e.g.\ oil revenue vs rural poverty).
\item For the Harrod--Domar model, state the \textbf{assumptions} (closed economy, fixed $c$, savings determine investment) and give at least two \textbf{limitations} (ignores human capital, technology, foreign aid).
\item Evaluate every model: examiners reward balanced judgement, not mere description.
\item Use diagrams where relevant (e.g.\ a rising trend of real GDP) and conclude with policy implications for Cameroon.
\end{itemize}
\end{bilan}

\findecours

\end{document}
